% Options for packages loaded elsewhere
\PassOptionsToPackage{unicode}{hyperref}
\PassOptionsToPackage{hyphens}{url}
\PassOptionsToPackage{dvipsnames,svgnames,x11names}{xcolor}
\PassOptionsToPackage{space}{xeCJK}
%
\documentclass[
  11pt,
  oneside]{book}
\usepackage{amsmath,amssymb}
\usepackage{iftex}
\ifPDFTeX
  \usepackage[T1]{fontenc}
  \usepackage[utf8]{inputenc}
  \usepackage{textcomp} % provide euro and other symbols
\else % if luatex or xetex
  \usepackage{unicode-math} % this also loads fontspec
  \defaultfontfeatures{Scale=MatchLowercase}
  \defaultfontfeatures[\rmfamily]{Ligatures=TeX,Scale=1}
\fi
\usepackage{lmodern}
\ifPDFTeX\else
  % xetex/luatex font selection
  \setmainfont[BoldFont=DejaVuSerif-Bold.ttf,ItalicFont=DejaVuSerif-Italic.ttf,BoldItalicFont=DejaVuSerif-BoldItalic.ttf]{DejaVuSerif.ttf}
  \setmonofont[BoldFont=DejaVuSansMono-Bold.ttf,ItalicFont=DejaVuSansMono-Oblique.ttf,BoldItalicFont=DejaVuSansMono-BoldOblique.ttf]{DejaVuSansMono.ttf}
  \ifXeTeX
    \usepackage{xeCJK}
    \setCJKmainfont[BoldFont=NotoSerifCJKsc-Bold.otf,AutoFakeSlant=0.15]{NotoSerifCJKsc-Regular.otf}
          \fi
  \ifLuaTeX
    \usepackage[]{luatexja-fontspec}
    \setmainjfont[BoldFont=NotoSerifCJKsc-Bold.otf,AutoFakeSlant=0.15]{NotoSerifCJKsc-Regular.otf}
  \fi
\fi
% Use upquote if available, for straight quotes in verbatim environments
\IfFileExists{upquote.sty}{\usepackage{upquote}}{}
\IfFileExists{microtype.sty}{% use microtype if available
  \usepackage[]{microtype}
  \UseMicrotypeSet[protrusion]{basicmath} % disable protrusion for tt fonts
}{}
\makeatletter
\@ifundefined{KOMAClassName}{% if non-KOMA class
  \IfFileExists{parskip.sty}{%
    \usepackage{parskip}
  }{% else
    \setlength{\parindent}{0pt}
    \setlength{\parskip}{6pt plus 2pt minus 1pt}}
}{% if KOMA class
  \KOMAoptions{parskip=half}}
\makeatother
\usepackage{xcolor}
\usepackage[paperwidth=176mm,paperheight=250mm,left=23mm,right=23mm,top=24mm,bottom=24mm]{geometry}
\usepackage{longtable,booktabs,array}
\usepackage{calc} % for calculating minipage widths
% Correct order of tables after \paragraph or \subparagraph
\usepackage{etoolbox}
\makeatletter
\patchcmd\longtable{\par}{\if@noskipsec\mbox{}\fi\par}{}{}
\makeatother
% Allow footnotes in longtable head/foot
\IfFileExists{footnotehyper.sty}{\usepackage{footnotehyper}}{\usepackage{footnote}}
\makesavenoteenv{longtable}
\setlength{\emergencystretch}{3em} % prevent overfull lines
\providecommand{\tightlist}{%
  \setlength{\itemsep}{0pt}\setlength{\parskip}{0pt}}
\setcounter{secnumdepth}{-\maxdimen} % remove section numbering
\ifLuaTeX
\usepackage[bidi=basic]{babel}
\else
\usepackage[bidi=default]{babel}
\fi
\babelprovide[main,import]{english}
\ifPDFTeX
\else
\babelfont{rm}[BoldFont=DejaVuSerif-Bold.ttf,ItalicFont=DejaVuSerif-Italic.ttf,BoldItalicFont=DejaVuSerif-BoldItalic.ttf]{DejaVuSerif.ttf}
\fi
% get rid of language-specific shorthands (see #6817):
\let\LanguageShortHands\languageshorthands
\def\languageshorthands#1{}
% Shared print conventions; the two books keep separate prose and pagination.
\usepackage{fvextra}
\usepackage{xurl}
\usepackage{etoolbox}
\usepackage{fancyhdr}
\setCJKmonofont{NotoSansMonoCJKsc-Regular.otf}[AutoFakeBold=1.5]
\setlength{\emergencystretch}{3em}
\setlength{\parindent}{1.8em}
\setlength{\parskip}{0.32em}
\linespread{1.12}
\renewcommand{\contentsname}{目录}
\renewcommand{\figurename}{图}
\renewcommand{\tablename}{表}
% Reserve enough room for three-digit page numbers in the mathematics contents.
\makeatletter
\renewcommand{\@pnumwidth}{2.4em}
\renewcommand{\@tocrmarg}{3.4em}
\makeatother
\DefineVerbatimEnvironment{verbatim}{Verbatim}{breaklines,breakanywhere,fontsize=\small}
\AtBeginEnvironment{longtable}{\small\setlength{\tabcolsep}{3pt}\renewcommand{\arraystretch}{1.15}}
\pagestyle{fancy}
\fancyhf{}
\fancyfoot[C]{\thepage}
\renewcommand{\headrulewidth}{0pt}
\fancypagestyle{plain}{\fancyhf{}\fancyfoot[C]{\thepage}\renewcommand{\headrulewidth}{0pt}}
\Urlmuskip=0mu plus 1mu
\widowpenalty=10000
\clubpenalty=10000
% Philosophy-only book typography. No shared build semantics are changed.
\makeatletter
\newcommand{\UEOTChapterHeading}[1]{%
  \vspace*{28\p@}%
  {\parindent\z@\raggedright\normalfont
   \hyphenpenalty=10000\exhyphenpenalty=10000
   \interlinepenalty\@M
   \fontsize{18}{23}\selectfont\bfseries #1\par\nobreak
   \vskip 24\p@}}
\let\@makechapterhead\UEOTChapterHeading
\let\@makeschapterhead\UEOTChapterHeading
\renewcommand\@pnumwidth{3em}
\renewcommand\@tocrmarg{4em}
\renewcommand\section{\@startsection{section}{1}{\z@}%
  {-3.0ex\@plus-1ex\@minus-.2ex}%
  {1.5ex\@plus.2ex}%
  {\normalfont\Large\bfseries\raggedright\hyphenpenalty=10000\exhyphenpenalty=10000}}
\makeatother
% A long note stays with a text-bearing page rather than leaving an orphan tail.
\interfootnotelinepenalty=10000
\counterwithout{footnote}{chapter}
% English conventions apply only to the philosophy English edition.
\renewcommand{\contentsname}{Contents}
\renewcommand{\figurename}{Figure}
\renewcommand{\tablename}{Table}
\hyphenation{self-maintenance object-hood counter-factual}
% Keep long contents titles intact; retain right-aligned page numbers.
\makeatletter
\let\UEOTOriginalTOC\tableofcontents
\renewcommand\tableofcontents{\begingroup
  \hyphenpenalty=10000\exhyphenpenalty=10000
  \UEOTOriginalTOC\endgroup}
\patchcmd{\l@chapter}{\rightskip\@pnumwidth}
  {\rightskip\@pnumwidth plus 1fil}{}
  {\PackageError{ueot-philosophy}{Cannot apply chapter TOC layout}{Check book.cls compatibility.}}
\patchcmd{\l@chapter}{\parfillskip-\@pnumwidth}
  {\parfillskip-\rightskip}{}
  {\PackageError{ueot-philosophy}{Cannot apply chapter TOC page alignment}{Check book.cls compatibility.}}
\patchcmd{\l@part}{\rightskip\@pnumwidth}
  {\rightskip\@pnumwidth plus 1fil}{}
  {\PackageError{ueot-philosophy}{Cannot apply part TOC layout}{Check book.cls compatibility.}}
\patchcmd{\l@part}{\parfillskip-\@pnumwidth}
  {\parfillskip-\rightskip}{}
  {\PackageError{ueot-philosophy}{Cannot apply part TOC page alignment}{Check book.cls compatibility.}}
\makeatother
\ifLuaTeX
  \usepackage{selnolig}  % disable illegal ligatures
\fi
\usepackage{bookmark}
\IfFileExists{xurl.sty}{\usepackage{xurl}}{} % add URL line breaks if available
\urlstyle{same}
\hypersetup{
  pdftitle={Objects, Evolution, and Meaning},
  pdfauthor={Degui Qian (钱德桂)},
  pdflang={en},
  colorlinks=true,
  linkcolor={MidnightBlue},
  filecolor={Maroon},
  citecolor={Blue},
  urlcolor={MidnightBlue},
  pdfcreator={LaTeX via pandoc}}

\title{Objects, Evolution, and Meaning}
\usepackage{etoolbox}
\makeatletter
\providecommand{\subtitle}[1]{% add subtitle to \maketitle
  \apptocmd{\@title}{\par {\large #1 \par}}{}{}
}
\makeatother
\subtitle{An Open Philosophy of Nature in Pursuit of Unity}
\author{Degui Qian (钱德桂)}
\date{2026-11-11}

\begin{document}
\frontmatter
\maketitle

{
\hypersetup{linkcolor=MidnightBlue}
\setcounter{tocdepth}{1}
\tableofcontents
}
\mainmatter
\begingroup\small\setlength{\parskip}{0.25em}

\textbf{Bilingual Scholarly Revision · English Volume · Version 3.0}\\
11 November 2026

The original ideas and paper are by \textbf{Degui Qian (钱德桂)}. This book takes \emph{A Universal Evolutionary Object Theory (UEOT): A Framework and Research Program for a Unified Science} as its intellectual source. It reorganizes the argument in light of subsequent work on the mathematical core, with assistance from ChatGPT in analysis and editorial revision. The original paper and its author statement remain part of the historical record. The conceptual argument developed here stands on its own; technical propositions, proofs, and their conditions belong to the companion mathematical work.\footnote{\hypertarget{ph-note-1}{}This is a collaborative revision based on the original paper and subsequent research. Attribution of the original authorship follows the contribution statement on page 43 of that paper. Appendices B--D locate the new chapter structure, normative arguments, and mathematical feedback in this edition. New formulations not individually confirmed by the original author are not retrospectively attributed to the original paper.}

\textbf{Edition, contributions, and rights.} The Chinese and English volumes together constitute Version 3.0. Chapter numbers, arguments, and mathematical qualifications correspond chapter by chapter; each volume has its own pagination. The English is a complete scholarly rendering of the Chinese text in this edition. Three native AI sub-agents undertook philosophical, Chinese-language, and English-language review respectively, with a lead agent responsible for the final text. This is an internal editorial review, not external academic peer review, and it does not mean that the author has individually confirmed every new formulation.

Expression validly licensed under CC BY 4.0 in the past remains available under that license, including its permissions for translation and adaptation. This edition does not revoke those rights. No additional general license is granted for new expression in this edition to the extent that it is protected by copyright; the respective rights holders reserve their rights. This does not limit prior licenses, statutory exceptions, or ideas and methods not protected by copyright. Third-party materials retain their respective rights. Source versions, the scope of revisions, and the same bilingual rights statement are given in \href{https://github.com/MurphyHoops/UEOT/blob/pub/philosophy-bilingual-edition/philosophy/PUBLICATION_RIGHTS.md}{Publication Rights} and this edition's editorial records.

\par\endgroup

\chapter{Preface: From the Wish for Unity to a Philosophy We Can Share}\label{ph-preface}

A person watches a river flow, enters a city shaped by institutions, returns home to the claims of affection, and later asks, alone, what makes a life meaningful. These experiences can seem to belong to different worlds. Matter obeys laws; living things act to survive; societies coordinate through rules; people choose for reasons. We have learned to describe these worlds separately, yet often struggle to place them within a single world. Does the success of physical explanation make purpose and value superfluous? Does meaning become less capable of public recognition the more closely it belongs to personal experience?

The original UEOT paper answered this division with considerable ambition. It sought to connect matter, life, mind, and meaning through the organization of objects, a proposed pair of evolutionary concerns, and interactions across levels. Its first merit lies in the question it refused to set aside: if human beings belong to nature, our purposes, commitments, and self-understanding must also have a place in the natural world. We should not acknowledge the real conditions of all our activities while excluding the experiences that matter most to us from reality.\footnote{\hypertarget{ph-note-2}{}Original paper, pp.~1--3 and 40--41. The shared world of matter, life, mind, and meaning is the overarching question taken up here. This edition separates the original universal and ultimate claims according to the distinct arguments they require.}

This book continues that inquiry. It argues that one promising route to understanding the world is to examine how relatively stable organizations arise through change, how relations sustain them, how they acquire capacities to engage with their surroundings, and how such capacities can develop into action guided by reasons. The force of this approach cannot rest on repeating that everything is unified. We must show which divisions a proposed unity overcomes, which differences it preserves, and which questions it enables us to answer.

Why, for example, should yesterday's promise still bind a person whose body has replaced much of its material? In what sense do interacting participants form an organization? Does being able to predict a system mean that we know what constitutes it? Does a living being's capacity to maintain itself supply every end worth pursuing? Is the long-term stability of a society enough to establish its justice? These questions are connected, but they cannot all be answered in the same way. A mature philosophy of unity explains how their differences are related.

The position developed here may be called \textbf{relational realism about objects}. The objects principally considered in this book are organizations in the real world that can withstand certain changes, preserve certain distinctions, and enter into interactions. Our questions determine which distinctions we seek to track. Whether those distinctions survive intervention, support prediction, and withstand shared scrutiny is not ours to decide at will. The position is also a form of \textbf{process-oriented naturalism}: persistence is an achievement within change, and the activities of mind and evaluation have bodily, historical, and social conditions. Naturalism here is a commitment to explanation---to understanding ourselves, as far as possible, within a common world. It does not stipulate in advance that every explanation must be reduced to one vocabulary.

A philosophy of this kind requires two forms of courage. One is the courage to ask large questions and acknowledge the connections still to be discovered between bodies of local knowledge. The other is the courage to revise an initial formulation when the argument demands it. Replacing the claim that all stable existence requires a geometric loop with an inquiry into the causal dependencies that sustain an organization does not abandon the question of existence. It gives that question a more definite object of dispute. Replacing the claim that every object has a unique ultimate purpose with a demand to identify the mechanisms and evaluative relations that give action its direction does not expel purpose from nature. It prevents purpose from becoming a label that can be attached retrospectively to any outcome.

The six parts follow a continuous course. Part I examines change, relations, and objects. Part II considers openness, repair, memory, and the possibility of action. Part III turns to purpose, value, and self-revision. Part IV concerns composition, scale, and evolution. Part V addresses mind, freedom, morality, and society. Part VI returns to method, the unity of the cosmos, and truth. The argument can be read without consulting the companion mathematics. A concise correspondence at the end directs readers who wish to check its formal background.

The unity sought here does not require every voice to agree. A philosophy that can accommodate real disagreement, explain its own revisions, and submit to correction by the world is better suited to a shared pursuit of truth than a system that declares inquiry complete.

This edition makes four public revisions to the original theory. Persistence does not universally guarantee a nonempty failure structure for every permitted family of interventions. A two-dimensional Π/Φ representation requires testable conditions of low rank and independent anchoring. Natural evolution does not universally guarantee either a unique optimal direction or a unique attractor. An advantage in growth does not confer moral legitimacy. These revisions retain responsibility for the original explanatory questions while acknowledging that some of the stronger initial claims cannot be sustained. The attribution identifies the source of the ideas and paper; it does not retrospectively ascribe the collaborative revisions to the original author's words.

There is also a positive result. In a class of organizations with an explicitly specified resource supply, in which both program code and execution components can be damaged, one and the same program can pay for recovery by using redundant information and then select growth-promoting actions whose ranking is derived from actual persistence and birth events. Faithful inheritance changes the distribution of programs under explicit finite-population conditions. The joint example of regeneration and optimal growth uses the same damaged initial state, action permissions, and resource account; population feedback separately specifies the founding population and resource scale within that same generative model. This is a limited but substantive achievement. It reduces the need to assume in advance that a system ``can maintain itself,'' ``has fitness,'' or ``must select the optimum.'' It does not establish that every organization behaves this way, and it does not make growth a substitute for all the questions of purpose and meaning pursued in this book.\footnote{\hypertarget{ph-note-3}{}The common-process organization and feedback developed in this edition are set out in the seven parts of the companion mathematical volume's definitive core overview and in the source constructions, formal proofs, and execution evidence at \nolinkurl{formalization/ueot-core/docs/compression/theory_completion/unified_closure/core_theory_closure_v1}. The original inputs are finite physical components, program code, permitted requests, noise, initial states, and actual supply. Success, conditional means, and optimality are derived from the same source. No universal support in nature is claimed; empirical status remains UNVERIFIED.}

Four kinds of discussion have distinct responsibilities here. Definitions explain how concepts are used. Arguments about mechanisms explain how conditions produce consequences. Interpretive claims ask why those consequences help us understand objects. Normative arguments ask which reasons agents should accept. Mathematical proof constrains claims about mechanisms; formal correctness alone cannot decide what matters philosophically. A complete account makes room for counterexamples and independent argument, and identifies universal claims without empirical support as unverified.

\part*{\phantomsection\hypertarget{ph-part-1}{}\addcontentsline{toc}{part}{Part I: Objects amid Change}{\Large Part I}\\[1.2em]{\LARGE Objects amid Change}}

\chapter{Chapter 1: Why Ask about Unity?}\label{ph-ch-01}

Consider a bridge. Engineers discuss its load-bearing capacity, urban planners the communities it connects, and residents the changes it makes to journeys to school or medical care. None of these descriptions excludes the others, and none can do all their work. The bridge's use depends on a physical structure; building that structure depends on institutions and purposes. A description confined to material motion would leave unanswered why this bridge was built. A discussion confined to public meaning would not tell us whether it might collapse.

Unity deserves investigation because real consequences cross disciplinary boundaries. A reason helps bring an institution into being; the institution changes material arrangements; those arrangements reshape people's opportunities. When explanation stops abruptly at a disciplinary border, we risk treating what remains unexplained as a mysterious power from another world. The first demand of unification is for connections between neighboring explanations that can be examined: how does an input at one level become a condition at another, and how much of a conclusion survives a change in description?

An objection immediately arises. Perhaps the world is not unified at all, and the search for unity merely expresses an aesthetic preference. The objection has force. Simplicity makes a theory easier to grasp; it does not oblige nature to be simple. Yet a more modest case for unity comes from practice. In different activities we encounter the same world, and the consequences of action do not stop propagating when our terminology changes. Tracing these consequences is part of our explanatory responsibility. It warrants a demand for compatibility, though not yet a demand that every law follow from one final formula.

We should therefore distinguish a shared world from a single theory. The former claim says that the bridge's physical processes and its residents' activities do not occur in unrelated universes. The latter says that one particular theory can adequately explain them all. This book adopts the former as a realist position. It treats the latter as a research ambition whose merits must be earned through results. We can affirm a shared world while leaving the prospects of a single theory open.

A second objection appeals to the success of specialization. If the sciences already work well, why seek more general concepts? Because specialization leaves questions that the division of labor cannot itself resolve. When does a system count as an object? When has the same object changed, and when has another object appeared? What do control, repair, selection, and understanding have in common? Such questions recur across disciplines. A clear shared vocabulary can prevent equivocation and may allow an argument developed in one setting to become useful in another.

\begingroup\interlinepenalty=10000

A shared vocabulary must still answer to the particulars of each case. Describing bridges, cells, and institutions with the same names explains no new relation by itself. A unifying framework must show what we understand better once the questions are connected: how an object's boundary limits its action, how maintenance in turn depends on resources and history, and why a particular success still cannot replace a judgment of value. If a connection can be sustained only by changing the meanings of its terms, the claim to unity should be withdrawn. This book seeks to make different explanations constrain one another. A stronger system must earn its standing through the understanding those connections provide.\footnote{\hypertarget{ph-note-4}{}Original paper, pp.~1--3; Core, Chapters 0, 28, and 31. The philosophical case for unity is argued in the chapter itself. The composability of modules is not treated as evidence that all theories of nature have already been unified.}

\par\endgroup

\chapter{Chapter 2: How Things Exist in Processes}\label{ph-ch-02}

Ordinary language separates things from their changes. First there is a cup; then the cup becomes warm, is damaged, or is moved. This convenient idiom can suggest that a thing has an inner core wholly independent of processes. From another perspective, however, the cup's shape depends on a material structure continuing to withstand its surroundings; a vortex depends on flow; an institution depends on people repeatedly enacting its rules. These things persist in different ways, but none can be explained apart from what actually happens.

A process perspective does not require us to abolish things. It asks how they become stable subjects of discourse. If the world consisted only of unrelated moments, we could not connect one observation with the next. The emergence of an object means that certain distinctions can be tracked through a course of change: it can serve a similar role after moving, preserve an organization after replacing some material, or respond predictably after disturbances within a specified range.

``Object'' here refers primarily to a concrete entity that can be tracked through an interval of change. It is not a definition that everything real must satisfy. A flash, a collision, or a formative process that has not yet acquired a stable shape can also be real. Relatively stable organization explains why some things can continue as the same objects of inquiry; it does not explain why only stable things exist. Different time windows and questions require different distinctions to be preserved.

The word ``certain'' matters. A block of ice can retain material continuity while losing its shape. An institution can retain its name while changing its function. A person can revise a belief while remaining responsible for a past action. Objecthood does not freeze every property forever. It requires us to say which properties support the tracking relevant to our question, and how actual processes sustain that tracking.

Does this reduce objects to convenient classifications made by observers? Consider a forest. Two people may divide it into different regions for different purposes; neither can stop a fire from spreading simply by naming the regions differently. Classification serves an inquiry, while the relations through which fire spreads constrain what the classification can accomplish. A description's selectivity does not make the whole reality of its object a product of selection.

A small model illustrates this objectivity. Suppose a device has a visible bit representing its operating condition and an internal bit. Two actions respectively make the next state satisfy or fail the operating condition, while the internal bit is sampled afresh. If the internal bit does not affect the consequences of these actions, grouping the four complete states into two classes by the visible bit answers questions about operation and control exactly. If we instead ask for the internal bit's particular value, that description is insufficient. The description depends on the question; whether it preserves the consequences of action remains strictly constrained by the same process. Chapter 14 returns to the model to ask how an exact description can also support judgments about persistence and purpose.

Chapter 7 introduces an organizational model in which recovery has a cost. Here we can first consider how it changes the partition into objects. The model contains two programs, each healthy or damaged, giving four initial states. The first always chooses the care action, which favors the original individual's continued existence. The second chooses an action favoring the production of new individuals when healthy, and repairs itself before providing care when damaged. If the sole question is whether a birth occurs next, the three initial states that perform care share a birth probability, while the healthy second program forms a separate class: exactly two classes are needed. If the description must preserve the joint records of birth counts at all finite future horizons, however, the damaged second program must also be separate. After recovery, its healthy successors perform a different action, and already in the second period their mean birth count differs from that of the first program. Exactly three classes are then needed. The first program's healthy and damaged initial states incur different first-period repair costs, but their bodily outputs after recovery and their subsequent birth records can be identical.

If the full-future birth task also retains the actual first-period repair fee, the difference between zero cost when healthy and three units when damaged separates the last pair, requiring exactly four classes. Full public observation likewise distinguishes all four initial states. This classification fixes the internal programs and resource rules; it does not replace an organizational description that preserves the entire external control authority. It shows that the choice of partition depends on the capacity being protected, while the adequacy of the partition is constrained by actual responses. The same response once does not guarantee the same future. Identical predictions do not, by themselves, erase differences in the cost of realizing them.\footnote{\hypertarget{ph-note-5}{}Part III of the definitive core overview in the mathematical volume, on minimal representation of one-step formation responses and finite futures, and Section 23.4, on the three-response boundary in the same regeneration model. This supplement provides a written proof within the common model and exact computational checks. It adds no Lean theorem covering the corresponding complete classification and supplies no empirical validation.}

A finite description does not require the whole observable world to be finite. A model may operate within a finite domain while other accounts assign probability to cases beyond it. Deleting those cases for convenience and then declaring complete agreement between model and account changes the question. A legitimate compression can retain an explicit ``outside the domain,'' so that omitted cases still carry their probabilistic consequences, even if differences among those cases are no longer distinguished. A finite representation becomes a tool of inquiry by stating clearly which comparison it preserves.

The process perspective thus resists two extremes. Treating things as isolated entities unaffected by relations leaves their formation and disappearance unexplained. Treating them as arbitrary narratives leaves unexplained why some classifications support prediction and intervention while others fail. A defensible middle position understands objects as relatively stable organizations within real processes. Their boundaries, time scales, and salient properties may vary with the question, but the world must supply sufficient continuity and constraint for them to exist as described.

Stability need not mean rest. A body regulating its temperature exchanges matter; a team keeping a commitment adjusts its division of labor; an effective rule is enacted in successive situations. Stability may consist in unchanging states or in a recognizable manner of changing. These cases belong together because they pose a common problem of persistence, not because their microscopic mechanisms must be alike.

Returning persistent things to the processes of their formation is not a direction first proposed here. Whitehead's process philosophy already placed the becoming and relations of actual occasions at its center. His ``actual entities'' also carry experiential and metaphysical commitments that this book does not adopt. A shared emphasis on process does not make the two concepts of an object synonymous. The approach here is more limited: which actual processes sustain an identifiable organization, and which interventions would interrupt that persistence? A process orientation supplies a question; the particular mechanisms of persistence still need to be explained.\footnote{\hypertarget{ph-note-6}{}Whitehead, \emph{Process and Reality} (1929; corrected edition, 1978), Part I, Chapter II, pp.~17--22. These passages concerning process, relations, and actual entities were read for this edition. No reading of the entire book is claimed, nor is its complete metaphysics adopted as an established premise of UEOT. Full bibliographic details appear in Appendix D.3.}

\begingroup\interlinepenalty=10000

Being and becoming are therefore no longer mutually exclusive. Processes of becoming produce relatively stable things, and stable organizations constrain what can subsequently become. An object is both a result of past processes and a condition of future ones. Understanding it requires us to ask how it arose and what it enables to happen next.\footnote{\hypertarget{ph-note-7}{}Core, Chapters 2--4, 8, and 11--12, on processes, representation, predictive identity, persistence, and transport through a history. Section 36.8 of the mathematical volume explains how a finite representation preserves all escaping probability mass and the corresponding probabilistic comparison under an explicit one-sided support condition. ``Relatively stable organization'' is the philosophical expression used here to bring these distinctions together.}

\par\endgroup

\chapter{Chapter 3: Relations, Boundaries, and a Shared World}\label{ph-ch-03}

A door separates two spaces while connecting them in a particular way. The boundary of an object often performs both roles. A bodily boundary regulates exchange; an institutional boundary determines membership; an instrument's boundary determines which signals enter a measurement. Understanding a boundary merely as an outer shell misses these functions of selection and transformation.

A relational perspective emphasizes that many capacities arise from connections with other things. A key's capacity to open a lock depends on the lock as well as the key's shape. An expression's meaning depends on habits of use and on hearers. An organization's power depends on other participants' recognition and cooperation. None of this entails that everything is a relation and that its relata need no explanation. A key still needs a material realization; recognition must be sustained by actual people and institutional processes. Relations change how we understand their terms without relieving us of the need to explain how they are realized.

Emphasizing relations does not establish that objects can be eliminated from the world. Worrall's discussion of structural realism asks what can survive a change in scientific theory: accounts of entities' natures may change while some mathematical relations remain effective. That is an argument about continuity in scientific knowledge. It does not directly supply the individuation of objects required here. We must still explain which activities realize a set of relations, how its boundary is identified, and under which interventions it retains the relevant capacities. A structural diagram alone cannot establish the objectivity of relations.\footnote{\hypertarget{ph-note-8}{}Worrall, ``Structural Realism: The Best of Both Worlds?'' (1989), especially the Fresnel--Maxwell discussion on pp.~115--121. The main text compares structural continuity across theoretical change with the task of identifying objects. It does not claim that Worrall's article supports UEOT.}

Three complementary questions help locate a boundary. Which changes can enter from outside while the organization under discussion persists? Which internal differences alter its outward responses? Which interventions would disable that way of responding? These questions concern openness, distinguishability, and constitutive dependence respectively. They may suggest similar boundaries or different ones. A database's access boundary need not coincide with the boundary of the physical equipment storing its data.

Saying that boundaries depend on questions can easily be mistaken for saying that reality has no boundaries. The question selects a mode of testing; the test reveals real connections. Treating a building and its heating system as one object can be useful for predicting indoor temperature, while repairing its pipes requires distinctions among subsystems. Both descriptions can be true because they track different, compatible dependencies. A partition that repeatedly excludes the exchanges decisive for an outcome, by contrast, needs to be revised.

Even when tasks and interventions have been specified, the evidence may support more than one way of individuating objects. The realism defended here does not guarantee a unique boundary for every question. Faced with this underdetermination, we should retain competing descriptions and seek circumstances in which their consequences differ. If they really cannot be distinguished for the present task, we should acknowledge that limit. Reality can constrain description by excluding errors without immediately selecting a uniquely correct picture.

We must also distinguish a relation not yet observed from a relation that does not exist. A limited capacity to detect things does not determine the extent of the world. An object may be affected through intermediate processes without being able to identify that influence directly. For human knowledge, this means that the shared world exceeds present experience. For science, it means that an account of an interface should leave room for omissions to be discovered.

The strongest conclusion of a relational view is not that everything determines everything else. It is that explanation must assign relations precise roles. Some transmit information, some constrain action, some maintain material conditions, and some are only statistical associations. Distinguishing these roles prevents us from moving unlicensed from correlation to causation, dependence to identity, or cooperation to a common interest.

\begingroup\interlinepenalty=10000

An object's boundary is consequently open to discussion and revision while remaining constrained by reality. It organizes shared experience and reminds us that any perspective which makes part of the world visible may leave another important connection outside its frame.\footnote{\hypertarget{ph-note-9}{}Core, Chapter 6, on openness, environmental interfaces, causal direction, and boundary conditions; Chapter 13, on the distinction between access and physical damage.}

\par\endgroup

\chapter{Chapter 4: Ω: How an Organization Sustains Its Conditions}\label{ph-ch-04}

Consider a device that requires power, control, and cooling. Power allows its controller to operate; the controller regulates cooling; cooling keeps the equipment capable of working. There is a returning dependence here. Understanding it requires more than drawing a circle on paper: we must identify the activities that jointly sustain the conditions of further activity. This is the central reason for retaining the idea of Ω.

The original paper used a closed loop to express the internal organization of persistence. The intuition captures something static descriptions often overlook: the order of many systems has to be continually realized. Their activities consume conditions while also re-establishing them. Life offers particularly clear cases, but related questions arise in technical systems and institutions.\footnote{\hypertarget{ph-note-10}{}Original paper, pp.~4--5, on Ω and existence, and pp.~37--38, on the proposed use of topological data. This chapter distinguishes the shape of a trajectory, a causal cycle, and constitutive dependence.}

Organizational closure, however, must acquire its meaning from the relations among activities. A device driven by an external rhythm may trace an elegant cycle without maintaining the conditions that drive it. A system regulating around a stable state may display no conspicuous circle in its macroscopic trajectory while still containing essential feedback. The shape of a trajectory, a cycle in a causal graph, and a regenerative dependence among mechanisms answer different questions.

How can Ω become an explanatory concept? One place to begin is with the ways an organization can fail. First specify what must be maintained; then ask which removals would prevent its maintenance. If two independent routes can each supply a required resource, removing one may do no harm, whereas removing both causes failure. Integrity therefore concerns redundancy, substitution, bottlenecks, and their combinations, as well as the presence or absence of a loop. A complex system can have many geometric cycles and depend on one fragile common support. An apparently simple one can have several routes to recovery.

Experiments involving removal themselves need interpretation. Removing a component can sometimes release an inhibition and activate a backup. In such a case, removing more does not necessarily make things worse. We should describe the actual response instead of requiring the world to conform to a convenient monotone picture. Integrity is always relative to the organizational capacity, interventions, and time window being examined.

When a declared capacity can in fact fail under permitted interventions, patterns of failure and redundancy can reveal the dependencies sustaining it. If those interventions cannot disable it, the failure structure for that family of operations may be empty. That an object has real conditions does not mean that every task and family of operations yields the same nonempty Ω.

A simple example makes the limitation clear. A device stores a bit that can be read later, together with a redundant bit that does not affect the readout. If permitted interventions alter only the redundant bit, the device does not lose its readout capacity under those operations. It remains a nontrivial predictive object. Enlarging the intervention family to cut the actual power or readout pathway may produce failure. But that enlargement changes the question; it cannot retroactively establish the same failure claim for the original one.

``Real processes have conditions'' is thus a philosophical position here, whereas a particular Ω is a mechanistic structure specified by a task, an operation family, and a time window. Its objectivity consists in responses being answerable to shared testing under the same conditions. An experiment that deliberately overlooks a crucial pathway may still be impoverished. We can criticize it by extending publicly comparable interventions and testing its predictions. Relativity to a boundary does not make all boundaries equally good; the reality of a dependence does not establish one uniquely correct boundary.

\begingroup\interlinepenalty=10000

For a philosophy of unity, Ω therefore sets a positive research demand: identify the constitutive dependencies of persistent organization and explain how they are preserved, replaced, and rebuilt. This permits comparisons among different forms of maintenance while leaving particular mechanisms to the appropriate accounts. A symbol can guide several inquiries. What connects them is the question they share and the explanatory links that survive examination.\footnote{\hypertarget{ph-note-11}{}Core, Chapters 12--14; companion mathematical volume, Chapters 32--34. Representations by minimal failure structures require the relevant finiteness, monotonicity, and interpretation of interventions. These conditions are not stipulated as the universal essence of every object.}

\par\endgroup

\chapter{Chapter 5: The Same Object: Identity as a History}\label{ph-ch-05}

A chair repaired many times, an institution whose entire staff has changed, and a person continually forming memories all make ``Is it still the same?'' a difficult question. If identity required identical material, many familiar cases of persistence would disappear. If it required only identical function, two simultaneous copies would count as one object. Identity demands more than either material or function supplies alone.

We can begin by distinguishing similarity, equivalence, and historical identity. Similarity concerns closeness in certain properties. Equivalence concerns indistinguishability for a class of questions. Historical identity requires a directed relation of continuation. Two files with identical contents are equivalent for a reading task but may have different origins. A file revised through an authorized process has different contents while belonging to the history of the same document. Our inquiry determines which relation matters; our argument must state which one it uses.

Historical identity need not rest on one permanently unchanging property. It may be supported by a sequence of warranted transmissions: material connections, preserved organization, inherited records, assumed responsibilities, and mutually recognized rules. At each step, we must explain why the later state continues the object being tracked. Successive transmission can accommodate substantial change overall, yet encounter genuine ambiguity in cases of division, fusion, or replacement.

Branching is a particularly revealing test of the distinction. Suppose a team's members, archives, and working relations continue in two successor teams. Each may have real historical continuity with its predecessor. That alone cannot establish that each is strictly identical to the one original team: given the transitivity of strict identity, it would make the two now separate teams identical to one another. Legal or organizational rules can determine who inherits the name, contracts, and responsibilities, but settling those questions is not the same as creating past facts. We should preserve both successors' historical relations and then specify whether the present judgment requires strict identity, an entitlement to inherit, or functional continuation.

This also illuminates personal identity. People can change their convictions while remaining responsible for their actions because the history on which responsibility rests usually extends beyond their present beliefs. Severe changes in capacity, coercive circumstances, or memory impairment may nonetheless alter what it is reasonable to attribute to them. Recognizing this complexity helps us avoid both fixing a person into an immutable essence and using change to erase that person's history with others.

Repair displays the directionality of identity especially clearly. Its aim is usually to restore a capacity of this particular object while preserving specified continuities. Producing a functionally identical replacement may accomplish a different task. The replacement can be equally valuable without pretending to have the same history. Reproduction is a third relation: descendants inherit aspects of an organization while entering the world as other individuals.

\begingroup\interlinepenalty=10000

Understanding identity as a constrained history explains how growth and constancy can coexist. What merits preservation need not be every past state. It may be the relations through which commitments, capacities, and a shared life continue. An object's past is not a static label. It continues to act in the present through transmissions that can be explained.\footnote{\hypertarget{ph-note-12}{}Core, Chapters 11, 13, and 15. For the subsequent precise relations among object roles, formation and binding, and genealogical identity, see Chapter 34 of the mathematical volume.}

\par\endgroup

\part*{\phantomsection\hypertarget{ph-part-2}{}\addcontentsline{toc}{part}{Part II: Openness, Memory, and the Possibility of Action}{\Large Part II}\\[1.2em]{\LARGE Openness, Memory, and the Possibility of Action}}

\chapter{Chapter 6: Open Existence and the Finite Body}\label{ph-ch-06}

A living being that must breathe, feed, and rest cannot preserve itself as a complete thing sealed off from the world. It must admit some things, release others, and sustain its organization through these exchanges. For many ways of existing, openness is a condition of existence itself, rather than an addition to something already complete.

Openness also brings exposure. An interface that admits nutrients may admit harmful substances; a channel carrying news may carry misinformation; a cooperative relationship may create dependence. The issue is not whether more exchange is always better, but how exchanges contribute to maintenance and development. Protecting a boundary and expanding a connection can each be reasonable in different circumstances.

The original paper distinguished relatively stable objects from those strongly dependent on environmental exchange, and used degrees of openness to describe a life cycle. This can guide comparison. Once formed, some structures depend comparatively little on continuing regulation; other organizations must keep working to maintain themselves. Yet openness is difficult to reduce to one scale on which everything can be ranked. What is exchanged, at what speed, through which interfaces, and how indispensable the exchange is to a particular capacity may all change the comparison.\footnote{\hypertarget{ph-note-13}{}Original paper, pp.~4--5 and 9--11. Openness is treated here as an organizational question whose contents and interfaces require specification, rather than as a confirmed universal law of a single life cycle.}

In living beings, the body gives openness a center that cannot simply be relocated at will. The body is a resource for action and a place where damage occurs. We can abandon a task more readily than we can abandon the body that bears its consequences. Because action carries this burden, the world is not an undifferentiated collection of information. Some changes make further activity possible; others permanently diminish a capacity.

This establishes an elementary sense in which something is at stake. For a being sustaining an organization, improving and damaging its conditions have different consequences. That difference grounds functional explanations of life: organs, regulation, and behavior can be understood through their contributions to organizational continuation. A full account of moral value still needs to explain who merits concern, how conflicts should be addressed, and which reasons can be publicly justified. Embodiment gives the question of value someone or something that bears the consequences; it does not supply every answer.

Finitude also makes it necessary to distinguish sustainability on average from what can be borne on this occasion. A device receiving enough resources on average may stop before the next supply arrives. Someone with sufficient income over the long term may lack support at the moment care is urgently needed. If averages conceal timing, reserves, and irreversible harm, judgments of sustainability lose essential content.

\begingroup\interlinepenalty=10000

The body's philosophical significance is that it returns abstract possibility to a reality with times, places, and costs. Action is no costless choice from an unlimited menu. Acting consumes conditions, changes surroundings, and alters what can be done later. Living openly requires attention both to the capacity for continued exchange and to the capacity to bear its consequences.\footnote{\hypertarget{ph-note-14}{}Core, Chapters 8--10 and 22; mathematical volume, Chapters 34--35, on the differences among resources, finite paths, actual survival, and trusted repair. The examples contrasting averages with what a particular path can bear explain this conceptual distinction.}

\par\endgroup

\chapter{Chapter 7: Repair, Life, and Autopoiesis}\label{ph-ch-07}

An object repaired from outside differs significantly from an organization able to participate in its own repair. In the latter, some conditions of maintenance become part of the organization's activity: detecting damage, mobilizing resources, rebuilding pathways, and checking whether recovery works. Self-maintenance thereby acquires a richer meaning. The object continues its operation and responds in an organized way when that operation is threatened.

Repairing oneself does not mean creating everything from nothing. Repair requires some remaining capacity to detect, act, and coordinate, together with a way of distinguishing success from failure. If the mechanism interpreting a repair program has entirely ceased to function, the program's text cannot repair the system merely by continuing to exist. If recovery materials are unavailable, even a correct sequence of steps may be impossible to carry out. Such dependencies are part of an explanation of repair, rather than objections to the very idea.

Suppose a device can restore its control program from a backup. We can reasonably say that it has recovered one of its capacities. Yet the backup reader, the rules executing the restoration, and the mechanism connecting program instructions with real actions all contribute to that success. Explaining how these supports themselves arose requires a further investigation of their origins. Distinguishing formation from recovery reveals levels within living organization instead of hiding every difficulty under the name ``autopoiesis.''

The origin of a program also has an informational aspect. Two local state records may be identical even though the program's meaning differs because of an unrecorded interpretive convention or formation history. Reading only the local record cannot distinguish what the two cases require. Reliable recovery must preserve enough of the relevant distinction or acquire it through another channel. Recovery thus depends on the organization of information and realization; this does not mean that a program can exist without physical conditions.

Autopoiesis is philosophically attractive because it seeks to move beyond an understanding of life as a collection of parts. The parts have roles within an organization they jointly sustain, and concrete processes within that organization produce and replace the parts. The explanation has a circular structure, but it need not be logically circular. We can explain how processes depend on one another at different times and test those dependencies through disturbance. An objectionable circle uses the capacity for self-production as both the conclusion and its sole reason. An explicable cycle is a regenerative network operating through time under real conditions.\footnote{\hypertarget{ph-note-15}{}Original paper, pp.~4--5, on self-maintenance; Core, Chapters 9 and 26; mathematical volume, Chapters 34--35, on formation, the origin of programs, and repair within the same underlying process. Gaps in program provenance concern information and the conditions of realization.}

A mathematical realization already allows this demand concerning origins to be checked step by step. Finite registers store a target program. An interpreter reads the actual registers and determines whether to leave them unchanged or write to them. Environmental faults damage target bits, while the program itself is protected by the model's assumptions. Repair, failure, cost, and offspring actually produced are preserved in one record. Only an actual birth entails that a child inherits the program and has the capacity for a subsequent repair. The advance consists in deriving ``the program can be executed'' and ``the offspring inherits a basis for repair'' as results of execution instead of supplying them as separate guarantees. The assumptions of a reliable interpreter, protection, and inheritance remain explicit. The construction does not make every support self-produced.\footnote{\hypertarget{ph-note-16}{}Mathematical volume, Sections 39.1--39.3; the actual executions and offspring applications of \texttt{RegisterRepairSource}, \texttt{RegisterFaultSource}, and \texttt{RegisterBirthMaintenance}. Entry into the fault carrier provides a capacity for repair, not permanent health.}

Autopoiesis is not another name for self-repair. Maturana and Varela analyzed networks of processes in which the components produced participate in realizing and regenerating the network, constituting it as a unity. Restoring a program from a backup therefore establishes a particular capacity for recovery. A claim to have explained life requires a fuller account of organizational relations. Early autopoietic theory also related attributions of purpose and function to an observer's description. Later organizational theories of function, such as that developed by Mossio and colleagues, ask a further question: when a component contributes to self-maintenance, how is it itself subject to production and maintenance by that organization? These approaches do not form one uncontested doctrine. Organizational existence, the attribution of function, and the explanation of normativity are connected tasks that no single definition can perform on its own.\footnote{\hypertarget{ph-note-17}{}Maturana and Varela, \emph{Autopoiesis and Cognition} (1980), pp.~78--81 and 134--135; Mossio, Saborido, and Moreno, ``An Organizational Account of Biological Functions'' (2009), especially pp.~828--833. The relevant passages in the original works were read for this edition. The observer's standpoint in early autopoietic theory and the later organizational account of function must not be combined as though they were one conclusion.}

An objection remains: if enough supports are listed, could any machine be called alive? This is a reason not to define life by one isolated capacity. Self-maintenance, metabolic exchange, organizational continuity through development, reproduction, and heredity can form connected clusters; different inquiries may focus on different clusters. Understanding life through a connected set of capacities has greater explanatory promise than asking one simple boundary to decide every case.

This affects our understanding of aging and death. Death may involve more than a numerical threshold being crossed: the relation between maintenance and recovery may cease to be organizable. Continuing activity in one component does not establish the continuation of the living whole. Nor does the reuse of its material by other organisms mean that the original individual persists. These distinctions allow us to recognize both the continuity of living processes and the irreplaceability of an individual life.

Recovery also exposes the difference between an endpoint and a history. An interrupted and subsequently restored operation may end in a state as satisfactory as that of an uninterrupted operation. Their histories remain different. Repair makes certain subsequent activities possible again, preserving or rebuilding the pathways they require. Damage already suffered, opportunities missed, and responsibilities incurred do not disappear when the endpoint is restored. An account of persistence must record capacities regained while retaining an irreversible past.

Repair therefore combines fidelity to the past with creativity directed toward the future. It must preserve enough continuity for the restored object to be the one that matters in the inquiry. Yet restoring its capacity may require different materials and routes. A developed form of self-maintenance preserves, within limits, an organization capable of further change rather than endlessly repeating one state.\footnote{\hypertarget{ph-note-18}{}The later formal autopoiesis result retains inputs concerning the life cycle, program provenance, encoding and decoding, physical dynamics, and bridges of identity; see Chapter 34 of the mathematical volume. The account of life as a connected cluster of capacities is a philosophical proposal advanced in the main text.}

A further model brings these questions into the same body. Damage affects not only the components looked after by the program but also the program code and the main switch that enables its execution. Each line of code is stored in three actual bits. A fault flips the bit in use and turns off the execution switch. Recovery still depends on a finite trusted boot component. It performs only majority reads and writes at a cost, without hiding an additional copy of the ``originally correct program.'' The remaining backups preserve the distinctions needed for recovery, available resources pay for it, and the main program subsequently reads the restored code and acts. For the damaged state examined by the model, cutting the boot channel prevents recovery. If every encoding is erased so that two different original programs leave the same current physical state, the recovery mechanism cannot reliably reconstruct both from that state alone. This distinguishes having a description of the steps of repair from having the conditions for carrying them out.

In this model, the action \texttt{care} favors retention of the original individual. The action \texttt{bloom} raises the birth rate while lowering the original individual's one-period probability of continued existence. The first program always chooses \texttt{care}. The second chooses \texttt{bloom} when healthy and, when damaged, pays three units for repair before choosing \texttt{care}. Its subsequent actions still follow its actual code, which its descendants inherit faithfully. The same program can therefore restore its conditions of execution and, when healthy, turn to an action with a greater consequence for population growth.

Calling the second program optimal requires us to specify the comparison. The class contains every policy that uses only the permitted public history and submits its requests before the current random influence occurs. It may choose care, birth-favoring action, or stopping, and may randomize with arbitrary real-valued probabilities; requests by members of the same generation may be correlated. Over-budget requests are actually refused, and the consequences of refusal are included. Every policy has the same resources, observation permissions, and rules of code inheritance. Within the model under its stated conditions, the second program attains the highest exponential growth factor of expected population in this complete comparison class. That does not maximize the original individual's chance of surviving the next step, guarantee that a finite population will never become extinct, or show that the program understands its purpose. Chapters 14 and 19 discuss the different evaluations and questions of lineage.\textsuperscript{\hyperlink{ph-note-3}{3}}

Keeping the same body, funding, refusal, and inheritance rules while varying only admissible probability parameters reveals where the fixed programs are optimal. Each program, in its own damaged body, can still pay the three-unit recovery fee and perform care. Their backup codes differ, but each program's repair and subsequent growth belong to its own continuous history. The second program, which favors birth, is optimal when the gain in births at least offsets the loss of persistence of the original identity and its long-run growth factor is at least one. The first, always-care program is optimal under the opposite ordering of actions when the care growth factor is at least one. At a tie with a common growth factor of at least one, both can be optimal. Even if care by itself produces fewer than one successor on average, the second program may still achieve supercritical growth through its different behavior in healthy and damaged states.\footnote{\hypertarget{ph-note-19}{}Part IV of the mathematical volume's definitive core overview, on exact growth and attainment by the internal programs over the strictly positive probability domain, and the \href{https://github.com/MurphyHoops/UEOT/blob/3d443ec29eb17af1fc00699a74e2d42d73516d81/formalization/ueot-core/docs/compression/theory_completion/unified_closure/core_growth_classification_v1/README.md}{growth-parameter classification package}. The original bodies, public observations, wallet, \texttt{stop}, refusal, and faithful inheritance are retained. The two fixed programs have growth factors \(k_c\) and \((1-q)k_b+qk_c\); the optimum over all permitted policies is the maximum of those two factors and \(1\). The formal Lean contract provides two-sided exponential bounds at every finite horizon, the least base for a common upper bound, and necessary and sufficient conditions for attainment by the internal programs. The nth-root interpretation is given in the mathematical text. The critical eleven-twelfths example, the two-period population counterexample, and a prefix containing one paid recovery are supported by written derivations and exact computation. Selection for growth does not automatically supply a normative reason.}

When the expected populations generated by both fixed programs decline, the original request menu still lets a damaged body submit an over-budget request, whose actual refusal preserves it. Both fixed programs then fall below the optimal growth factor of one available to the full control class. The admissible probability conditions of this classification ensure that the first care action has a positive probability of leaving a successor. The result limits continuous operation under these two fixed rules. A permitted history-dependent policy can pay for one actual care action on the same damaged founder and then choose continuation or refusal according to the state, while still attaining a factor of one. The possibility of death at the first step, the actual fee, and all subsequent resource supplies remain in the account. Such history-dependent control is a different behavioral rule from either fixed internal program. Their failure to be optimal does not prove that all actual recovery is incompatible with optimality. Refusal preserves the body because the model explicitly says so; this is no basis for prescribing universal inactivity to natural systems.

How can this count as self-regeneration if the same mechanism does not rebuild its own boot component? The objection cannot be dismissed by redescribing everything nearby as internal to the organization. External supply, a trusted boot component, and the execution conditions actually restored must be specified separately. Material dependence on an environment also differs from which constitutive relations the organization can reproduce. The new contribution presented here is a mechanistic example of stated scope: damage, retained information, recovery actions, actual payment, and offspring inheritance are connected within one generative process, and some execution conditions genuinely enter the organization's own maintenance activity. The boot conditions that are not rebuilt remain explicit. Incorporating deeper conditions requires a new mechanistic proof; the existing result does not meet every condition of life or full autopoiesis.

This also distinguishes regeneration from dependence on an environment in general. A component held up by a fixed support depends on it without necessarily producing or restoring it. An organization that pays for materials, reads backups, and reconstructs an execution gate actually changes its conditions of maintenance. Its execution and responses to intervention justify attributing a regenerative component to it. They do not thereby establish life, consciousness, or the full standing of a subject. A specified capacity has been restored; further classifications still need justification.

\begingroup\interlinepenalty=10000

The restored organization can nevertheless encounter noise that causes death. Successful local repair does not eliminate the possibility of overall extinction, and a healthy endpoint does not establish an uninterrupted history. The positive conclusion is that finite resources and sufficient information can realize certain recoveries. The boundaries of failure identify the distinctions that must be preserved.\textsuperscript{\hyperlink{ph-note-3}{3}}

\par\endgroup

\chapter{Chapter 8: How Memory Becomes Part of the World}\label{ph-ch-08}

Imagine entering a corridor twice. On each occasion the same two doors are visible, but a different sign was shown at the entrance: the first indicated that the left door was passable, the second that the right one was. An agent who remembers the sign can make different, justified choices despite the identical present scene. History has entered the present. Through memory, it changes which possibilities still deserve consideration.

Memory does more than store a finished past. It turns experience into constraints on the future, allowing an object to distinguish situations that look alike but differ in the consequences of action. Some consequential differences in the world are absent from present observation yet can be preserved in a continuing record of observations and actions.

This shows why informational sufficiency is always relative to a question. If we ask only whether an action satisfies a constraint in every case still possible, we may need only to know which cases have not been ruled out. Comparing risks, expected benefits, or costs generally requires more: probabilities, evaluations, and temporal conditions. The same memory can be sufficient for one judgment and insufficient for another.\footnote{\hypertarget{ph-note-20}{}Core, Chapters 4--5 and 27; mathematical volume, Chapter 37, on support derived from actual histories and finite feasibility. A support set suffices for the proved result of zero occupancy in uncertified regions. This does not establish its sufficiency for preserving arbitrary probability predictions or rankings of value.}

The distinction can be tested within an actual process. Fix the permitted actions and observation scheme, and take the complete distribution of observations producible by each finite action sequence as the predictive response. Its update can then be derived from the actual posterior. The response suffices to predict those future experiments and permits comparison of the information retained by other sufficient memories under the probability law of actual pasts. An existing controller may, however, read records from much earlier. Compressing the predictive state cannot silently delete memory that the controller still uses. Two actual histories with the same present appearance may each have positive probability while requiring this controller to choose different actions. A comparison of complete action records between that controller and the compressed process must preserve the full history used by the same controller.\footnote{\hypertarget{ph-note-21}{}Mathematical volume, Sections 39.7 and 39.11--39.12. The actual response recursion for general finite action words and minimality modulo null sets on the source are each stated under their own conditions. Path comparisons for full-history programs preserve the same program and actual history. Identical responses under different control memories are not treated as identical arbitrary actions.}

What counts as enough to remember can also be derived from actual action. Suppose a finite process returns only two recorded outcomes: whether a birth occurred and whether a writing fee was paid. For an object that still exists and is healthy, the absence of a birth after choosing not to repair for the moment---the action called \texttt{idle} in the model---may provide evidence of a fault. The same zero record from a damaged or nonexistent object may supply no new information at all. Memory that merely accumulates outwardly identical signals can record evidence that was never obtained. What it needs to preserve is what a signal could distinguish given the action and circumstances in which it occurred.\footnote{\hypertarget{ph-note-22}{}Mathematical volume, Sections 39.82--39.88. The observer, count posterior, shared graph, and underlying execution are proved for the fixed source, the 0/1 menu, and a finite clock. Reachable primitive examples establish the insufficiency of bodily state alone, the possibility of omitting tail information, and the failure of that omission when address 2 becomes available.}

The same present can also carry different futures. Two actual histories can leave the same body and the same wallet at the same time, yet induce different posteriors about an unknown environment and hence different probabilities of the next birth. Conversely, two other histories can retain different individuals at the inaccessible end of a population record and different parent--child links while permitting exactly the same future decisions under the specified dispatch permissions. A compact controller may merge the latter distinction while the underlying world retains both histories. The omission has a justification and a limit: if previously inaccessible addresses become available, the omitted differences may again affect action, and memory must be refined.

The resulting compact policy is not merely an abstract existence claim. It can update from permitted histories, select actions at shared nodes, and be executed step by step through the original transactions. Every request can be paid for before execution; failures, empty requests, and the world's evolution after intervention stops remain in the account. A smaller decision memory does not imply a smaller complete world or establish that two objects have the same identity. This finite realization shows how to reduce duplication in an action table. It does not prove that an external computer has become an object's internal cognitive organ, still less derive two universal drives from two counts.

What remains possible must likewise come from the adopted world model and actual records, rather than merely from what appeared in training data. A rare transition's absence from a finite record does not establish its impossibility. Treating absence from a sample as impossibility in the world produces an unduly optimistic view of available action. A useful compression of memory states which distinctions it retains and what justifies excluding others.

Memory is thereby connected with identity. An organization may change its internal coding while retaining its predictive and practical roles, provided it preserves the relevant historical distinctions. Preservation of those roles does not, by itself, decide whether it remains the same historical object. Informational continuity in memory and historical continuity of its bearer must be connected in the particular inquiry. Two identical memory copies can have the same predictive capacity and begin different histories from the moment of copying.

Memory also has an easily overlooked social form. Promises, archives, language, and shared habits keep histories active that no individual could remember alone. A person can travel according to rules in an unfamiliar city because the results of earlier coordination have been incorporated into the public environment. This external memory extends agency while creating dependencies. Who keeps a record, who can alter it, and whose experiences it excludes can change the possibilities of a shared world.

\begingroup\interlinepenalty=10000

The information at issue here is information available for memory and action: distinctions preserved and transmitted within an organization and a task, with an effect on later responses. It requires an actual carrier but cannot be understood simply by counting the material involved. This use does not exclude statistical information or other concepts of information, nor does it treat every physical trace as a memory an agent can read. Memory connects a changing object with the world's history, so that the present is more than an isolated moment.

\par\endgroup

\chapter{Chapter 9: Possibility before Optimality}\label{ph-ch-09}

When asking what should be done, we often turn immediately to the best, quickest, or most profitable action. A prior question remains: which plans can actually be carried out while retaining the required conditions through subsequent changes? Without such plans, an optimization may merely compare options that do not exist.

Consider a task with two stages. Two first-stage operations both satisfy the immediate requirement. One leads to a second-stage state from which no further operation can comply; the other preserves a subsequent choice. Acceptability at the first step does not automatically make the whole task feasible. Conversely, the existence of a fully feasible policy does not imply that every initially acceptable choice can be rescued.

A policy here is not a lucky trajectory. It is a way of responding: it can choose appropriately after different observable consequences occur. If an action can have two outcomes outside the agent's control, the agent needs a way to complete the task after each outcome that has positive probability. Displaying one favorable outcome does not establish the action's full feasibility.

Limits of observation make the requirement stricter. An agent may not know which hidden state obtains and so cannot choose differently for states it cannot distinguish. What it can implement is a single choice based on its observation history, satisfying the requirement for all states that history has not ruled out. New observations can then update the set of remaining possibilities and guide the next action.

Feasibility becomes an organizational capacity extending into the future. In a model with explicitly specified finite states, possible transitions, and a time limit, we can work backward from the endpoint. Which sets of states admit one common acceptable action whose every possible observational successor still allows the remainder of the task to be completed? The resulting claim that a way forward exists is more definite than general optimism. It depends on an accurate account of possibilities and on the agent actually possessing the relevant interfaces for observation and action.\footnote{\hypertarget{ph-note-23}{}Core, Chapters 8 and 18; mathematical volume, Chapter 37. With a fixed finite source, specified observations, a time horizon, and explicit constraints, that chapter proves the equivalence between feasibility of the initial support and the existence of a full history-dependent policy. It also supplies a rigorous example of a reachable dead end.}

This may sound like a demand to eliminate every risk. We first need to specify the constraint at issue. An action can lie within a range whose model error has been certified while its outcome still has a positive probability of failure. It can preserve procedural permissibility without guaranteeing that the body will never be harmed. Feasibility acquires content only through an explicit aim. Conflating different kinds of safety makes it easier to misunderstand what has been guaranteed.

This finite conception of feasibility also matters in human practice. Preserving future options need not mean postponing decisions forever. Some opportunities arise only after an irreversible commitment. What matters is identifying what must be retained in advance, which risks are reasonable, and who bears failure's consequences. An account of freedom should begin here as well: choices must have significance, but there must also be choices that can actually be made.

The organizational model reveals a detail easily concealed by the word ``optimal.'' A healthy body can prepay either of two actions. A damaged body must first be repaired and can afford only the cheaper action. A request for the more expensive one is refused by the transaction, which preserves the original body. An action's name alone does not tell us what it actually does. If the environment cannot support growth, repeated refusal may even preserve the body better for the time being than continued execution. This failure case belongs in the classification. Changing refusal into death or removing it from the menu merely to rescue an attractive optimality claim would change the problem.

\begingroup\interlinepenalty=10000

Comparisons made one row at a time also require joint implementability. That each of two individuals requests no more than six units does not mean a shared wallet containing six units can pay both requests at once. Observability is limited too: a controller cannot act differently in observationally identical situations merely because an analyst knows their hidden difference. Only after permissions and resources have been accounted for can a theory of optimality describe its mathematical choices as realizable actions.\textsuperscript{\hyperlink{ph-note-3}{3}}

\par\endgroup

\chapter{Chapter 10: Learning Must Return to the World of Action}\label{ph-ch-10}

When a system performs well on past data, we are inclined to entrust it with new tasks. Yet three questions remain distinct: how were the data generated, which situations will the new actions encounter, and how do the data determine those actions? The claim that the system has learned becomes clear only when these questions are placed within the same generative process.

Suppose training tries only one class of mild operations, while deployment repeatedly uses another class because the model selects it. Even accurate knowledge of the first class does not justify treating its average error as the error of the second. Knowledge of responses must extend to the places where it will actually be used. Insufficient coverage is not simply another name for too little data. Repeatedly collecting familiar cases may still leave untouched the situations that determine deployment risk.

Selection itself changes the probabilities that matter. Different training records may select different forms of control, which in turn generate different futures. An evaluation of learning must consider how training occurs, how it brings about selection, and how the selected control generates deployment outcomes. Separate averages for training and deployment can miss the relation created by selection.\footnote{\hypertarget{ph-note-24}{}Core, Chapters 16--17, 20, and 27; mathematical volume, Chapters 36 and 38. The probability measure of the training experiment, the rows sampled during training, and occupancy during actual deployment have distinct mathematical roles.}

These distinctions permit a reasonable form of guarantee. If one training experiment can jointly calibrate all registered responses with high probability, a policy subsequently selected from those data can be assessed for fit, coverage, and risk on the same calibration event. We need not pretend that all possible training records are equally good. Rare, misleading records may remain. A guarantee should account explicitly for their probability while using the stronger conclusions available on reliable records.

Registration does not require allocating the same number of observations to every question. Before the experiment, we may assign different sample sizes according to importance and cost, then account jointly for the possibility that some responses remain inaccurately known. The observations must actually come from the registered process, and subsequent selections must stay within the scope of this shared evidence. Conclusions needed only on reliable records must also connect on the same set of reliable records. Combining advantages obtained under different screening conditions may promise a result that no single actual experiment can deliver.

Unity in learning, understood this way, does not collapse training, representation, and action into one activity. We first establish the reliable range of source responses, then explain which distinctions a new representation preserves, and finally subject selected actions to that knowledge. Evidence already obtained jointly can support several subsequent choices without charging the same uncertainty again whenever the description changes. A new description that omits a connection decisive for success must, however, answer for that omission. What can be reused is an acquired justification, not untested correctness.

The calibration event is a condition on actual error in the analysis. An operator may not know whether a particular training run satisfies it. Observable empirical quantities must be distinguished from probabilistic guarantees that those quantities approximate the true responses. A complete analysis allows us to account clearly for the risk of the whole experiment even when its actual error is not directly observable.

The mathematical work underlying this book provides a complete finite example. Training is an explicitly specified random experiment; the data genuinely change the selection; both selected controls have positive probability and generate different actual deployment distributions. The risk of the whole process can be calculated within one joint experiment. The example shows that these components can actually be connected, rather than assuming that a conclusion simply remains valid after selection. It also shows that coverage of every response row actually visited does not make the failure probability zero.\footnote{\hypertarget{ph-note-25}{}The registered selection example in Chapter 38 of the mathematical volume; numerical values and scope are given in Appendix C here. It establishes actual joint risk in the specified generative experiment. It does not infer an arbitrary external deployment mechanism from the marginal training law.}

Epistemologically, this challenges an overly simple picture of learning. Learning does not extract the world's truth intact from data and apply it without loss. It builds representations with a scope, chooses how to use them, and lets the consequences of new actions test those representations. The model's world and the world in which we live require continuing calibration.

Uncertainty therefore has a legitimate place in intellectual progress. It may indicate that evidence cannot yet distinguish two accounts, or that a model lacks sufficient coverage somewhere. Requiring every judgment to be affirmative or negative can force a knowledge system to disguise its blind spots as certainty. Mature learning includes recognizing when action is justified, when further information is needed, and when the question itself should change.

Learning also has ethical consequences. Who decides to explore, who pays for trial and error, and whether those affected can challenge a model's omissions all shape the meaning of progress. A system may continually improve its predictive score while ignoring the consequences borne by excluded people. It may then be learning technically while becoming less perceptive about shared life.

When learning, repair, and action draw on the same finite resources, this responsibility extends further. An exploratory step is no free preparation outside the world of action. It may repair the object or cause failure; spent resources cannot reappear at deployment. Evaluation should begin with the initial unknown source and continue to the agreed endpoint, carrying forward the actual object after exploration, its remaining resources, and the knowledge gained. Stopping paid intervention does not stop the environment from acting.\footnote{\hypertarget{ph-note-26}{}Mathematical volume, Sections 39.75--39.81. A fixed target, actual information permissions, an unknown source drawn once and retained, and a finite wallet and clock jointly delimit the result. Mathematical optimality does not replace normative justification of the purpose.}

An unknown cause that persists over time connects successive experiences. A failure may increase doubt about the environment's reliability, making a retry that was previously worthwhile no longer worth its cost. More information does not therefore always have positive value. A procedure may already work well enough for every possible source, so identifying the source more precisely changes no action. What matters in another case may be knowing which existing body remains usable, rather than finding a more precise name for an unknown mechanism. Whether further inquiry is worthwhile depends on the consequences it can change, the time left, and its cost.

The companion mathematics illustrates these distinctions in a finite-population model with actual births and task assignment. Whether another retained object can serve as a backup depends on whether the environment damages objects not yet dispatched to act and on which individuals the operator may access. In one three-step task, a procedure fixed in advance already attains the joint decision optimum. In another four-step task, switching objects in response to actual recorded outcomes strictly improves the result. Adaptation matters through these particular relations; its value is not guaranteed by the mere description ``able to learn.'' Nor does an optimum in an accounting model exhaust the value of life.

Before deciding how to act, we may not even know the extent of a fault: does it damage only the part accessed on this occasion, or every backup at once? Two mechanisms can produce exactly the same records for the root object while calling for different subsequent actions. Repeating the experiment on this original object cannot replace a test of the backup relations. Which mechanism is operating and how the various bodies now stand must then remain jointly uncertain. A reconstruction valid under only one mechanism cannot be presented as an observed fact.\footnote{\hypertarget{ph-note-27}{}Companion mathematical volume, Sections 39.89--39.95, and the \href{https://github.com/MurphyHoops/UEOT/blob/3d443ec29eb17af1fc00699a74e2d42d73516d81/formalization/ueot-core/docs/compression/theory_completion/unified_closure/mechanism_decision_closure_v1/README.md}{mechanism-decision delivery package}. The mechanism is fixed; diagnosis and action share the original wallet and time horizon. A weighting witness is not an identified natural prior, and worst-case regret is not the only normatively admissible evaluation.}

In an explicit three-step model, the companion mathematics proves an exact upper bound on the discriminating power of paid experiments. When a decision has to be made, the available evidence still cannot determine the mechanism. Yet action has not thereby lost its justification. A policy using independent private randomness attains the minimum worst-case regret among all permitted history-dependent policies. Regret here is the best expected return over the whole horizon if the fixed mechanism were known in advance, minus the evaluated policy's expected return under that same mechanism. Both obey the same observation permissions, budget, and time horizon; knowing the mechanism does not mean foreseeing future random outcomes. Random choice does not secretly observe the environment, and the analysis does not let an adversary see the private randomness before choosing the most unfavorable mechanism. The policy bears uncertainty under specified information conditions and limits the possible resulting loss to the best attainable bound.

\begingroup\interlinepenalty=10000

The capacity to discriminate, the capacity to act without decision loss, and the worth of adopting the objective remain three different judgments. At two steps, the mechanisms are indistinguishable yet share an optimal action with a nonzero payoff. At three steps, a strictly positive decision loss is unavoidable. Evidence separating the mechanisms only at the end of the third step cannot be booked as a learning benefit already realized within that same horizon. Such results make epistemic limits and practical responsibility clear together, without treating every unknown as a reason to stop acting.

\par\endgroup

\part*{\phantomsection\hypertarget{ph-part-3}{}\addcontentsline{toc}{part}{Part III: Purpose, Evaluation, and Self-Revision}{\Large Part III}\\[1.2em]{\LARGE Purpose, Evaluation, and Self-Revision}}

\chapter{Chapter 11: How Purpose Enters Nature}\label{ph-ch-11}

When a plant grows in a particular direction, we may say that it is seeking light. When a person takes a detour, we may say that she is keeping a promise. Both explanations invoke a possible later result to make present activity intelligible, but their grounds differ. The plant's functional relations require an account of growth and maintenance; the promise also involves the person's understanding of reasons. The language of purpose can connect these domains, provided we explain how the connection works.

This book holds that purpose need not enter the world as a force from outside nature. It can belong to organized activity: certain differences count as success or failure, a system responds to departures from them, and these responses help shape what it does next. In such an organization, a purposive explanation identifies a real relation of directedness. No future event has to act backward on the past. Present mechanisms must instead preserve, enact, or revise some requirement concerning the future.

The first level concerns function. It examines the role a regulatory activity plays in maintaining an organization and what makes that role its function. A second level concerns history: some forms of organization have been retained because of their past effects. A third concerns intention: an agent can represent outcomes, compare reasons, and organize activity accordingly. A fourth concerns shared commitment: participants can establish a common task through rules. These levels can connect without becoming identical.

A pump can be broken and still have the function of pumping. A fragment that happens to lodge in a leak may temporarily save the equipment without thereby having leak prevention as its function. A beneficial present effect is not a complete criterion for attributing function. We must also ask why the part arose or was retained, what role it plays in explaining the system's capacities, and how it enters the organization's own maintenance. Wright emphasizes how an effect explains why something is there; Cummins asks how the capacities of parts jointly realize a capacity of the whole; organizational accounts of function examine the relation between a part's contribution and the organization's production and maintenance of that part. These approaches do not answer exactly the same question. One expression, ``for the sake of,'' cannot erase their differences.\footnote{\hypertarget{ph-note-28}{}Wright, ``Functions'' (1973), pp.~155--161; Cummins, ``Functional Analysis'' (1975), pp.~760--764; Mossio, Saborido, and Moreno (2009), pp.~828--833. The accidental blockage of a leak is an illustrative construction of this book. An accidental origin does not itself prevent a contribution from acquiring functional status under further conditions. The distinction used here is that a currently beneficial effect does not suffice on its own.}

Accordingly, saying that a regulatory process ``ought'' to work first requires specifying its function and operating conditions. Malfunction is not moral wrongdoing, and functional success does not establish something as worth pursuing. A mechanism that maintains an organization well can serve an unjust institution. A person can also reject the action most favorable to continued existence in order to honor a better-justified commitment. The transitions from function to intentional purpose and from purpose to moral reasons each need a further argument.

The same dynamics can serve different purposes. A device that moves forward or backward can be used to approach a target or to move away from it. Knowing how it moves does not determine which evaluation to adopt. Even if it consistently moves in one direction, we must distinguish its control rule, external constraints, chance initial conditions, and any reasons for acting. Calling every actual outcome what the system really wanted would deprive purposive explanation of its discriminating force.\footnote{\hypertarget{ph-note-29}{}Original paper, pp.~6--8; Core, Chapters 18 and 26--27; mathematical volume, Chapters 34--35. The same process can serve different goals. Identifying a purposive mechanism requires the relevant faithfulness conditions and an independent source.}

How can an account of purpose be more than an arbitrary attribution? It must yield consequences that can be compared. If the proposed reward, signal, or conditions of maintenance change, does behavior change accordingly? If two purposes are equivalent with respect to past behavior, can a new situation make their predictions diverge? In explaining human action, we must also take agents' statements of their reasons seriously and examine how those statements relate to their conduct. A report is neither an unquestionable last word nor noise that may be dismissed in advance.

An objection holds that if purpose depends on evaluation, it can only be a subjective projection. That conclusion is too quick. Evaluations can themselves be realized through organizations, institutions, and historical relations. That a departure from a set temperature triggers regulation is part of a mechanism; that a contract makes conduct answerable to obligations is part of a shared practice. Both have real effects. What must be resisted is the further inference that, because evaluation has a real basis, every existing thing determines one unique evaluation.

Purpose in nature is consequently layered and revisable. Life establishes basic stakes; learning enlarges the futures that can be represented; language makes reasons available for discussion; institutions let commitments extend beyond a single encounter. Human purposes neither descend from nowhere nor follow by automatic deduction from the desire to survive. They arise in nature and history, and can in turn evaluate the conditions under which they arose.

Mechanism and ranking can also be connected in a concrete way. In the register model described earlier, actual births and paid writes are produced by the same execution step. Once price, discounting, and the time horizon are specified, the same program is, under stated cost conditions, no worse than any causal policy in the original comparison class. When opportunities for repair have nonzero probability, the mean births resulting from actually writing each of two values, together with the mean payment, can also identify the target, cost, and fault rate, thereby determining policy rankings for the specified evaluation. The evaluation still requires justification. What the mechanism supplies is a testable relation between operations and the resulting ranking. Mathematical access to exact means does not by itself solve learning from finite data or establish ethical legitimacy.\footnote{\hypertarget{ph-note-30}{}Mathematical volume, Sections 39.3--39.5; \nolinkurl{RegisterFaultControlModel}, \texttt{RegisterFaultJointRecord}, and \nolinkurl{RegisterFaultPurposeInformation}. Cost, price, and discounting are specified separately. At total fault, parameter readout degenerates; this does not establish nonidentifiability of the purpose-ranking class in which all policies are indifferent.}

The same mechanism can constrain what finite observations support. A judgment based on a complete batch of probes, preserving its actual births and payments and using a cost interval fixed in advance, can certify an existing repair program's future performance under the specified conditions. It can also return an unresolved judgment when the data contain no useful information or the cost lies close to the threshold. The costs already paid during training must remain in the account. The evaluation's time horizon matters as well: a repair may incur a larger immediate cost while enabling the same program to obtain a positive return over a longer future. An unfavorable short-term judgment does not automatically justify abandoning repair forever. The mathematical contribution is a finite, checkable assessment; why its evaluation is worth adopting remains a further question.\footnote{\hypertarget{ph-note-31}{}Mathematical volume, Section 39.16; \nolinkurl{RegisterFaultFinitePurposeProbe} and \nolinkurl{RegisterFaultFinitePurposeInfinite}. Certificates from complete finite data, actual infinite records, indistinguishability under a fixed readout, and reversals between one-step and infinite-horizon evaluations are each stated with their conditions. Optimal future evaluation does not include optimal total net return from training and does not imply risk-free execution.}

The timing of payments also limits what purposes can be realized. In one model with no external replenishment, a strictly positive fault probability, and a strictly positive repair price, the same program, starting from a damaged state with any finite reserve, has a positive probability of overspending. This does not establish almost-sure bankruptcy.\footnote{\hypertarget{ph-note-32}{}Mathematical volume, Section 39.18, \nolinkurl{RegisterFaultFeedbackResource}.} A different training protocol explicitly includes an external supply. It accounts for every fault, birth, and cost during training and ensures that resources cover each training prefix. When training returns a favorable repair judgment, it returns the same policy, which continues from the parent actually left at the end of training and retains the actual ledger. The optimality assessment begins at deployment. It does not establish optimal total net return across training and the subsequent future, nor does it account for every form of energy consumption. Why this evaluation should be adopted remains a separate normative question.\footnote{\hypertarget{ph-note-33}{}Mathematical volume, Section 39.24, \texttt{RegisterFaultFundedProbe}: a complete new protocol of \(4N+1\) steps. The actual final parent differs from that in the older source.}

A mechanism can supply a grounded direction without first being assigned an arbitrary reward. In the regenerating organization developed in this edition, the next census actually retains some parents and adds some new individuals. Once every failure and death is included, we can compare how actions affect these contributions. If the question is how quickly expected population grows over the long run, the ranking arises from actual persistence and transmission. Repeated choices change how many organizations remain able to act, establishing a traceable connection between evaluation and feedback in the world.

The question chosen still matters. The continued existence of an individual parent, the number of descendants, mean population growth, the avoidance of catastrophe, and the conservation of materials are different concerns. One action in the model increases offspring contribution and long-run growth while reducing the original individual's probability of surviving the next step. No algebraic simplification removes that difference. A growth objective is one branch of purpose with a basis in mechanism. It neither supplies UEOT with a unique ultimate purpose nor determines how a person ought to live.

An exponential growth factor must also be understood through what it leaves out. Two expected-population sequences that differ throughout by a fixed positive multiplier can have the same growth factor, although their finite-horizon populations and the loss represented by that multiplier remain different. In one critical example, the second program, starting from its own damaged founder, has expected population eleven-twelfths at every positive horizon. A policy that always invokes actual refusal retains one individual at every such horizon. Both attain a growth factor of one, while refusal yields the larger expected population at every positive finite horizon. In another example, the second program attains the highest long-run growth factor, but the policy maximizing population after two steps chooses refusal for damaged bodies at the final step.\textsuperscript{\hyperlink{ph-note-19}{19}}

These differences give assessments of purpose testable content. Changing the comparison horizon or the quantity evaluated may change the optimal action. The mathematical classification gives a common upper bound for all originally permitted history-dependent policies and identifies where the two existing programs attain it and where they do not. It does not claim that these programs maximize population at every date. A long-run growth advantage does not extinguish reasons concerning finite losses, the original individual, or material conservation. Whether those concerns should take precedence depends on the purpose under discussion and on normative argument.

\begingroup\interlinepenalty=10000

This limited naturalization nevertheless adds explanatory content. The claim that some activities help an organization persist need not be an observer's retrospective compliment. We can specify how a damaged body, actual operations, resources, and inheritance jointly produce the difference, and which interventions would alter the ranking. This account does not show that the organization subjectively understands the ranking, still less that it can give reasons for its own goals. Functional directedness, behavioral goals, and reasoned intention can connect, but each connection needs its own argument.

\par\endgroup

\chapter{Chapter 12: Π and Φ: Two Concerns of Action}\label{ph-ch-12}

Someone preparing for a long journey faces two different difficulties: obtaining enough resources and knowing where to go. Abundant supplies may not rescue a traveler with a false map; an accurate map may not rescue one without water or rest. Action commonly involves both the conditions for accomplishing something and the ability to discern its consequences. The original proposal of Π and Φ captures these two concerns, neither of which can readily substitute for the other.\footnote{\hypertarget{ph-note-34}{}The original paper describes Φ as ignorance or lack of adaptation on p.~6, then discusses exploration, information, culture, and value on pp.~12--15 and 24--28. This chapter maintains connections among these themes while requiring each measure to be defined independently.}

This book uses Π to prompt questions about resources, capacities for realization, and the returns from action. It uses Φ to prompt questions about residual information, the identification of structure, and the adequacy of prediction. A particular model must specify their meanings. Does ``resource'' refer to energy, time, materials, or an exchangeable return? Does an information measure concern uncertainty, predictive error, distinguishable responses, or model complexity? Useful connections can be established between these answers, but a common letter does not confer a common unit.

The distinction helps explain how efficient activity can damage long-term capacity. A system that commits every available resource to immediate output may lose the ability to monitor, learn, or maintain itself. Conversely, continual inquiry without action consumes the conditions that support knowing. Resources and knowledge sustain one another within a process: investigation takes time, sound judgment can save resources, and new action produces new information.

They are not, however, two forces that always oppose one another. A better arrangement may reduce consumption and improve information together; an injury may undermine both. A trade-off must arise from actual feasible conditions, rather than from an a priori demand that every change move along a line of opposition. We can compare costs and benefits only after establishing which changes can really be made.

Exploration also exposes the ambiguity of ``more information is better.'' Someone investigating an unfamiliar field may discover many more unknowns than were previously recognized. Knowledge of the extent of one's ignorance increases, while some measures of uncertainty rise. This need not be a failure to learn. It may reveal that earlier certainty depended on omission. The amount of information, the degree of confidence, and the accuracy of understanding require separate assessment.

Truth, morality, and culture likewise cannot simply be treated as quantities of information. A person can know many facts and refuse responsibility; an institution can gather data exceptionally well while silencing the people concerned. If Φ stands indiscriminately for information, order, morality, and civilization, the breadth of the concept conceals the connections it ought to explain. We should first specify how some kind of knowledge contributes to a worthwhile activity, then examine the corresponding evaluation.

Information's value for decisions also depends on relevance. News revealing the result of a coin toss wholly unrelated to the task increases what is known without improving the choice. If information is free, its acquisition leaves the opportunity set unchanged, and it can be ignored, its addition cannot force the original decision to worsen. Actual inquiry, however, usually consumes time, attention, and other resources. Concern for information must eventually return to its relation to action and life.

The same regenerating model provides a concrete test. Under a fixed care program, public records permit three separate questions: did the original individual survive, was a new individual born, and was a new individual born already damaged? Let \(s\) denote survival under care, \(b\) the birth probability, and \(q\) the output-damage probability. The three probabilities are \(s\), \(b\), and \(bq\). Over an open parameter range allowed by the model, their local changes have three independent directions. No single smooth two-dimensional intermediary can preserve all three completely. Section 23.4 of the mathematical volume gives the short proof, with determinant \(b>0\). This is a calculable counterexample within a model; empirical support is a separate matter. It limits the claim of a two-dimensional structure while leaving intact the model's existing regeneration and growth results. Treating resources and information as two important families of questions carries a different burden of argument from claiming that every actual response has only two degrees of freedom.\textsuperscript{\hyperlink{ph-note-5}{5}}

\begingroup\interlinepenalty=10000

The philosophical strength of the dual concern lies in drawing attention to the material conditions of action and the capacity to discern the world. They are related, yet neither can wholly perform the other's role. It does not lie in deciding beforehand that nature has exactly two drives. A stronger shared theory of two drives would have to establish the relevant common mechanisms across environments and withstand independent measurement and counterexamples. That open question does not deprive the dual concern of value as a starting point for analysis.\footnote{\hypertarget{ph-note-35}{}Core, Chapters 23 and 27; mathematical volume, Chapters 33--34. A common two-dimensional rank condition is necessary in particular models. Low rank alone does not identify physical axes, and information gain is not identical to decision value.}

\par\endgroup

\chapter{Chapter 13: Scales of Evaluation and What Remains Invariant}\label{ph-ch-13}

Two teachers can use different grading scales while ranking the same assignments in the same order. Numerical representation and evaluation are therefore distinguishable. Knowing what someone chose rarely tells us the exact scores used to make the choice, still less how many independent components of value contributed to a total.

If an overall evaluation combines a return with a cost, both terms may admit simultaneous changes that cancel in the total. Every final choice then remains the same while the interpretations of return and cost change. Identifying the components requires further grounds: independent measurement, fixed units, constraints imposed by the mechanism, or interventions that can act on them separately. A good fit to the overall objective does not by itself identify its components.\footnote{\hypertarget{ph-note-36}{}Core, Sections 18.4, 23.2, and 23.5; mathematical volume, Chapters 33 and 35. A gauge characterization preserving exactly the same numerical objective must be distinguished from purpose equivalence that preserves only the ordering of policies.}

Such freedom of representation does not deprive evaluation of objective content. It directs inquiry toward what survives admissible reformulations: which futures retain priority, which constraints remain binding, which interventions change choices, and which numerical differences are fixed by independent measurement. Knowledge may first identify an equivalence class and then narrow the uncertainty through further anchoring. Demanding a unique numerical representation at every stage invites an overstatement of the evidence.

Relations and distances differ in a similar way. Knowing that two states count as the same in some respect does not tell us how the cost of departing from that relation should increase. Different scales can record the same zero-error relation while assigning different tolerances to the same number. An organization comparing losses across levels must explain how these scales are calibrated. Otherwise, a ``small change'' in one place may represent serious damage in another.

This has a direct philosophical bearing on social evaluation. Entering health, time, dignity, return, and ecological effects into a single scorecard already makes value choices. Such choices are sometimes necessary, but they must remain explainable and contestable. The precision of a table does not make its weights facts supplied by the world. A common evaluative institution is justified through how it recognizes relevant interests, handles reasons, and permits revision.

Some value conflicts should not be forced into a complete ranking in advance. A person may find irreplaceable goods in two forms of life while lacking sufficient reason to judge one universally superior. Such incomparability need not be irrational. It may reflect plural ends, limited information, or the creative role of choosing itself. That action eventually requires a choice does not mean that a unique maximum was already waiting to be calculated.

\begingroup\interlinepenalty=10000

A more exact account of evaluation acknowledges both comparability and its conditions. Quantities that can appropriately be measured deserve careful measurement; weights that require explanation should be made public; justified disagreement should remain visible. Mathematical representation clarifies argument, while an evaluation's legitimacy still depends on reasons. Precision and judgment support each other without assuming the same task.

\par\endgroup

\chapter{Chapter 14: Optimality, Steady States, and an Unwritten Future}\label{ph-ch-14}

Imagine a person alternating each day between work and rest. Over time, a stable average proportion may emerge, although the person never occupies a state in which change has ceased. The example distinguishes several forms of long-term stability. An average structure, a region visited repeatedly, the distribution of conditional survivors, and a state in which a process finally remains are not the same fact.\footnote{\hypertarget{ph-note-37}{}Core, Chapters 10, 19, and 21; original paper, pp.~6--7 and 13--15, on GOD, GOA, and evolutionary direction. The question of long-run direction is retained while invariance, averaging, cycling, and attraction are distinguished.}

The original paper linked action and purpose through the Global Optimal Direction and the Global Optimal Attractor. The proposal expresses an important question: how can local activities form a longer-term organizational direction? But a direction assumed in advance to be a unique gradient, or a future assumed to end at a unique fixed point, excludes many possibilities. Several actions may be equally good, long-run behavior may cycle, environments may change, and purposes may conflict.

Return to the small device from Chapter 2. One visible bit records whether the operating condition is satisfied; an internal bit is resampled at every step. The internal bit may change or stay the same. If we are concerned only with the visible condition and the consequences of actions, we need not distinguish every internal value. The device has two actions: maintenance makes the operating condition hold at the next step, while interruption makes it fail. Despite its simplicity, the example permits careful distinctions among object description, action selection, and continued existence.

It first offers a positive account of objecthood. Grouping the four complete states into two classes according to the visible bit preserves the consequences of every action exactly. This description has a definite use: it lets us explain what we are operating on with fewer distinctions. Its accuracy does not depend on whether an observer likes the classification, but on whether differences within each class really leave the relevant consequences unchanged. If the content of the internal bit later matters, the description must be enlarged. A common world can admit more than one accurate description of objects because it also admits more than one question.

Let the device start in a state satisfying the operating condition. After every preceding history, the controller chooses interruption with probability one-eighth. After the first step, the proportions satisfying and failing the condition are already stable: seven-eighths and one-eighth, respectively. The same distribution recurs at every subsequent time. Stability here does not even require a long wait.

A different answer follows if the concern is whether interruption has ever occurred since the start. After two steps, the probability of at least one interruption is fifteen-sixty-fourths, exceeding the one-eighth probability at either individual time. Some runs are interrupted and then recover; observing only the endpoint misses their earlier interruption. As the horizon extends, the probability of an interruption at some point approaches one. A distribution can thus be stable from the first step while every run continues to face fresh opportunities for interruption.

An assessment of persistence must specify what it tracks. A marginal distribution answers how many systems are operating at a given time. Whether a task can be completed without interruption requires attention to the whole history. That history preserves differences the marginal cannot capture. Recovery has genuine value because it makes later activity possible again; it does not erase an interruption that has already occurred.

We can also examine only the records with no interruption during a specified interval. In those records, the visible condition always holds, and from the first update onward the internal bit has a stable uniform distribution. This conditional description is accurate too. It tells us how the system is organized given continued survival. But the interrupted records were excluded before the description began. Stability among survivors does not by itself tell us how likely the original controller was to produce survivors.

Conditioning therefore need not weaken explanation. It makes room for two worthwhile questions: the probability of survival and the organization found among survivors. Both can be rigorously derived from the same operating mechanism and can thereby constrain one another. Trouble arises when the condition is forgotten and an answer about organization during continued operation is presented as a guarantee that every original run will continue.

Optimality introduces another dimension. Suppose the scoring rule rewards interruption and gives maintenance no score. Under that explicit rule, always interrupting maximizes the score. If we also require the device to preserve its operating condition from a satisfactory starting state, always maintaining is the permitted choice. Both judgments can be calculated. Their conflict calls for scrutiny of the evaluation and the practical constraint; an assurance that evolution necessarily seeks the optimum cannot settle it. Optimality requires feasible options, a specified evaluation, a time horizon, and a comparison class. Changing any of these may change the conclusion.

Even the scope of ``always maintaining'' needs attention. One restricted form of control requires only that a satisfactory condition must not be destroyed. It says nothing further about a state that has already failed. Under the reward just described, this controller maintains satisfactory states and continues to earn interruption rewards in failed states. Starting from a satisfactory state, the device remains satisfactory forever, while the full state space retains two different long-run regions. The persistence of the first does not require the whole space to converge to a single endpoint.

This distinction suggests a more defensible account of directedness. Actions can acquire a coherent direction around conditions already undertaken. That direction rests on the current organization, the responses it can implement, and the evaluation it recognizes. It need not imply that every possible initial state shares one destiny. Responsibility for a particular history often begins with preserving the relations that it can continue to sustain, rather than making the whole world admit only one outcome.

Stability, feasibility, and optimality each perform useful work. Stability lets us identify order that recurs through change; feasibility lets us ask whether a commitment has an actual path to fulfillment; optimality helps compare actions within stated evaluations and constraints. Placing them in one process allows inquiry into both their mutual support and their conflict. This form of unity adds content to understanding because each concept must confront consequences disclosed by the others.

Why a particular order is worth maintaining requires further reasons. An arrangement can last a long time while imposing indefensible costs on some participants. A scoring rule can be followed accurately while recording none of the losses that matter most. The assessment of a goal cannot stop at whether it is optimized. It must ask who enters the evaluation, which consequences are omitted, and whether those affected can raise objections capable of changing the arrangement.

The relation between meaning and persistence then becomes more concrete. Where a commitment must be fulfilled over time, preserving bodies, memory, cooperation, and routes of repair can give a wish the conditions for realization. These capacities deserve care. Yet a good action that can occur only once, a relationship that has ended, or an undertaking left unfinished by a finite life can also have meaning. The capacity to keep operating is not a necessary condition of every kind of meaning, and the end of a history does not imply that it was never worth having.

A unified natural philosophy can develop by connecting, within a common world, how organization is maintained, how actions become practicable, how purposes receive justification, and how histories bear consequences. Natural processes provide conditions for practice and limit what it can do; human reasons and shared life give some possibilities a significance worth undertaking. Understanding these relations preserves the ambition of a unified account of truth while making every step answerable to a specific question.\footnote{\hypertarget{ph-note-38}{}Sections 21.7 and 35.7 of the mathematical volume give the complete derivation for the same four-state source and its verified parameterized Lean construction. The same policy chooses interruption with probability one-eighth conditional on every history. Its unconditional invariant law, conditional-survival distribution, risk of past failure, and safety-constrained optimum are distinguished. The scope of Lean verification for the new numerical values and the complete constrained example remains as stated in the mathematical text. The philosophical chapter uses these results to discuss persistence and purpose; its normative judgments about defensible costs, objections, and meaning are argued independently.}

Long-run structure especially requires clarity about what is observed. The mean population composition across many replicated worlds, an ancestral path sampled with reproductive-value weighting, and the actual proportions in this finite population today may have different distributions. A weighted genealogy emphasizes paths that contribute more to the future; it does not show that the original controller acted with the reweighted probabilities. A regularity produced by sampling must not be substituted for the system's own rule of action.

Positive results about finite populations require careful quantifiers. First fix a finite observation time; then introduce enough genuinely independent founder cohorts, paying the actual costs of each, to bound the probability that the observed proportions depart from the mean ranking. This cannot show that a fixed small population will necessarily reach the optimum over an infinite future. A faster-growing program may be extinguished by the first random event, and faithful copying cannot recreate a program that has already disappeared.

\begingroup\interlinepenalty=10000

An explicit version of a theory can therefore complete the long-run account it promises while leaving the future open. It establishes which conditions guarantee which conclusions and which boundaries defeat stronger claims. It does not remove possibility, chance, or history from the world.\textsuperscript{\hyperlink{ph-note-3}{3}}

\par\endgroup

\chapter{Chapter 15: Agency and the Revision of Purposes}\label{ph-ch-15}

Executing a given goal effectively and asking whether that goal is appropriate are different capacities. The first makes a system flexible within an assigned task; the second makes the task itself available for reflection. Much of human agency consists in moving between these levels.

Levels of agency need to be distinguished here. Research on minimal agency asks how a system's own organization regulates its interaction with the environment. Feedback alone does not settle the dispute over the distinction between a thermostat and an autonomous agent. Reflective agency additionally requires the ability to examine reasons for action, programs, or goals themselves. Barandiaran, Di Paolo, and Rohde analyze agency through individuality, interactional asymmetry, and normativity, showing how the question can arise before explicit reflection. This chapter's discussion of goal revision concerns the stronger capacity. Lacking it does not entail the absence of every form of agency or every kind of moral standing.\footnote{\hypertarget{ph-note-39}{}Barandiaran, Di Paolo, and Rohde, ``Defining Agency: Individuality, Normativity, Asymmetry, and Spatio-temporality in Action'' (2009), pp.~367--386. This edition checked the formal publication metadata and the abstract and introduction of the authors' manuscript. These materials introduce the problem of minimal agency; no claim is made that a UEOT program satisfies all of the authors' criteria.}

Self-revision does not require an agent to stand outside her entire history. We examine reasons through a language already learned, imagine alternative lives with capacities already acquired, and learn within social relations what is worth questioning. Reflection has conditions, but it can change how those conditions shape what comes next. No one chose every circumstance of birth; a person may nevertheless learn to reject demands that those circumstances presented as self-evident.

An entirely unified supreme controller is a poor model of the self's internal conflicts. Habits, emotions, time horizons, and social commitments can make demands at once. Mature agency may organize ways of dealing with these differences rather than eliminate them: delaying a response, redescribing a problem, bringing a longer commitment to bear on a present choice, or allowing a new reason to revise an old commitment. Unity of self is an achievement of coordination that must be renewed.

Such coordination cannot be evaluated solely by asking which part controls the whole. A narrow aim may organize every activity very effectively while depriving a life of other important capacities. Someone who devotes everything to winning a ranking may be highly consistent while becoming unable to take friendship or bodily needs seriously. A richer agency includes the capacity to identify a misspecified goal and recover excluded reasons.\footnote{\hypertarget{ph-note-40}{}Core, Chapter 27, on reflexive structure, goal revision, and interfaces of agency; original paper, pp.~15--18 and 24--28. The main text gives the reasons for its claims about the normative significance of personal reflection and education.}

Knowing how to achieve a goal likewise differs from knowing why to achieve it. Technical proficiency can enlarge the space of action, but it can also strengthen a mistaken aim. Education should therefore develop more than execution. It should enable people to identify the sources of information, evaluation, and power, explain their choices, and attend to the responses of those affected.

An objection is that if every goal is revisable, action can never begin. Actual reflection does not require a fresh judgment on all values before every act. Most of the time we rely on relatively stable habits and commitments; conflict, failure, or important new evidence brings part of that background into the foreground. Rationality is layered in this respect: it includes both skillful continuation and justified interruption.

The reflective agency examined in this chapter is a complex, revisable relation between an object and the world. An agent forms predictions, undertakes certain futures, observes the consequences of acting, and may change representations, programs, or goals. This occurs within bodily and social conditions. Because the agent can reorganize how these conditions shape the future, the presence of causation does not make the capacity unreal.

A beneficial program's increasing share through inheritance does not mean that it has learned to reflect. Natural selection changes which organizations become more common; learning changes how an organization uses experience to adjust its actions; revision in light of reasons also involves understanding goals, acknowledging error, and responding to others. These processes may interact, but the growth proof within the common model directly supports only the connection between the first kind and the permitted controls.

\begingroup\interlinepenalty=10000

An agent who forgoes some opportunity for growth in order to care for someone else is not thereby a ``mathematically mistaken object.'' Whether a mistake has been made depends on the purposes and commitments adopted and the reasons that can be given. Mechanistic theory helps clarify costs and consequences; normative inquiry asks which consequences are worth bearing. Keeping their responsibilities clear makes the revision of purposes a real problem, rather than an exercise in labeling every choice optimal after the event.

\par\endgroup

\part*{\phantomsection\hypertarget{ph-part-4}{}\addcontentsline{toc}{part}{Part IV: Composition, Scale, and Evolution}{\Large Part IV}\\[1.2em]{\LARGE Composition, Scale, and Evolution}}

\chapter{Chapter 16: When Many Individuals Form a Whole}\label{ph-ch-16}

Putting several people in one room does not yet create a cooperative organization. If they begin to divide tasks, exchange necessary information, take up one another's actions, and sustain that coordination, a further question arises: does the new whole have capacities that cannot be read directly from their isolated activities? The reality of the whole must be explained through the consequences of this shared activity.

The simplest connection among parts is a common environment. Rain makes many people open umbrellas, but it does not make them an organization collectively controlling the weather. Interaction supplies a stronger connection: one person's action changes another's conditions. Participants can go further and jointly realize a durable interface, perhaps through mutual checking, common decision-making, or the reliable production of an output. Moving from common change to common organization requires additional mechanisms.\footnote{\hypertarget{ph-note-41}{}Core, Chapter 15, on commonality, joint information, formation at a higher organizational level, and the composition window; mathematical volume, Chapters 32 and 34. The two-card example is a thought experiment constructed to explain joint information.}

Shared information may also exceed what the parts supply separately. Imagine a sign encoded across two cards. Neither card alone identifies the sign, yet comparing whether the cards match reveals it. The additional readable information neither violates nature nor appears mysteriously from nothing. It belongs to the joint relation. Measuring only each part's isolated performance would miss the capacity.

An account of the whole can thus have real content without invoking an extra entity detached from the parts. Its activity is realized through the parts and their relations, but an inventory of their independent properties may be insufficient to predict the joint process. An orchestra performs through particular people and instruments; the organization of its performance is not adequately explained by the fact that each musician can make sounds. Rhythm, division of labor, and mutual adjustment belong to the causal explanation.

A sharper objection remains. If every step is realized by the activities of the parts, does saying that the whole acts count the same effect twice? The account requires no second force. Changing traffic rules changes the coordination of vehicles and pedestrians, but the change must be realized through signs, learning, enforcement, and actual driving. An explanation at the level of the whole is useful if it identifies how the organizational difference changes the result and remains explanatory across relevant changes of material realization. If a proposed variable for the whole has no mode of realization and adds nothing to prediction or intervention, its larger scale is no reason to retain it.

A joint capacity does not automatically give the whole a lasting identity. A temporary collaboration may complete a task and dissolve; an enduring organization must also maintain rules, membership, resources, and historical continuity. Producing one shared output, establishing a repeatable mechanism, and forming an object at a higher organizational level are distinct stages that may connect. At each stage, we should ask what testable organizational content has been added.

This distinction is particularly important in human life. A group cannot cancel its members' independent interests simply by declaring itself a whole. Even where common institutions are real, questions remain about which capacities must be centralized, which decisions should be distributed, and how members retain opportunities to leave, speak, and revise the arrangement. Wholeness and legitimacy require separate arguments that must then be connected.

The best forms of composition often create new common capacities while preserving sufficient individuality. Without constraint, cooperation may not last; without difference, an organization may lose sources of learning and correction. Composition can be understood as a window that requires maintenance: connections must be strong enough to support common tasks, while differences must remain sufficient for participants to continue as capable agents. The location of that window depends on the task, its risks, and public reasons.

A whole is accordingly a real achievement. It depends on local realization and can change local futures through the rules and interfaces it has established. In cases we can explain, so-called top-down effects are the operation of these concrete constraints. Showing how they are realized gives a stronger account of the whole than assigning it a mysterious extra force.

\begingroup\interlinepenalty=10000

Backup relations make the point especially clear. One additional healthy individual does not automatically provide one additional reliable capacity. If a fault affects only the body dispatched for the present action, an undispatched offspring may still perform the task. If a single environmental shock damages all existing individuals, the same apparent headcount offers different redundancy. Explaining the common capacity requires an account of fault correlations and of which parts the operator can actually access. Calibration confined to an isolated root object cannot validate a decision to deploy its backup offspring.\textsuperscript{\hyperlink{ph-note-27}{27}}

\par\endgroup

\chapter{Chapter 17: What Changes with Scale}\label{ph-ch-17}

On a city map, streets become lines, buildings become blocks, and the residents' experiences nearly disappear. The map is not thereby a failed photograph of the city. It retains relations important to one kind of task and leaves out others. A change of scale can yield new understanding, but it can also conceal distinctions that determine whether an action succeeds.

For a coarse description to become a reliable effective description, it must answer a question: can the situations it groups together still be treated alike in the subsequent process that matters? Two roads that look alike cannot be unconditionally merged for navigation if one is open and the other closed. Legitimate compression depends on responses and tasks, not just resemblance.\footnote{\hypertarget{ph-note-42}{}Core, Chapters 7, 14--15, and 20; mathematical volume, Chapter 32, on eight kinds of level relation. Predictive quotients, control quotients, composition, history, selection, and long-run structure each have conditions of their own.}

Even a compression that preserves prediction may fail to preserve control. Two microscopic states may give the same distribution for one observation while allowing different actions or incurring different costs. A coarse description used to organize control must account for these differences too. Otherwise, the macroscopic plan may demand actions the microscopic situation cannot implement.

Several kinds of ``level'' are also easily confused. A city's composition from districts is a part--whole relation. Replacing a detailed map with a sketch changes resolution. Tracing the city's name and institutions through time concerns historical continuity. The movement of researchers between cities concerns lineages or transmission. Comparing the cities' public goals involves evaluation. Each concerns relations across levels, but they are not different names for one operation. A unified account must explain when these relations are compatible rather than combine them in advance into an all-purpose hierarchy.

Scale must also be distinguished from time. Shrinking a map does not make the city live through its future. Slow variables may sometimes describe a long-term process, but the relation between fast and slow activities needs to be established. Omitting a state at one scale does not mean it has ceased to exist in actual time. Conversely, a change in organization over time is more than replacing a map with a coarser one.

Actual systems usually admit approximate descriptions. The task is then to specify where errors enter, how they propagate, and how much the task can tolerate. A small initial error cannot be disregarded if it is repeatedly amplified. Transferring an error budget across scales or units also requires calibration. Agreement of the underlying structure may establish that certain zero-difference relations coincide; it does not determine the magnitude of every nonzero difference.\footnote{\hypertarget{ph-note-43}{}Mathematical volume, Chapter 33, on relational defects, two-sided calibration, and transport of budgets. The discussion of different scales of measurement is not used as a proof that a physical renormalization flow has been established.}

On this account, macroscopic explanation has a conditional autonomy. Within an established task domain, we may use coarse states without repeatedly restating every microscopic detail. Their autonomy rests on reliable compatibility, not separation from the microscopic world. Biology, psychology, and the social sciences can retain effective concepts of their own, while their use across levels requires an account of the relevant interfaces.

The most valuable achievement of multiscale unification is therefore to identify what is preserved, what information is lost, and what additional mechanisms make macroscopic regularities possible. Each successful connection clarifies part of the common world. It need not place every level in a cosmic sequence already declared complete.

A legitimate simplification must say what it preserves. In the common model used in this edition, a physical label that is neither read by the sensor nor used in any transaction can be merged away. Under the same noise, repair, payment, survival, birth, and permitted observations remain identical item by item. Bodies with different labels remain distinct realizations, while the control problem for this purpose is preserved. Actual individual identities and parent pointers remain in the history; they do not disappear because finite types have been grouped together.

By contrast, retaining only a decoded program can omit whether the main execution gate is damaged. Two bodies may agree in their decoded program while only one can execute it. Removing damage diagnosis can similarly merge an action affordable for a healthy body with one unaffordable for a damaged body. Statistical accuracy for one readout does not establish suitability for questions of repair, resources, or inheritance.

\begingroup\interlinepenalty=10000

This offers an answer to arbitrary relativism. Representations can vary, but translation between them must preserve the relevant operations, information, and consequences. Some transformations meet those demands and others do not; the comparison is grounded in actual relations rather than a describer's preference. It also challenges a rigid doctrine of one unique hierarchy. The same process can admit several effective representations, each performing a different task within a stated use.\textsuperscript{\hyperlink{ph-note-3}{3}}

\par\endgroup

\chapter{Chapter 18: Coordination and Conflict in Shared Life}\label{ph-ch-18}

Central direction can make an organization more efficient, and it can suppress local knowledge. Members pursuing their own aims may improve collective performance or transfer costs to others. The relation between a whole and its parts therefore cannot be understood through obedience or noninterference alone. It calls for inquiry into how goals, information, power, and consequences are organized.

The suggestion that parts should act along the whole's optimal direction is particularly liable to mislead. It conceals at least two questions: who supplies the whole's evaluation, and do the parts' changes actually improve it? A direction that improves one local measure need not, when combined with every other local improvement, improve a global objective. Different demands may cancel one another or form a cycle.

Imagine three institutional arrangements, each supported by some participants and each favored over another under a local comparison. Discussion may move from the first to the second, then to the third, and finally back to the first. If such a cycle exists, we cannot assign each arrangement a fixed total score while maintaining that every local revision strictly raises it. We must change the coordinating mechanism, specify the comparison conditions, or acknowledge an unresolved conflict of values.\footnote{\hypertarget{ph-note-44}{}Core, Chapter 24, on multiple objectives, global potentials, evaluation at a higher organizational level, and thresholds for coordination; original paper, pp.~8 and 18--28. The three-arrangement cycle is a conceptual example showing why local improvement requires a further connection to global evaluation.}

Coordination must consequently be built; it is not the destination already contained in every interaction. Sharing information can reduce misunderstanding; compensation can redistribute costs; commitment can extend the time horizon; public discussion can revise goals. Each has conditions. Information may be reported selectively, compensation may fail to cover irreversible loss, and commitment may lock participants into a mistaken institution. The quality of coordination must be assessed together with these conditions.

Organizations spanning several levels also need to accommodate different time scales. Some local feedback requires an immediate response, while longer-term rules establish stable expectations. Requiring every local action to await the highest authority can destroy responsiveness; letting every short-term signal overturn a long-term constraint can make common commitments impossible. An appropriate arrangement does more than average the fast and the slow. It specifies which decisions belong to whom and over what horizon.

An important ethical consequence follows: macroscopic persistence does not automatically give members a reason to obey. An institution can maintain itself by persistently transferring losses to disadvantaged participants. Observing only whether the whole survives removes their predicament from the evaluation. A community's legitimacy depends on how it treats its members and whether they can help decide what is to be maintained.

This does not reduce every common obligation to immediate voluntary agreement. Shared resources, intergenerational responsibilities, and mutual dependence can extend obligations beyond a single act of consent. Yet the harder a relationship is to leave, the stronger the need for public reasons, appropriate representation, and channels for correction. A whole's real existence makes responsibility more pressing; it does not settle it.

\begingroup\interlinepenalty=10000

A unified philosophy should therefore understand coordination as a revisable arrangement for maintaining common capacities. Conflict need not be an eternal natural destiny, and harmony is not an inevitable outcome requiring no mechanism. The work of shared life is to find organizational forms that can sustain both difference and dependence.

\par\endgroup

\chapter{Chapter 19: Reproduction, Inheritance, and Individuality within a Lineage}\label{ph-ch-19}

One life continues through repair; another begins through reproduction. Both transmit organization, but their relations differ. Repair aims to preserve the continuity of this life, while reproduction brings a new individual into an inheritance relation. Calling both the immortality of one and the same object conceals the formation and ending of individual lives.

A lineage has a reality of its own. Structures, capacities, and modes of development enter descendants through reproductive mechanisms, are changed by their activities, and pass onward again. What is tracked is a relation of inheritance across individuals. It may endure for a long time without making any two individuals the same life. Parents and children can profoundly continue one another while each has a life of their own.\footnote{\hypertarget{ph-note-45}{}Core, Chapters 25--26; mathematical volume, Chapter 34, on a genealogical bridge derived from the same reproductive operator. Mean growth, actual extinction, conditional survival, and individual identity are treated separately.}

To understand inheritance, we must distinguish what is copied, how the copying is realized, and which consequences contribute to retention. The same abstract pattern may be realized in different carriers, but its transmission depends on concrete mechanisms. An offspring equivalent to its parent for one predictive task need not inherit every function. A widely transmitted pattern need not improve the welfare of everyone who carries it. Success in transmission and justified value are not interchangeable.

The unit of selection also requires a mechanistic account. Some changes affect an individual's chances of leaving descendants; others may matter only through group organization or longer-term transmission. Asking which level is selected requires examining how differences enter reproduction. The answer does not follow from whatever we draw as a whole. Object scale, organizational level, and level of selection can be related without coinciding.

Population averages also differ from particular histories. A process that grows on average can become extinct in some realizations. Its mean composition may also differ from the composition observed conditional on survival. These are not merely technical details. They determine to whom we are promising continuation and who bears the corresponding risk.

The original paper used evolutionary lineages to connect the brevity of life with the persistence of organization. That connection can be retained without treating evolution as an inevitable ascent toward a fixed optimal destination. Competitors, resources, and mutual dependence can change the conditions of selection. An organization's present advantage does not reveal a form eternally preferred by nature. Lineage history includes accumulation, loss, branching, and paths that happened to remain.

For human self-understanding, an awareness of lineage has two values. It reveals our dependence on predecessors we never met, and it warns against turning inheritance into possession of those who follow. Transmitting language, knowledge, and institutions is better when it also transmits the capacity to examine and revise them. A rich inheritance supplies conditions for continued life while leaving successors room to become different.

Actual birth histories constrain how models index objects. If every birth enters a strictly later generation, a finite permanent table of individuals cannot give every individual a normalized offspring channel: the individual of greatest generation has no later-generation offspring available within the table. This strict obstruction preserves the problem and requires a distinction between repeatable types and unique historical individuals. An average reproductive formula over types can have an actual birth process as its source. It still cannot, by itself, provide each offspring with a global identity or establish the existence of an infinite genealogy.\footnote{\hypertarget{ph-note-46}{}Mathematical volume, Section 39.2; \nolinkurl{RegisterRepairHistory.no_finite_normalized_strict_birth} and \nolinkurl{actual_birth_price_bridge}. A label history is not a globally unique identity for concurrent individuals across time.}

Labels may recur, while the time of an event supplies further historical evidence. If the transaction counter advances through success, faults, and failures, two offspring born to the same parent at different times can be distinguished by parent identity, actual count, and label, even when their inherited labels match. Their identities arise from events that occurred. Erasing the time metadata leaves the original distributions of actions, failures, and costs unchanged. The result still depends on a reliable and accessible serial record and a specified initial parent; it does not solve identity assignment for every individual in an arbitrarily concurrent population.\footnote{\hypertarget{ph-note-47}{}Mathematical volume, Section 39.8, and \nolinkurl{RegisterTimestampedBirthSource}. Serial event identity for the same retained parent and the interpretation of parent identity at a successful birth differ from the problem of global individual identity in an arbitrary concurrent population.}

A further model permits several parents to act in sequence. It derives parent identity from records of actual lookup, repair, and newborn locations, with failures also advancing the clock. Copying a protected program and drawing a random birth label are recorded separately. The Price identity, used to decompose change between generations, is applied only to the previously fixed queue of old parents and their actual offspring. Empty generations remain in the record, and this offspring population must not be mistaken for the entire population present at the end of the run.\footnote{\hypertarget{ph-note-48}{}Mathematical volume, Sections 39.20--39.21, on serial population identity and the Price identity for actual offspring. Empty generations and the original P5 program-matching conditions remain in place.}

We must also know whose program each child inherits. In a different environment, where the fault target is determined by each object's own protected program, actual offspring of two types can enter their respective maintenance processes. This result leaves intact the limitation of the original model with its fixed global target. Execution records and a separately specified noisy readout retain all faults and erroneous readings; only on that basis can the risk of identifying the actual parent be calculated. Only when the readout process satisfies its correctness event and a birth has actually occurred does the very offspring returned by the procedure acquire the relevant maintenance conclusion. If arbitrary unknown confounder classes are permitted and the noise parameter is one-half, different parents can still yield exactly the same readout distribution. Addresses, logs, and permission to read them are conditions of these results. They neither allow ancestry to be guessed from appearance nor guarantee optimality or moral legitimacy.\footnote{\hypertarget{ph-note-49}{}Mathematical volume, Section 39.25, \nolinkurl{RegisterOwnProgramJointJournal}. The fault target determined by an object's own program and permission to read the marked journal are specified separately. \texttt{None} does not identify the party that attempted a failed transaction.}

A parent's continued presence at the next census and its production of a new individual must be recorded separately. The first continues an identity; the second establishes a new parent--offspring relation. Adding the two can support a study of population contribution, but it does not turn every contribution edge into a strict birth. Confusing the relations changes organizational lifespan, generation length, and ancestral paths.

The common model in this edition accordingly retains a unique parent pointer for each actual birth and also records the edge along which the parent continues to act. Weighted genealogical path relations are then proved from these records, rather than inferred backward from a mean matrix. What descendants copy is the actual repaired program, after which they undergo the stated damage process. Faithfulness follows from the particular copying mechanism and majority-code recovery; it need not hold after complete erasure or under a different mutation rule.

\begingroup\interlinepenalty=10000

Selection feedback has a positive scope that can be stated precisely. Two actually different programs each face nonzero damage and birth noise under the same supply and clock. The block with greater growth increases its normalized expected share. Sufficiently many independent founder cohorts then connect this ordering to actual proportions at a finite time. Yet faithful copying cannot introduce a program absent at the outset, and a finite population may lose the superior program first. Selection has material and historical conditions. It supplies no blank guarantee of innovation or inevitable progress.\textsuperscript{\hyperlink{ph-note-3}{3}}

\par\endgroup

\chapter{Chapter 20: Innovation, Fragility, and Open Evolution}\label{ph-ch-20}

An organization increasingly effective at a familiar task may become less able to cope when the task itself changes. Skill saves resources, while exploration can interrupt skill, expose uncertainty, and temporarily reduce performance. The original paper's distinction between exploration and optimization provides a useful starting point for thinking about innovation.\footnote{\hypertarget{ph-note-50}{}Original paper, pp.~12--15. Core, Chapter 26, already distinguishes transitions in individuality, mechanisms of selection, and unconditional evolutionary optimization. This chapter develops the question of how spaces of possibility arise.}

Some innovations find a better option among existing choices; others change what can be chosen. A new way of cooperating creates a common capacity that previously did not exist. A new representation makes formerly indistinguishable situations available as different tasks. Portraying every innovation as an ascent over a fixed landscape treats the formation of the space of choices as though it had already been completed.

Two levels of change therefore require explanation: how performance improves under fixed conditions, and how action transforms the conditions themselves. The first asks which action is better; the second also asks where new actions, relations, and evaluations come from. They may alternate within a history, but no argument requires every system to follow the same life cycle in the same sequence.

Innovation is not automatically worthy of praise. A new mechanism can enlarge cooperation or increase the power to injure; a stronger adaptation can transfer costs to the surrounding environment. Novelty and value are different judgments. Evaluating innovation solely by its continued spread mistakes successful expansion for sufficient justification.

Fragility, meanwhile, warns that long-term success can conceal concentrated dependencies. Many apparently independent activities may rely on one channel, allowing a local interruption to have extensive effects. Conversely, redundancy and diversity that seem repetitive in the present task may preserve opportunities for future recovery. Whether they should be maintained depends on cost, fault structure, and who bears the losses. More is not invariably better.

An open account of evolution must accommodate both reversible adjustment and irreversible transformation. Restoring an existing organization, forming a new whole, producing offspring, and complete dissolution have different historical meanings. An original object may disappear while some of its materials and capacities enter another organization. The relation of continuation is real, but it is not necessarily a victory for the original individual. Taking endings seriously allows us to ask what is worth passing on.

Human culture gives this openness a reflective dimension. We can recognize that we are changing the conditions of choice, and thereby discuss the direction of innovation, the scale of experiments, and repair after failure. Such discussion cannot turn history into a wholly designable project. It can, however, reduce the temptation to portray every actual victor as inevitable progress.

Hope in evolution then has a more concrete meaning. It is confidence that possibilities for improvement remain, together with a willingness to build the conditions for finding and realizing them. Hope does not require the world to guarantee ultimate success. It requires us to take seriously our capacity to affect what follows.

Copying an existing correct program differs from creating a new one. Majority recovery depends on coding redundancy that remains available. When two different histories are erased into exactly the same physical state, the recovery procedure has no information by which to choose the target it should restore. Adding an action called ``create the correct answer'' merely moves the unexplained problem into the definition of the action. This impossibility result protects the problem of innovation itself: an explanation must identify where new differences come from, how they are evaluated, and how they enter actual transmission. A repair theorem cannot simply be extended without limit.

\begingroup\interlinepenalty=10000

Environmental change can also change rankings. A program favored under a fixed supply may lose its advantage when the supply stops; a lineage with faster average growth may be less robust under a worst-case risk criterion. A core theoretical edition can state these conditions and failure cases without first exhausting every ecology and environment. Making its boundaries public gives subsequent inquiry a genuine starting point and prevents every counterexample from being redescribed as something the theory predicted all along.

\par\endgroup

\part*{\phantomsection\hypertarget{ph-part-5}{}\addcontentsline{toc}{part}{Part V: Mind, Freedom, and Shared Life}{\Large Part V}\\[1.2em]{\LARGE Mind, Freedom, and Shared Life}}

\chapter{Chapter 21: Why Experience Has a Subject}\label{ph-ch-21}

Seeing a color, feeling pain, and remembering an encounter are more than descriptions of external events. They are happening for someone. We can observe a person's activities from the outside and listen to their reports, but both forms of access differ from the first-person experience they are undergoing. That difference calls for serious explanation.

This book begins from a naturalistic position: subjects are embodied beings within the world, with histories and relationships. Experience is connected with the organization of perception, memory, action, and bodily regulation. A subject does not first exist, fully formed, inside the mind and then collect images of an external world through channels of access. Capacities, perspectives, and self-understanding develop through engagement with the world. This starting point places mind within a shared nature rather than outside the causal world.\footnote{\hypertarget{ph-note-51}{}Original paper, pp.~15--18; Core, Section 27.6, distinguishing functional, intentional, and normative significance. The naturalism about subjects defended here is a philosophical position. A model of neural organization is not equated with a complete explanation of phenomenal consciousness.}

Considered in organizational terms, experience has a certain unity. Sounds, sensations, intentions, and memories need not remain unrelated fragments; they can jointly affect an action. Yet this unity does not require a central observer to survey every content at every moment. Much activity takes place outside attention, and some conflicts become apparent only when action is interrupted. The unity of a subject may consist in an integration continually achieved among multiple processes.

Further argument is nevertheless needed to move from the integration of functions to an explanation of why there is experience. A model might accurately describe how information travels, which contents can be reported, and how behavior changes, without thereby answering why suffering is experienced as suffering. Acknowledging this problem preserves the explanatory responsibility of naturalism. Renaming unexplained subjectivity a ``complex network'' does not complete the explanation.

This is also the challenge raised by Chalmers in distinguishing questions about functions from the question of experience. Even after discrimination, reporting, and behavioral control have been explained, we can still ask why these activities are accompanied by experience. Recognizing the question does not require accepting his entire theory of consciousness. It requires this book to say how far its own explanation reaches. A functionalist may, in turn, maintain that a sufficiently complete organizational account will eventually dissolve the apparent additional requirement. This book has not settled the dispute: no theorem about organization supplied here establishes sufficient conditions for experience, and no argument here proves that such a sufficient explanation is impossible in principle.\footnote{\hypertarget{ph-note-52}{}David J. Chalmers, ``Facing Up to the Problem of Consciousness,'' \emph{Journal of Consciousness Studies} 2(3) (1995): 200--219, especially Sections 2--3. The author's publicly available text was checked. This book takes up its explanatory challenge without treating it as proof that naturalism has been refuted.}

One objection might ask whether consciousness should be placed outside nature if organizational explanation has not resolved every question about experience. That conclusion would also be premature. An incomplete explanation does not by itself provide grounds for positing another kind of entity. We can seek stronger connections among bodies, processes, reports, and experience, test competing organizational hypotheses, and retain first-person evidence as a constraint on inquiry. The strength of naturalism lies in allowing these kinds of evidence to address a shared problem, rather than in denying the problem in advance.

Nor is the difficulty resolved by calling every object with feedback conscious. A device may maintain a variable, and a lineage may preserve information, without thereby possessing what we mean by experience. Objecthood, functional agency, reflective selfhood, and phenomenal experience must be connected by specific arguments. Continuities in the conditions under which concepts become applicable do not erase every distinction among them.

Subjectivity also involves a situation that no one else can simply replace. My body is injured here; my memory connects this encounter with a particular history; my actions change the world from this position. Others can understand, help, and explain me, but incorporating my experiences into a comprehensive model does not make me dispensable. A first-person standpoint has a place in a shared world precisely because that world contains different beings who actually undergo its consequences.

\begingroup\interlinepenalty=10000

This chapter proposes a naturalistic program of explanation that still needs to be developed: investigate the bodily conditions of experience, first-person evidence, and the consequences of action together, allowing each to constrain the others. It has not proved that physical description exhausts experience, nor inferred another kind of entity from an explanatory gap. If the unity of nature is to include subjects, an account must explain their real place within it. It cannot begin by removing the experience it is meant to explain.

\par\endgroup

\chapter{Chapter 22: Attention, Emotion, and a World That Matters}\label{ph-ch-22}

If every change reaching the senses mattered equally, action would be difficult to organize. Attention brings some contents into the foreground and leaves others in the background. Emotion gives the world different kinds of weight: urgency, warmth, danger, loss, or promise. Attention and emotion help constitute the world an agent actually encounters; they do more than add color to reasoning already completed.

The selectivity of attention depends on history. In the same scene, an experienced practitioner may notice crucial clues that a beginner has yet to distinguish. Caring about someone may make us more attentive to their needs than we were before. Goals, memories, and feelings jointly shape salience, while changes in salience can influence new goals. This reciprocal relation makes attention part of agency.

Emotion can also have an epistemic role. Worry brings a loss that has not yet occurred into consideration; shame makes our conduct as seen by others a question for us; anger may disclose an injustice that has been overlooked. Yet emotion can also misidentify its object, exaggerate a threat, or remain bound to an earlier experience. Taking emotion seriously requires neither treating it as an absolute authority nor excluding it before rational thought can begin. Reason can learn what is at stake through emotion, while new information and reasons can revise the interpretation toward which an emotion directs us.

We should distinguish regulation from suppression. A person's ability to delay an immediate response and allow other reasons to enter a choice can enlarge freedom. Learning to remain silent about every discomfort, however, may weaken the ability to recognize injury. The value of regulation depends on what it preserves, not simply on whether outward behavior remains composed. A calm expression can accompany considerable autonomy, or a situation in which expression is forbidden.

Self-understanding can assist this regulation. An agent may discover that a recurring response comes from an earlier experience: it once worked, but may no longer suit the present circumstances. Understanding this does not immediately make the response disappear, but it can create conditions for change. New habits, changes to one's surroundings, and support from others can also help form different responses. The work of the will often consists in organizing these conditions, rather than in overpowering the body through an unconditionally strong inner command.

Empathy extends this understanding to others. Through their expressions and circumstances, and through comparable experiences of our own, we try to understand the world from another perspective. Imagining what someone feels is not the same as knowing it. Their corrections matter. Understanding requires channels through which objections can be heard, especially when we have become highly confident in our own model of the other person.

The original paper places attention, emotion, memory, will, and empathy within a dynamic organization. This direction deserves further development. Specific neural mechanisms require appropriate evidence. At the conceptual level, common organizational patterns can already help frame questions: what is integrated, which feedback is controlled, what information is ignored, and which goals can be revised?\footnote{\hypertarget{ph-note-53}{}Original paper, pp.~15--18, on attention, emotion, memory, and empathy. This text does not adopt neural quantities, clinical diagnostic measures, or specific claims of therapeutic efficacy that have not been independently checked.}

\begingroup\interlinepenalty=10000

A world that matters to someone is not a counterfeit version of the objective world. It is how a finite subject enters reality from a particular situation and with particular concerns. The question is whether that way of engaging remains open to correction by the world and by other people. Intellectual maturity includes making the organization of one's attention more responsive, rather than merely intensifying what already commands it.

\par\endgroup

\chapter{Chapter 23: Freedom Is Not the Absence of Conditions}\label{ph-ch-23}

A person saying ``I agree'' under threat and a person saying the same words after understanding and discussion perform actions with different significance. Both utterances occur within natural causal processes, yet their degrees of freedom plainly differ. An explanation of freedom must therefore ask how causes operate through an agent, as well as whether the behavior has causes at all.

This book understands freedom as a real capacity that depends on conditions: an agent can understand the relevant situation, form and revise intentions, act for reasons among genuinely feasible alternatives, and have appropriate opportunities to resist coercion. Such freedom can develop gradually and can also be weakened. It needs no starting point outside body and history, but it does require that the agent's contribution to the process remain real.

One objection holds that an action cannot be free if it is influenced by character, environment, or past experience. On that requirement, understanding and learning would become enemies of freedom: the more effectively reasons influenced a person, the more their actions would be shaped by causes. A more defensible distinction separates influences that develop the capacity to judge and choose from those that bypass or damage it. Persuasion and manipulation do not differ because one is caused and the other uncaused. They differ in the place they give to understanding, correction, and refusal.

A stronger objection does not depend on present coercion. Imagine someone whose judgments are now coherent and who sincerely endorses their goals, but whose upbringing was deliberately closed to alternatives. Competing reasons were excluded, questioning was repeatedly punished, and the person never had an opportunity to examine how these goals had been formed. Saying that action ``passed through their reasons'' is insufficient to establish how free they are. Their present self-endorsement may itself be the product of manipulation.

An answer must connect present control with its formative history. Access to different reasons, opportunities to question without punishment, and the gradual revision of inculcated goals are important ways of regaining autonomy in action. Not every influence one did not choose cancels freedom. What matters is whether an influence continues to close off understanding and revision, and whether the agent subsequently has opportunities to make a goal into a commitment for which they can take responsibility. These conditions do not constitute a checklist that mechanically settles every case. Fischer and Ravizza's discussions of responsiveness to reasons and of an agent's ownership of the mechanism of control provide an important point of reference for this historical question. This book adopts a limited compatibilist position here; it has not disposed of the objection that freedom requires being the ultimate source of oneself.\footnote{\hypertarget{ph-note-54}{}John Martin Fischer and Mark Ravizza, \emph{Responsibility and Control: A Theory of Moral Responsibility} (Cambridge University Press, 1998), Chapter 7, ``Responsibility and History.'' The publisher's book description and chapter abstract were checked; no page-by-page reading of the whole book is claimed. The reference locates the issues of reasons-responsiveness and historical ownership. The specific manipulation example and this book's response are given in the main text.}

Another objection worries that, if freedom comes in degrees, people can be ranked by a single score of capacity. That inference does not follow. Assessing control over a particular action should help allocate support and responsibility appropriately, rather than decide who has basic standing as a person. Someone with limited capacity for a task still undergoes the consequences of their life. Respect is not conditional on attaining the highest complexity of control.

Judgments of responsibility must consider knowledge, control, opportunities for choice, and reasonably foreseeable consequences. A person unable to obtain crucial information should not be assessed in the same way as someone who knows the consequences and conceals the information. An institution that imposes an impossible task cannot shed its own responsibility for failure merely by saying that a participant ``made a choice.'' The effective scope of freedom is constituted jointly by personal capacities and public conditions.\footnote{\hypertarget{ph-note-55}{}Original paper, pp.~24--28; Core, Chapters 18 and 27, on implementable policies and components of agency. The public reasons concerning freedom and responsibility belong to this book's independent normative argument.}

This also connects reflection with the possibility of action. A person may imagine many options, but if each is blocked by the same coercion, imagination alone does not amount to adequate freedom. Conversely, many formally available options may yield little more than random movement if their differences cannot be understood. Freedom requires epistemic capacities to meet feasible opportunities.

Personal freedom and shared rules need not be opposed. Reliable commitments, public safety, rights of expression, and predictable procedures can extend the plans each person can carry out. Rules can also be used to dominate. What matters is how rules are made, how their effects are distributed, and whether those subject to them can raise objections that have consequences.

\begingroup\interlinepenalty=10000

Naturalism can therefore take freedom seriously. An object capable of changing the future through reasons, memory, and shared institutions possesses real causal powers. Developing those powers is a practical task within the world. Freedom does not require the world to cease affecting us. It consists in gradually acquiring, through our relations with the world, a position from which we can understand, take responsibility for, and revise our actions.

\par\endgroup

\chapter{Chapter 24: The Commitments That Give Life Meaning}\label{ph-ch-24}

When someone says that life has lost its meaning, they are usually reporting more than a decline in efficiency. They may still complete their tasks without knowing why those tasks are worth completing. Or they may sincerely care about something while circumstances prevent them from taking part in it. The question of meaning concerns the relations among subjects, action, and value.

This book understands meaning as a relation in which agents can see their activities as participation in something worth caring about, and can give that participation a place in the history of their lives. This involves felt engagement, an actual connection, and value that is open to discussion. The three can support one another, but they can also come apart. A person who feels fulfilled can still ask whether their chosen cause harms others. Work with real value may fail to provide a sense of wholeness to someone who performs it under adverse conditions.

Connecting engagement with value in this way is not distinctive to UEOT. Susan Wolf likewise understands meaning through a relation between subjective engagement and value independent of personal preference. This book continues that inquiry by asking what bodily capacities, memories, responses from others, and conditions of repair enable meaningful participation to become part of an actual life. Citing this related view does not mean that ``worthwhile'' has already been fully defined. Particular reasons are still needed: reducing avoidable suffering, honestly advancing understanding, or sustaining relationships that do not depend on domination, for example. Such reasons can be challenged and can conflict with one another.\footnote{\hypertarget{ph-note-56}{}Susan Wolf, \emph{Meaning in Life and Why It Matters} (Princeton University Press, 2010). The author's Tanner Lectures under the same title, delivered at Princeton on 7--8 November 2007, were also checked, especially the first lecture's connection between engagement and value and the second's discussion of the limitations of evaluative judgment. This book does not equate a meaningful life with moral faultlessness, nor claim that the source supplies a complete theory of value.}

The immediate subject here is meaning as experienced by agents able to reflect on their own lives. This is not a test of whose life has value. Infants, people who have lost the ability to express themselves, and those temporarily unable to organize their experience into a life narrative do not thereby lose their standing as beings worthy of concern. Nor must someone feel fulfilled at every moment to participate in a meaningful relationship. Meaning, the value of a life, and comprehensive self-understanding should remain connected without being reduced to a single capacity.

Meaning need not await a unique final destination for a life. Caring for a particular person, completing an honest inquiry, preserving an important memory, or repairing a damaged relationship can matter within a finite history. Their meaning does not require the universe to have assigned them an overarching purpose in advance. It depends on real beings for whom things matter, actual relationships, and value for which reasons can be given.

This does not make meaning entirely a matter of arbitrary personal creation. A person can declare anything meaningful to them, but public discussion can still ask what the activity actually connects them with. Does it rest on deception? Whose lives does it help, and on whom does it impose costs it refuses to acknowledge? Subjective engagement is important to meaning, but intensity of engagement cannot establish by itself that its object deserves such devotion.

The original paper places dynamic coordination between the individual and the world at the center of meaning. That insight can be developed further: personal plans depend on bodies, other people, and a public environment, and meaning often arises through participation in activities extending beyond an isolated self. Coordination, however, must not mean unconditional accommodation to an existing whole. Someone who refuses an unjust order may come into conflict with their organization, yet find profound meaning in taking responsibility for a reasoned objection. The direction of social stability is not always a direction worth endorsing.\footnote{\hypertarget{ph-note-57}{}Original paper, pp.~24--28 and 40--41; Core, Section 27.6. This text distinguishes the experience of meaning, actual participation, and grounds for finding something worth pursuing. Justified dissent tests the language of ``alignment with the whole.''}

Commitments that require continuing performance need corresponding conditions if they are to be realized. Care requires someone's presence; research requires memory and communication that can continue; shared plans require responses that can be trusted. Preserving these conditions allows wishes to become practices for which people can take responsibility. Repair can be valuable here because it sometimes reopens opportunities to fulfill a commitment. These are conditions for carrying out particular commitments, however, rather than necessary conditions of all meaning. A justified refusal or a unique act of kindness can acquire its full significance in the moment it occurs.

Meaning must also make room for commitments that cannot be completed. A person's work may be interrupted, a relationship may end with death, and capacities may change. If meaning is equated with steady progress toward a measurable endpoint of success, such phases of life are too readily judged failures. Faithful concern, appropriate mourning, and entrusting unfinished work to those who follow can also have irreplaceable meaning.

Finitude makes this meaning more specific. Because time is limited, commitment involves relinquishing some possibilities and entering others seriously. Because people can be hurt, care is not an indifferent exchange. Because history cannot be fully repeated, honesty and repair carry weight. The meaning of life need not be preserved by erasing its finitude.

\begingroup\interlinepenalty=10000

Meaning is thus a connection that needs to be sustained and can also be revised. People can form continuity through commitments and change those commitments when they discover better reasons. A community has a responsibility to make room for this formation and revision: enabling people to contribute, to be heard, to obtain support, and, when necessary, to begin again. A unified worldview that helps us understand these connections more accurately has already made a substantive contribution to the question of meaning.

\par\endgroup

\chapter{Chapter 25: Morality Requires Reasons We Can Offer One Another}\label{ph-ch-25}

An organization can maintain order through fear or through fair cooperation. If we ask only which arrangement enables the organization to persist, we have not yet answered which arrangement people have reason to accept. Natural processes establish stakes and dependencies. Ethics must explain why something that matters to one being should count as a reason for another's action.

This book adopts three connected normative positions. Real suffering and the loss of capacities deserve serious consideration. People in comparable circumstances are entitled to comparable reasons and should not be arbitrarily excluded merely because of who they are. Shared rules should, as far as possible, enable those affected to participate in their explanation, criticism, and revision. Bodily vulnerability, interdependence, and practices of public reasoning support these positions. They still require normative defense and cannot be presented as theorems derived directly from biological survival.

Why accept these positions? First, they require an explainable consistency between first-person and third-person perspectives. If I hold that my suffering cannot be ignored merely because others benefit, I must explain why comparable suffering loses all standing as a reason when someone else undergoes it. Second, these positions can support a more reliable common life, in which no one must depend entirely on another's occasional goodwill for their basic circumstances to enter a decision. Finally, they leave room to correct errors, because those who are harmed can bring forward facts and reasons we have not yet recognized.

This argument does not assume that all interests can always be fully reconciled. Limited resources, conflicting commitments, and emergencies may force people to bear losses. Precisely for that reason, we must distinguish necessary burdens from arbitrary sacrifice, compare alternatives, explain decision procedures, and take compensation and repair seriously. Moral reasoning concerns both outcomes and the standing with which people participate in bringing them about.

An egoist who consistently treats only themselves as an ultimate object of concern can remain formally consistent. We cannot deduce from ``this person cares about their own suffering'' that ``this person must care equally about everyone's suffering.'' This book openly adopts a further normative premise: comparable injury should not lose its standing as a reason for serious consideration simply because it happens to someone else. We support this premise with reasons concerning opposition to arbitrary domination, protection of people's basic circumstances from dependence on the goodwill of the powerful, and opportunities for those affected to correct decisions. These are substantive moral reasons, not necessities disguised as conclusions of natural law. Someone who rejects them all need not have made a logical mistake, but that rejection gives them no grounds to demand that others regard their domination as already justified to those others. Ethical practice also requires education, protection, and institutions. It cannot simply wait for a deductive starting point that nobody could refuse.

Care and justice therefore need to work together. Care brings particular people into view; justice reminds us that sensitivity cannot be reserved for those we know. A sympathetic group may still neglect people who are distant or silent, while a procedure concerned only with abstract shares may miss concrete difficulties. Institutions should connect general reasons with particular circumstances, rather than require a choice between indifferent equality and partial warmth.

The ability to give reasons is a capacity needed for participation in certain decisions, not an admission requirement for moral concern. Infants, people unable to communicate, and future generations still bear the consequences of decisions. The experiences of nonhuman animals that may be capable of suffering cannot be excluded merely because they cannot debate with us. Natural processes that lack experience may be protected for other reasons, including ecological dependence, historical relationships, and responsibilities across generations. We need not assume that all beings have exactly the same moral standing. Representation does not license representatives to decide interests arbitrarily on behalf of others: caregivers and public institutions must remain answerable to evidence, scrutiny of conflicts of interest, and correction, and must listen seriously whenever those represented can respond.

This ethical approach, centered on justification to one another, has a clear affinity with Scanlon's contractualism: the justification of action must face what others could reasonably reject. This book does not claim to reproduce or replace his full theory, nor does it treat actual unanimous agreement as sufficient for justification. Actual consent may be shaped by fear or missing information, while justified objections from a minority may remain suppressed within a stable institution. Mutual justification is a demand by which institutions can be criticized, rather than a certificate for existing consensus.\footnote{\hypertarget{ph-note-58}{}T. M. Scanlon, \emph{What We Owe to Each Other} (Belknap Press of Harvard University Press, 1998). The author's Harvard bibliography and the publisher's book description were checked. This book is situated only with reference to the publicly stated themes of mutual justification and reasonable rejection; arguments not read in full are not asked to carry its burden of proof.}

Naturalism can make a positive contribution here. It helps explain how bodies and environments affect capacities, how institutions expose some people to vulnerability, and how care and education can extend agency. Ethics then asks which extensions we have reason to support together and which costs cannot lightly be transferred to others. Empirical inquiry and normative judgment meet through explicit reasons, and each can change the questions the other asks.\footnote{\hypertarget{ph-note-59}{}Original paper, pp.~24--28, on naturalistic ethics; Core, Chapters 18, 24, and 27, which explicitly retain evaluative bridges and constraints between subjects. The principles of care, mutual recognition, and public reasons advanced in this book are defended in the main text.}

Moral universality does not require every life to be the same. We can demand mutual respect and consistency in basic reasons while allowing people to pursue different goods. A shared world is worth sustaining precisely because it contains many lives that matter. A more mature philosophy of unity should make room for this diversity in its account of organization, rather than abolish it in the name of unity.

Theories of growth are especially unable to dispense with this work of justification. That a mechanism expands faster means only that, under the given environmental, material, and transmission conditions, it increases a specified count more effectively. It does not explain why the suffering of those sacrificed may be ignored or why the claims of the weaker deserve no hearing. To turn advantage directly into justification is to move without argument from a factual relation to normative authority.

This book's positive position on normativity begins with agents who differ from one another yet are interdependent: we undergo consequences, can make claims, and seek a degree of reciprocity in the reasons offered within a shared world. A rule that can defend itself only by appealing to the identity of the powerful, while refusing to consider others in comparable circumstances, is open to criticism on grounds of equal concern and opposition to domination. That argument depends on commitments concerning experience, equal concern, and shared life. It should face objections openly, rather than use biological growth or a theorem of formal control to evade them.

\begingroup\interlinepenalty=10000

Knowledge of mechanisms remains indispensable. It can reveal whether a moral commitment requires impossible resources, how an institution transfers risks, and how apparently universal rules impose injuries from which some people cannot recover. Factual knowledge helps make reasons answerable to reality; reasons determine how we assess those consequences. Holding these responsibilities together keeps meaning from being reduced to how many beings continue to exist.

\par\endgroup

\chapter{Chapter 26: Institutions, Markets, and the Conditions of Collective Rationality}\label{ph-ch-26}

An institution connects many encounters. Commitments can extend through time, rules can operate without particular people being present, and a role can continue after its occupant changes. An institution is therefore a real organization: through information, resources, entitlements, and expectations, it changes what actions are possible. Treating institutions as objects helps us investigate how they form, persist, and can be repaired.\footnote{\hypertarget{ph-note-60}{}Original paper, pp.~8 and 18--28, on hierarchy and social objects; Core, Chapters 15, 24, and 29. An institution's organization and its legitimacy are treated as connected but distinct questions.}

The objecthood of institutions also depends on participants' practices. Rules on paper must be understood, enacted, and contested; statistical indicators must be measured and used; roles must be taken up. Institutions thus possess a continuity extending beyond any one person while remaining open to questions about their concrete realization. ``The institution requires it'' cannot indefinitely serve as the last word that stops explanation.

Markets offer an important example. Prices, commitments, property rights, information, and rules of exchange connect many dispersed actions. Under appropriate conditions, local signals can direct resources toward new needs and permit cooperation without requiring participants to know everything about the world. Those conditions are themselves part of an institutional organization. Which costs enter prices, who has information, what commitments can be trusted, and who can participate in exchange all affect the result.

The transition from individual gain-seeking to common benefit therefore requires an account of mechanisms. A person's gain can come from providing real value to others, or from shifting losses onto people without bargaining power. If consequences do not enter the assessment of transactions, local optima can accumulate into public harm. The original paper's concerns about market bubbles, concentration, and crises find a more precise philosophical place here: we should investigate how institutions connect local feedback processes, without assuming that those processes must lead to a collective optimum.\footnote{\hypertarget{ph-note-61}{}Original paper, pp.~18--24. Its discussions of bubbles, concentrated wealth, and crises supply the questions pursued here. Particular distribution laws, risk measures, and policy effects are not presented as universal facts already established by UEOT.}

Efficiency and justice must likewise be distinguished. Imagine two people dividing a fixed quantity of a resource. Both regard having more as an improvement, and one of them already holds everything. If the comparison permits only changes to their shares, improving the other's position requires reducing the current holder's share. The arrangement may satisfy the condition that no one can be made better off without making someone else worse off, while leaving its justice entirely open to question. The rules of comparison used to judge efficiency cannot, by themselves, determine the legitimacy of rights or initial distributions.

The accumulation of wealth and concentration within networks also require investigation through specific mechanisms. Returns can be reinvested, informational advantages can open further opportunities, and rules can reinforce established positions. These may produce cumulative processes. The existence of feedback, however, neither supplies a distributional law common to every society nor determines a unique reform. Comparing explanations requires attention to channels of transfer, time horizons, alternative mechanisms, and affected groups.

Systemic risk often arises from mismatches between time scales. Rapid local responses may improve an individual participant's position while collectively removing support on which the whole depends. Slow institutional adjustment may stabilize expectations yet fail to respond in time to changes already taking place. Drawing a network identifies connections; we must still establish what obligations each connection carries, how losses spread through it, and what buffers are available. An elegant aggregate indicator cannot replace those questions.

From this perspective, a basic task of public governance is to bring shared consequences into shared decisions. Disclosure of information, clear responsibility, appropriate dispersal of power, channels for objection, and capacities to recover after local failures can all be posed as questions of institutional design. Their specific realization should be informed by evidence and public discussion. Philosophy can clarify evaluative aims. A unified framework alone cannot calculate one policy for every society.

\begingroup\interlinepenalty=10000

Collective rationality does not require everyone to share a single mind. It can arise through reliable divisions of work, mutual checking, and procedures open to correction. A whole worth pursuing is an organization in which dispersed knowledge can make a difference, those who suffer harm can speak, and public errors can be acknowledged. Such a whole will contain disagreement, but disagreement itself can become part of its capacity to know.

\par\endgroup

\chapter{Chapter 27: Civilization, Technology, and a World Larger than Ourselves}\label{ph-ch-27}

Individual action depends on language, tools, institutions, and natural conditions, most of which no individual has created alone. Civilization can be understood as a process that sustains and transforms the conditions of action across generations: preserving knowledge, forming public rules, and building infrastructure, while also creating new dependencies and risks. It connects many lives within a longer history.

The original paper pictures civilization moving toward more highly organized information and more extensive coordination of the whole. Such a picture can encourage responsibility, but civilizational progress cannot be measured solely by the quantity of information, resource efficiency, or system size. Stronger organizational capacities can improve care or extend surveillance; faster communication can share knowledge or spread misinformation. To assess progress, we must ask how capacities are distributed, whose lives improve, and whether new powers can be subject to shared constraints.\footnote{\hypertarget{ph-note-62}{}Original paper, pp.~24--28 and 40--41, on the subject--world complex and the direction of civilization. Civilizational maturity is proposed as a direction requiring evaluative justification. No monotonic historical increase in an information-to-resource ratio is assumed.}

Technology makes this question especially clear. An automated system may display complex prediction, selection, and self-adjustment within a specified task, and those capacities deserve careful study. Performing a task better does not itself make the task appropriate or determine who is entitled to deploy the system. The purposes and responsibilities of technology need to be explained through actual institutional arrangements. They cannot be assigned to a supposedly neutral tendency toward efficiency.

In systems with complex learning capacities, understanding, agency, and the possibility of experience should also be considered separately. We can investigate what history a system preserves, how it selects actions, and how it responds to reasons. Whether it has subjective experience requires further argument. Careful distinctions help avoid two practical errors: ignoring real effects because experience remains uncertain, and transferring all responsibility to an unexplained ``autonomous subject'' because its abilities are impressive.

The natural environment reveals a deeper openness in civilization. A society can grow on its balance sheet while depleting the conditions of recovery on which it depends. Treating the external world solely as a supply of replaceable inputs overlooks the time scales and irreversibility of many changes. Ecological thinking asks which dependencies are not recorded directly in current transactions and which consequences extend across places and generations.

This does not require imagining the Earth as a super-person with a single will. We can acknowledge real ecological relations without assuming an overall purpose that can simply be read off from them. Human responsibilities toward nonhuman life and natural processes can arise from several kinds of reason: they support human life, they include lives and experiences worthy of concern, and they embody histories and relationships that immediate interests cannot adequately measure. Particular ethical positions remain disputable, but real dependencies mean that the discussion cannot be deferred indefinitely.

Responsibility across generations makes the open boundaries of a community especially apparent. Those who come later cannot participate in all of today's decisions, yet they must live within the conditions those decisions leave behind. We cannot determine a complete conception of a good life for them. We can try to preserve their capacity to form, revise, and carry out plans. There are therefore reasons beyond present preference to sustain possibilities of recovery, the checkability of knowledge, and the revisability of public institutions.

What deserves preservation here is a set of conditions for life that we have reason to sustain, rather than the capacity of any institution to restore itself. Oppressive arrangements may be highly effective at recovering after a crisis. Nor can future choice justify indefinitely postponing help in the present. Immediate irreversible harms and the opportunities of later generations must both enter public comparison. Responsibility across generations requires us to explain what we are preserving and who bears the costs; it does not establish continuation itself as the highest good.

The maturity of civilization can accordingly be understood as a reflexive capacity. Alongside increasing its power to change the world, civilization can come to recognize the conditions, distribution, and consequences of that power. A more unified worldview should make us more attentive to these connections, rather than give civilization a mission whose endless expansion is automatically justified.

\begingroup\interlinepenalty=10000

A world larger than ourselves does not need a cosmic command to deserve serious attention. It already enters our lives through dependencies, encounters, vulnerability, and shared history. Acknowledging that fact while offering public reasons for concern is work that philosophy of nature and ethics can pursue together.

\par\endgroup

\part*{\phantomsection\hypertarget{ph-part-6}{}\addcontentsline{toc}{part}{Part VI: The Methods of Unification and Progress toward Truth}{\Large Part VI}\\[1.2em]{\LARGE The Methods of Unification and Progress toward Truth}}

\chapter{Chapter 28: Making Explanations Answerable to Failure}\label{ph-ch-28}

A framework that can produce a retrospective account of any result has difficulty telling us what we did not already know. Explanatory strength comes from assuming a risk of failure before the results arrive: specifying which differences matter, what consequences should follow, and what observations would lead us to revise the account.

The original paper proposes an analytical route from identifying objects and profiling their structure, through the proposed pair of concerns, direction, and long-run organization, to integrated diagnosis. This route can be reconstructed as a disciplined way of asking questions. Specify what is being tracked; identify its actual interfaces and dependencies of maintenance; propose evaluations and mechanisms; then use independent results to test the links among these explanations. Unified concepts can make explanatory responsibilities more visible rather than remove the need for concrete work.\footnote{\hypertarget{ph-note-63}{}Original paper, pp.~9--11 and 37--40; Core, Chapters 16--17 and 28--31; and the scientific-interface documents distinguishing independent certification, counterexamples, and unresolved conclusions.}

An explanation must first keep its question stable. If a task has failed and the researcher then declares that this outcome never belonged to the object's own future, object boundaries become a device for avoiding refutation. Boundaries may be revised, but a revision must state which earlier claim it changes. Continuing to use the same name cannot erase a change in the object, protocol, or evaluation.

Second, conceptual relations must be distinguished from quantitative measures. Saying that a pressure diverges from an organization's direction is insufficient to define a universal score of divergence. We must examine how the direction is established, how zero pressure is handled, and whether the measure mistakenly classifies a directly opposite direction as aligned. Ratios that combine resources and risks measured in different units also require a calibration procedure. The original diagnostic indicators may be tools for further investigation, but a tool's name cannot establish its validity.\footnote{\hypertarget{ph-note-64}{}Original paper, pp.~9--11, on preference intensity and deviation measures. This chapter raises questions about definition and calibration; it supplies no new universal empirical thresholds for these measures.}

Third, success within a model must be connected with sources of evidence in the world. A proof concerning a specified process can show that a conclusion necessarily follows if the world satisfies certain conditions. A simulation can show these conditions being realized together in a particular construction. An empirical experiment must test whether an external process actually possesses the required properties. These activities can support one another while answering different questions. A world generated by researchers does not become confirmation about the external world simply because its construction is complex.

With finite data, comparisons fixed in advance, retention of all attempts, explicit accounts of sampling, and independent examination of failures can reduce selective explanation. The point is not a ritual multiplication of tables. It is to let others determine what researchers have seen, which results have not been used, and how many conclusions rely on the same evidence. Repeated use of a single calibration event must not be counted as multiple independent successes.

Failures also require accurate attribution. An instrument malfunction does not directly overturn every theory. Failure of a model's applicability conditions is not a contradiction in its deductive proof. Failure of an evaluative bridge to identify an actual purpose does not imply that the object under investigation does not exist. But these distinctions must remain open to criticism. A theory's proponents cannot unilaterally relegate every unfavorable result to peripheral assumptions. If multiple reasonable applications fail systematically, more central assumptions need to be reconsidered.

A sharper objection asks whether a unified theory can ever count as having failed if every counterexample merely leads it to narrow its scope. Open revision has epistemic value, but it does not automatically rescue the original claim. A revision must record the scope that has been lost and show whether the remaining proposition still has nonempty instances, independent content, and consequences that distinguish it from alternatives. If all that remains is ``whatever meets the conditions meets the conditions,'' its explanatory contribution has been hollowed out. The finite realizations retained in this book support the corresponding conditional inquiries; they do not retrospectively vindicate the original universal theory. If UEOT produces no further explanatory gains or empirical distinctions, it may remain primarily a way of organizing questions.

Discriminating comparisons are especially valuable in research aimed at unification. If several theories predict some form of cycle, finding a cycle cannot, by itself, support one of them over the others. Their different predictions about interventions, time scales, or new environments must be compared. An experiment is decisive because of the way it distinguishes competing explanations, not because of the grandeur of the object it observes.

Making explanations answerable to failure also means recognizing that ``uncertain'' is sometimes the best current conclusion. Certainty must be earned through evidence; it cannot be imposed on the world as an attitude. A theory that locates what is unknown accurately has already reduced the cost of finding a direction for those who follow.

One test for the definitive core developed in this edition is to ask which independent premise a conclusion has actually eliminated. Giving ``successful maintenance'' in advance and then deriving maintainability, supplying fitness and then showing that greater fitness confers an advantage, or first supplying a conditional mean and then naming it an actual source does not advance the explanation. The new common source explicitly specifies programs, damageable components, read permissions, resources, and noise, then derives recovery, ordering, and finite feedback from execution records. Existing mathematical tools establish the verified connections within that process. This is a checkable explanatory contribution, rather than merely another aggregate certificate or a larger collection of files.

Counterexamples must also retain the conditions under which they arise. Under subcritical conditions, refusal can favor temporary continued existence more than continued action; a program optimal for growth may become extinct at the outset; a representation sufficient to predict program code may omit the execution gate. Once these facts are found, the positive class and the text must be explicitly revised. One cannot silently add free resources, remove permitted actions, or substitute a more accommodating world. Proportionate verification should examine the connections most liable to distort the conclusion, rather than manufacture an appearance of certainty through repeated large-scale testing.

\begingroup\interlinepenalty=10000

Formal proof, actual program execution, independent exact checks, and empirical evidence from nature carry different kinds of warrant. The first three can jointly support the claim that a publicly specified model realizes the conditional propositions stated of it. They have not established that nature universally follows the same model. This book therefore explicitly marks empirical support for universality as ``UNVERIFIED.'' The distinction gives completion of the finite model an accurate meaning while preserving the indispensable role of external empirical tests.\textsuperscript{\hyperlink{ph-note-3}{3}}

\par\endgroup

\chapter{Chapter 29: Cosmic Unity as an Explanatory Horizon}\label{ph-ch-29}

People can feel connections among the stars, living things, and their own minds. That experience is not automatically a scientific theorem, but it raises real philosophical questions. Does the diversity of the world have deeper common conditions? Can topology, information, dynamics, and organization yield compatible explanations across scales? The original UEOT paper extends these questions to fundamental physics and the universe as a whole, revealing the widest reach of its ambition.\footnote{\hypertarget{ph-note-65}{}Original paper, pp.~29--39, on conjectures in fundamental physics and cosmology. This text discusses the explanatory obligations they place on a unified philosophy of nature. Those conjectures are not included among the core mathematical results completed in the work reported here.}

This book retains that horizon. The search for unity can lead us to recognize recurring structures, ask whether apparently independent principles arise from common mechanisms, and try to understand more phenomena with fewer assumptions. Philosophy can give conceptual direction to this work while requiring a clear account of what would constitute an actual derivation.

A common vocabulary is not yet a common theory. Calling uncertainty, biological adaptation, and social ignorance ``information'' merely invites a comparison. A unified explanation requires definitions of the respective informational objects, an account of their correspondence, and reasons why a theorem in one domain can carry over to another. Likewise, using Ω for cycles in trajectories, causal dependencies, and spatial topology may preserve a family resemblance, but it cannot replace bridges among these mathematical and physical objects.

The original paper proposes three groups of conjectures in fundamental physics: understanding quantum structure through organization and information; explaining particle properties through topological organization; and understanding spacetime and gravity through collective informational relations. They remain directions for unified inquiry in this book, rather than an already completed picture of nature. Moving from statistical distinguishability to spacetime geometry requires explaining why the metric, causal structure, and dynamics take the required forms. Moving from topological classifications to particle theory requires correspondences with observables, interactions, and numerical predictions. Connecting theories through changes of scale requires genuinely compatible processes. Each is a distinct explanatory responsibility in which progress can be made.

We should especially distinguish the ability to construct a description from the claim that reality realizes its corresponding mechanism. Mathematics generally allows many objects and many representations. Physical explanation must also say why the relevant conditions apply, how parameters are determined, and how alternatives are excluded. The appeal of unity may guide a search; derivation and observation must jointly guide the choice.

Toward its conclusion, the original paper suggests that the organization of objects and the proposed pair of concerns may be different expressions of a deeper principle of action. As a metaphysical conjecture, this expresses the possibility that order in the world has common generative conditions. As a program in mathematical physics, it requires a specified domain for the action, an admissible class of variations, dynamical equations, and an interpretation in the world. Only as these elements are developed can a unified action become a theory with consequences rather than a name.\footnote{\hypertarget{ph-note-66}{}Original paper, p.~41, proposing an ultimate unified action principle. It is retained as a research prospect, not assumed as an already established mathematical or physical premise.}

Purpose at the level of the whole universe faces an additional difficulty. We can discuss how a living being sustains exchanges with its environment because its boundaries and conditions can be distinguished. Transferring the same account to the entire universe requires explaining what environment, maintenance, and success mean there. Even if the universe has an overall structure, it does not thereby become a person with intentions. Organization of a whole, self-maintenance of a living kind, and purpose as experienced each require their own arguments.

This does not impoverish the universe of meaning. Human beings can understand parts of nature, living beings can develop concerns, and societies can jointly inquire into what is real. These activities already take place within the universe. The existence of meaning does not require every part to possess the same kind of meaning, or the universe as a whole to make ethical choices for us. A natural world containing real subjects and shared practices of value is already richer than a picture of unrelated events merely accumulating.

\begingroup\interlinepenalty=10000

This book's attitude toward ultimate unity is therefore open but directed. It is worth seeking deeper structures that make different explanations compatible, proposing bold conjectures, and relinquishing attractive bridges when necessary. If nature is unified, that unity should be able to withstand questions that become progressively more precise.

\par\endgroup

\chapter{Chapter 30: Truth and Co-evolution}\label{ph-ch-30}

A person can sincerely believe something false; a group can unanimously accept something false; a belief that favors short-term survival can also conflict with the facts. Truth therefore cannot be defined solely by conviction, consensus, or adaptive success. This book adopts an open realism: whether a claim about the world is true depends on how the world actually is, while our understanding approaches it through evidence, reasoning, and revision.

This position does not require us already to possess complete knowledge unrestricted by any perspective. Every judgment uses a language and a method, and observation is organized through instruments, questions, and background understanding. Yet mediation by methods does not mean that the results are ours alone to determine. Whether a bridge will bear a load, whether a response recurs, and whether a counterexample satisfies the premises can all force changes in our expectations. The world enters inquiry through consequences that cannot be negotiated at will.

The objectivity of truth can thus coexist with the historical character of knowledge. Within history, we form concepts, discover omissions, and alter ways of measuring and proving. Some earlier conclusions are overturned, some restricted, and others strengthened by wider testing. Historical change itself does not decide which conclusion is correct, but it can create better conditions for seeking truth.

The co-evolution of philosophy and mathematics has a clear place here. Philosophy proposes conceptual relations and universal ambitions; mathematics clarifies their premises and supplies positive results and counterexamples; models and experiments investigate how those premises apply to the world. Philosophy then reconsiders its concepts, while mathematics addresses structures raised by the new questions. The two books remain separate so that each form of argument can bear its own responsibility. They refer to one another because the same question can cross the boundaries between those responsibilities.\footnote{\hypertarget{ph-note-67}{}Core, Chapter 31 and Appendices C and E; the three scientific-interface documents; mathematical volume, Chapters 34 and 38. This chapter adopts and defends an open realism concerning objective truth and historically situated knowledge.}

Formal proof should also be understood within this process. It allows an inference to be checked step by step, making it harder for rhetoric to conceal the difference between what has been given and what has actually been obtained. It cannot choose every question worth caring about on a researcher's behalf, and successful compilation alone cannot establish that the external world fits a model. Yet when an apparently natural conclusion is shown to require additional premises, or a small counterexample reveals a gap in an ambitious account, philosophy acquires a specific reason for revision.

Philosophy likewise does more than wait for technical results and then append a commentary. It can ask whether an indicator measures what originally mattered, whether the way an object is identified conceals affected parties, whether an objective function imports an unacknowledged value judgment, and why a certain success deserves to count as progress. These questions can change the design of subsequent models. Even a perfect solution to an excessively narrow question may fail to meet the original need for understanding.

Like the other organizations examined in this book, the organization of shared inquiry has real conditions of maintenance. Materials must be open to inspection, critics must be able to speak, errors must be publicly acknowledged, and credit and responsibility must be properly recorded. Independent checking matters because different positions can reveal different omissions, not because another investigator is guaranteed to be right. Honestly preserving failures gives knowledge a memory more complete than a narrative of success alone.

Locating the objectivity of knowledge within social processes of public criticism is an idea developed by philosophers of science such as Longino. The lesson this book draws is that correcting individual limitations requires an organized diversity of perspectives, rather than simply greater personal confidence. Conditions for seeking truth must still be distinguished from truth itself. Open discussion can improve the chances of detecting error, but it does not automatically make the majority right. Nor is a claim necessarily false simply because there is currently no way to test it.\footnote{\hypertarget{ph-note-68}{}Helen E. Longino, \emph{Science as Social Knowledge: Values and Objectivity in Scientific Inquiry} (Princeton University Press, 1990). The publisher's original description was checked. Its treatment of objectivity and the social processes of knowledge locates this chapter's discussion. A table of contents or publisher's summary is not represented as a reading of the book's complete argument.}

Progress can occur before a final answer is reached. Dividing an unclear claim into discriminating questions, proving one part, finding a counterexample to another, and specifying what further evidence from the world is needed already change our epistemic situation. Knowledge grows through more than an increasing number of affirmative statements. It also grows as the distance between claims and their grounds becomes smaller.

\begingroup\interlinepenalty=10000

The worldview developed here returns to its starting point. Objects form amid change and persist through relations; some also participate in the future through memory and action. Within these natural conditions, human life develops reasons, commitments, and shared responsibilities. The desire for unity makes us seek connections; the resistance of reality requires those connections to be carefully demonstrated and explained. What we can sustain is a capacity for better reasons, more accurate evidence, and more just forms of shared life to keep emerging, rather than an intellectual endpoint that will never need revision.

\par\endgroup

\chapter{Appendix A: How the Concepts Are Used}\label{ph-app-a}

This table clarifies the book's terminology. An explanation of a concept does not prove that the relevant kind of organization exists in every real system.

\begin{longtable}[]{@{}
  >{\raggedright\arraybackslash}p{(\columnwidth - 2\tabcolsep) * \real{0.5000}}
  >{\raggedright\arraybackslash}p{(\columnwidth - 2\tabcolsep) * \real{0.5000}}@{}}
\toprule\noalign{}
\begin{minipage}[b]{\linewidth}\raggedright
Term
\end{minipage} & \begin{minipage}[b]{\linewidth}\raggedright
Meaning adopted in this book
\end{minipage} \\
\midrule\noalign{}
\endhead
\bottomrule\noalign{}
\endlastfoot
Object & The book focuses on concrete entities that can be tracked through a period of change. The relations that must be preserved depend on the question, timescale, and range of interventions. This definition does not deny the reality of brief events or processes of formation. \\
Process & The formation, continuation, and alteration of states and relations over time. Persistent objects are realized through processes. \\
Relational realism & A description is selected in light of a question, but its adequacy is constrained by real relations in the world. The involvement of a perspective does not amount to the arbitrary creation of facts. \\
Boundary & An interface distinguishing internal from external activity. Questions about prediction, intervention, resources, and identity may call for different boundaries. \\
Ω & An inquiry into the constitutive dependencies, redundancies, and modes of failure relevant to sustaining an organization. Geometric cycles are only one potentially relevant kind of structure. \\
Historical identity & Tracking an object through relations of formation and continuation. Historical continuity alone does not establish strict identity; branching requires particular care in distinguishing inheritance, functional similarity, and numerical identity. \\
Openness & An object persists through exchanges and external conditions. The role of openness must be specified in terms of interfaces, content, timing, and dependencies. \\
Repair & The recovery of a capacity while preserving specified relations of continuity. The formation of a new object, copying, and reproduction require separate accounts. \\
Memory & The preservation of historical differences as distinctions available for subsequent prediction and action. Its sufficiency is relative to a particular question. \\
Feasibility & The existence of a complete strategy that can actually be implemented and can respond to the relevant subsequent branches. The constraints and time horizon must be specified. \\
Purpose & A relation through which future outcomes, distinguished as successes, failures, or objects of preference, enter into the organization of present action. Functional, historical, intentional, and jointly committed forms require separate accounts. \\
Π/Φ & A dual concern with resources for implementation and informational adequacy. A quantitative theory requires independent definitions, units, mechanisms, and calibration. \\
Representational freedom in evaluation & Different numerical decompositions or representations may preserve the same overall evaluation or choice relation. Additional anchoring determines the extent of identification. \\
Long-term organization & Invariance, cycles, occupation proportions, conditional survival, and other specified forms of temporal behavior. Each requires appropriate conditions of existence and stability. \\
Agency & The question of how a system regulates action and its relations with the environment through its own organization. Criteria for minimal agency still require a defense. Reflective agency also involves examining reasons, programs, and goals; this further capacity is not a necessary condition of all agency. \\
Meaning & A relation in which a subject understands an activity as participation in something worth caring about, and incorporates that participation into the history of a life. Experiential engagement, real connections, and normative evaluation have distinct roles. \\
Freedom & A conditional capacity involving understanding, responsiveness to reasons, and real opportunities for choice. It must be assessed together with the social arrangements that support or undermine these capacities. \\
Unification & The establishment of warranted connections and relations of compatibility among different explanations. A single universal theory is a stronger research ambition. \\
Truth & The correspondence of a claim with how the world actually is. Knowledge approaches truth through evidence, reasoning, and shared correction over time. \\
\end{longtable}

\chapter{Appendix B: Correspondence with the Original Paper}\label{ph-app-b}

The original paper supplies the questions from which this book begins. The book reorganizes its parts according to the relationships among arguments and retains all its principal themes. Retaining a question does not mean retaining every strong claim made about it in the original. Page numbers refer to the physical page sequence of the \href{https://github.com/MurphyHoops/UEOT/blob/3d443ec29eb17af1fc00699a74e2d42d73516d81/UEOT_v3_top.pdf}{original 43-page PDF}; themes extending across pages therefore overlap.

\begin{longtable}[]{@{}
  >{\raggedright\arraybackslash}p{(\columnwidth - 4\tabcolsep) * \real{0.3333}}
  >{\raggedright\arraybackslash}p{(\columnwidth - 4\tabcolsep) * \real{0.3333}}
  >{\raggedright\arraybackslash}p{(\columnwidth - 4\tabcolsep) * \real{0.3333}}@{}}
\toprule\noalign{}
\begin{minipage}[b]{\linewidth}\raggedright
Location and theme in the original paper
\end{minipage} & \begin{minipage}[b]{\linewidth}\raggedright
Where this book takes it up
\end{minipage} & \begin{minipage}[b]{\linewidth}\raggedright
Principal revision
\end{minipage} \\
\midrule\noalign{}
\endhead
\bottomrule\noalign{}
\endlastfoot
Pages 1--3: divisions among matter, mind, and meaning; common questions of systems thinking, control, and information & Preface; Chapters 1 and 30 & Begin with the motivation for unity and argue for explanatory compatibility. Distinguish a common world from a single ultimate theory. \\
Pages 4--5: topological existence, Ω, types of objects, openness, and life cycles & Chapters 2--7 and 20 & Develop the account through persistent organization, constitutive dependence, and modes of failure. Do not require every object to exhibit the same geometric loop or life cycle. \\
Pages 6--7: Π/Φ, GOD/GOA, coupling, and fields & Chapters 3 and 11--14 & Retain the dual concern with resources and knowledge, while supplying arguments about the sources of purpose, measurement, and long-term behavior. Explain particular effects through actual couplings. \\
Page 8: cooperation and competition at one level; hierarchies of wholes and components & Chapters 16, 18, and 26 & Distinguish joint activity, the formation of a parent object, and public justification. The persistence of a whole does not by itself give its members a reason to obey. \\
Pages 9--11: structural profiles, local preference indices, and five-step diagnosis & Chapters 17 and 28 & Reformulate diagnosis in terms of questions, interfaces, calibration, and explanations that can fail. Ratios and directional indices require independent validation. \\
Pages 12--15: biological evolution, exploration and optimization, and lineage continuity & Chapters 19--20 & Separate repair, individuals, inheritance, and units of selection. Innovation may change the space of possibilities; an evolutionary direction requires a specific mechanism. \\
Pages 15--18: selfhood, consciousness, attention, will, memory, empathy, and psychological disorder & Chapters 8, 15, and 21--23 & Develop an account of subjects with bodies and histories. Connect functional organization with the question of experience, while leaving specific clinical and neuroscientific claims to independent evidence. \\
Pages 18--24: markets, wealth concentration, bubbles, network risk, and fast and slow levels & Chapters 18 and 26 & Explain local and collective consequences through institutional mechanisms. Assess efficiency, justice, stability, and predictive capacity separately. \\
Pages 24--28: subject--world complexes, meaning, value, ethics, and freedom & Chapters 23--27 & Offer independent normative arguments through commitment, care, public reasons, and real opportunities. Preserve reasoned disagreement and the irreplaceability of individuals. \\
Pages 28--29: open questions concerning core concepts, mathematics, and measurement & Chapters 10, 13, 28, and 30 & Integrate conceptual anchoring, model applicability, finite evidence, and the organization of correction into the methodology. \\
Pages 30--33: conjectures about quantum phenomena and particle topology & Chapter 29 & Retain the direction of research toward unification. State that moving from structural analogy to a physical theory still requires observables, dynamics, and numerical correspondences. \\
Pages 34--37: information geometry, gravity, singularities, the cosmological term, black-hole information, and unification across scales & Chapters 17 and 29 & Discuss these as the original paper's prospects for fundamental physics. Scale mappings, temporal processes, and compatible dynamics carry separate explanatory responsibilities. \\
Pages 37--39: proposals concerning topological data, particle resonances, and cosmological observations & Chapters 4 and 28--29 & Experiments must distinguish competing explanations and begin with specific predictions. Structural indications alone cannot confirm an entire worldview. \\
Pages 40--41: a unified worldview, purpose and meaning, and an ultimate action principle & Preface; Chapters 24 and 29--30 & Retain an open naturalistic ambition. Meaning has real bearers, while ultimate unification requires evidence for each proposed connection. \\
Pages 41--43: references, acknowledgments, and author contributions & This book's attribution statement and Appendix D & Retain Degui Qian as the author of the original paper. Its references remain fully available in the original PDF. Citations not independently checked are not repurposed as this book's empirical evidence. \\
\end{longtable}

\chapter{Appendix C: Interfaces with the Independent Mathematical Volume}\label{ph-app-c}

\textbf{The mathematical specification adopted by this edition.} The mathematical volume's \href{https://github.com/MurphyHoops/UEOT/blob/3d443ec29eb17af1fc00699a74e2d42d73516d81/mathematics/UEOT_MATHEMATICS.md\#core-theory}{Definitive Account of the Core Theory} connects objects, organization, permitted action, resources, actual inheritance, and compatibility of representations in seven parts. The \href{https://github.com/MurphyHoops/UEOT/blob/3d443ec29eb17af1fc00699a74e2d42d73516d81/formalization/ueot-core/docs/compression/theory_completion/unified_closure/core_theory_closure_v1/README.md}{Core definitive delivery package} records a positive class derived from the same process, exact counterexamples, explicit restrictions of scope, and the corresponding verification. The conditional mathematical specifications completed in Sections 39.70--39.74 and the \href{https://github.com/MurphyHoops/UEOT/blob/3d443ec29eb17af1fc00699a74e2d42d73516d81/formalization/ueot-core/docs/compression/theory_completion/unified_closure/whole_theory_formal_closure_v1/FINAL_ACCEPTANCE.md}{whole-theory acceptance record} remain valid. The joint-decision, decision-representation, and mechanism-discrimination results of Sections 39.75--39.95 are also retained. Historical entries marked OPEN below describe their respective stages and do not supersede those later specifications. Final adoption is determined by the release receipt that matches the actual generative model, the two books, and the continuation record. Mathematical completion does not establish universality in the natural world. External empirical support remains UNVERIFIED.

The companion is the \href{https://github.com/MurphyHoops/UEOT/blob/3d443ec29eb17af1fc00699a74e2d42d73516d81/mathematics/UEOT_MATHEMATICS.md}{UEOT mathematical volume}. It retains Chapters 0--31 of the original Core text and adds work on hierarchies, calibration, object--purpose interfaces, control derived from the same process, actual occupation, feasibility based on finite support, and joint risk after selection. Either book can be read independently. The following table locates technical background; technical conclusions do not substitute for philosophical reasons.

\par\ifdim\dimexpr\pagegoal-\pagetotal\relax<12\baselineskip\newpage\fi

\section{C.1 Topics and Chapter Correspondences}\label{ph-c-1}

\begin{longtable}[]{@{}
  >{\raggedright\arraybackslash}p{(\columnwidth - 4\tabcolsep) * \real{0.3333}}
  >{\raggedright\arraybackslash}p{(\columnwidth - 4\tabcolsep) * \real{0.3333}}
  >{\raggedright\arraybackslash}p{(\columnwidth - 4\tabcolsep) * \real{0.3333}}@{}}
\toprule\noalign{}
\begin{minipage}[b]{\linewidth}\raggedright
Philosophical question
\end{minipage} & \begin{minipage}[b]{\linewidth}\raggedright
Precise work in the mathematical volume
\end{minipage} & \begin{minipage}[b]{\linewidth}\raggedright
Entry points
\end{minipage} \\
\midrule\noalign{}
\endhead
\bottomrule\noalign{}
\endlastfoot
How are observation, intervention, and the individuation of objects connected? & Protocols, canonical predictive states, and minimality relative to declared responses & \href{https://github.com/MurphyHoops/UEOT/blob/3d443ec29eb17af1fc00699a74e2d42d73516d81/mathematics/UEOT_MATHEMATICS.md\#ch-2}{Chapter 2}, \href{https://github.com/MurphyHoops/UEOT/blob/3d443ec29eb17af1fc00699a74e2d42d73516d81/mathematics/UEOT_MATHEMATICS.md\#ch-4}{Chapter 4} \\
How can patterns of damage reveal organizational integrity? & Object certificates, minimal failure structures, and conditions relating carriers to physical interventions & \href{https://github.com/MurphyHoops/UEOT/blob/3d443ec29eb17af1fc00699a74e2d42d73516d81/mathematics/UEOT_MATHEMATICS.md\#ch-12}{Chapter 12}, \href{https://github.com/MurphyHoops/UEOT/blob/3d443ec29eb17af1fc00699a74e2d42d73516d81/mathematics/UEOT_MATHEMATICS.md\#ch-13}{Chapter 13} \\
How can identity continue through change? & Historical transport, formation and binding, and distinctions among object roles & \href{https://github.com/MurphyHoops/UEOT/blob/3d443ec29eb17af1fc00699a74e2d42d73516d81/mathematics/UEOT_MATHEMATICS.md\#ch-11}{Chapter 11}, \href{https://github.com/MurphyHoops/UEOT/blob/3d443ec29eb17af1fc00699a74e2d42d73516d81/mathematics/UEOT_MATHEMATICS.md\#ch-34}{Chapter 34} \\
When can different hierarchies be combined? & Eight kinds of hierarchical relation, descent to quotients, error transport, and domain adaptation & \href{https://github.com/MurphyHoops/UEOT/blob/3d443ec29eb17af1fc00699a74e2d42d73516d81/mathematics/UEOT_MATHEMATICS.md\#ch-15}{Chapter 15}, \href{https://github.com/MurphyHoops/UEOT/blob/3d443ec29eb17af1fc00699a74e2d42d73516d81/mathematics/UEOT_MATHEMATICS.md\#ch-32}{Chapter 32} \\
Why does an identity relation not determine a numerical scale? & Defect rescaling, two-sided calibration, and composed budgets & \href{https://github.com/MurphyHoops/UEOT/blob/3d443ec29eb17af1fc00699a74e2d42d73516d81/mathematics/UEOT_MATHEMATICS.md\#ch-33}{Chapter 33} \\
How do purpose and objecthood meet? & Control with a specified evaluative bridge, purpose equivalence, and different goals for the same process & \href{https://github.com/MurphyHoops/UEOT/blob/3d443ec29eb17af1fc00699a74e2d42d73516d81/mathematics/UEOT_MATHEMATICS.md\#ch-18}{Chapter 18}, \href{https://github.com/MurphyHoops/UEOT/blob/3d443ec29eb17af1fc00699a74e2d42d73516d81/mathematics/UEOT_MATHEMATICS.md\#ch-34}{Chapter 34}, \href{https://github.com/MurphyHoops/UEOT/blob/3d443ec29eb17af1fc00699a74e2d42d73516d81/mathematics/UEOT_MATHEMATICS.md\#ch-35}{Chapter 35} \\
How can Π and Φ acquire independently grounded meanings? & Gauge freedom, independent anchoring, and necessary tests of a shared two-dimensional hypothesis & \href{https://github.com/MurphyHoops/UEOT/blob/3d443ec29eb17af1fc00699a74e2d42d73516d81/mathematics/UEOT_MATHEMATICS.md\#ch-23}{Chapter 23}, \href{https://github.com/MurphyHoops/UEOT/blob/3d443ec29eb17af1fc00699a74e2d42d73516d81/mathematics/UEOT_MATHEMATICS.md\#ch-33}{Chapter 33} \\
What distinguishes optimality from forms of long-run stability? & Constrained policy classes, optimal-action correspondences, invariant laws, occupation, and conditional survival & \href{https://github.com/MurphyHoops/UEOT/blob/3d443ec29eb17af1fc00699a74e2d42d73516d81/mathematics/UEOT_MATHEMATICS.md\#ch-19}{Chapter 19}, \href{https://github.com/MurphyHoops/UEOT/blob/3d443ec29eb17af1fc00699a74e2d42d73516d81/mathematics/UEOT_MATHEMATICS.md\#ch-21}{Chapter 21} \\
How are survival and repair connected in one world model? & Actual generative kernels, infinite paths, constrained optimality, and trusted repair mechanisms & \href{https://github.com/MurphyHoops/UEOT/blob/3d443ec29eb17af1fc00699a74e2d42d73516d81/mathematics/UEOT_MATHEMATICS.md\#ch-35}{Chapter 35} \\
How does memory preserve the possibilities needed for action? & Complete observation histories, actual posterior support, and policy-independent support recursion & \href{https://github.com/MurphyHoops/UEOT/blob/3d443ec29eb17af1fc00699a74e2d42d73516d81/mathematics/UEOT_MATHEMATICS.md\#ch-37}{Chapter 37} \\
How does local acceptability differ from complete feasibility? & Necessary and sufficient conditions for finite-horizon policy existence, executable relation checks, and reachable dead ends & \href{https://github.com/MurphyHoops/UEOT/blob/3d443ec29eb17af1fc00699a74e2d42d73516d81/mathematics/UEOT_MATHEMATICS.md\#ch-37}{Chapter 37} \\
How is learning connected with actual deployment risk? & Actual occupation under the same policy, registered calibration, obligations on the Good event, and the actual joint law & \href{https://github.com/MurphyHoops/UEOT/blob/3d443ec29eb17af1fc00699a74e2d42d73516d81/mathematics/UEOT_MATHEMATICS.md\#ch-36}{Chapter 36}, \href{https://github.com/MurphyHoops/UEOT/blob/3d443ec29eb17af1fc00699a74e2d42d73516d81/mathematics/UEOT_MATHEMATICS.md\#ch-38}{Chapter 38} \\
Why are inheritance and lineage different from persistence of the same individual? & Reproductive mechanisms, a positive-operator bridge derived from the same process, and individual--descendant relations & \href{https://github.com/MurphyHoops/UEOT/blob/3d443ec29eb17af1fc00699a74e2d42d73516d81/mathematics/UEOT_MATHEMATICS.md\#ch-25}{Chapter 25}, \href{https://github.com/MurphyHoops/UEOT/blob/3d443ec29eb17af1fc00699a74e2d42d73516d81/mathematics/UEOT_MATHEMATICS.md\#ch-26}{Chapter 26}, \href{https://github.com/MurphyHoops/UEOT/blob/3d443ec29eb17af1fc00699a74e2d42d73516d81/mathematics/UEOT_MATHEMATICS.md\#ch-34}{Chapter 34} \\
Where do knowledge and value require further explanation? & Reflective learning, the value of information, and functional, intentional, and normative interfaces of meaning & \href{https://github.com/MurphyHoops/UEOT/blob/3d443ec29eb17af1fc00699a74e2d42d73516d81/mathematics/UEOT_MATHEMATICS.md\#ch-27}{Chapter 27} \\
\end{longtable}

Each mathematical result has its own conditions concerning spaces, probability, horizons, interfaces, and measurability. Policy existence in a finite model does not establish a feasible policy for every actual system. A permitted recovery does not establish that every program can be generated from local physical information. Differences of purpose within one process, subjective experience, and public legitimacy also require the conceptual and normative arguments developed in this book.

The learning example cited in Chapter 10 is given in full in the mathematical volume. It has two registered source rows, each sampled independently 256 times under the specified experiment. Data select the two actions with probabilities \(7/8\) and \(1/8\); their actual failure probabilities are \(1/8\) and \(1/4\). The actual joint failure probability of training, selection, and deployment is \(9/64\), while the general calibration upper bound is \(17/32<1\). The former is the model's exact value; the latter is the general guarantee used. They are not presented as one and the same sharp bound. An all-true training record has positive probability and makes the empirical budget equal to 1. The claim that every training outcome satisfies the same smaller budget therefore actually fails in this example. This is a mathematical example within an explicitly generated model, not a test using data from an external system.

The general bridge from support feasibility to joint risk connects the very policy constructed from initial support feasibility with actual occupation and the joint law. It thereby removes the need to supply a separate upper bound on actual uncertified occupation. The numerical example directly proves that its selected policy has zero uncertified occupation. It shares the general bridge's interface but is not said to implement the same optimal learning algorithm. Zero uncertified occupation removes only the coverage penalty; positive reference failure and model error can remain.

\section{C.2 Historical Source Records: Read at Their Respective Stages}\label{ph-c-2}

This section retains the previous edition's complete chain of technical provenance for consultation. It is not a compulsory sequence of steps in the book's philosophical argument. Statements such as ``not yet completed'' and ``work continues'' describe the stage concerned. Later results discharge only the obligations they actually address. The current overall status is stated at the beginning of Appendix C and in C.3. Readers may start with C.3 and consult this record as needed.

The addition of Chapter 39 to the mathematical volume connected program interpretation, actual records, paid resources, offspring repair, and rankings of purpose within one fault process to the existing theorem applications. Supply remained a design condition, and the cost account covered controlled writes only. The law of the actual infinite interpreter's records and integrable returns was connected with the same policy; that interpreter still did not generate its own resource supply. A new predictive construction derived the joint law of finite future experiments, minimal sufficiency modulo null sets, and response updates from the actual past and retained latent states. It preserved the countable observation record of the same controller using complete histories; different controller memories could not thereby be merged. Timestamped birth records gave distinct birth events of one retained parent derivable identities. The continuous construction built a law over all times and an actual continuous modification from finite Gaussian laws, connecting them to the earlier quadratic Ornstein--Uhlenbeck (OU) recovery result. It did not take the required Brownian certificate as an input or thereby complete general stochastic integration, Itô calculus, or Girsanov theory. Within the declared protocol class, minimality for arbitrary actual historical pasts and fresh adaptive futures was derived from the same process. Protocol and source classes not covered, together with other global obligations, remained open at that stage. These local provenance results did not establish formalization of the whole Core theory, still less empirical confirmation.

The same chain then began to construct the stochastic action itself. Conditional moments were derived from the entire natural past of one noise process. Different partitions of the same time input were shown to yield the same actual output. Density on the genuine predictable σ-algebra then supplied the terminal integral on the full predictable \(L^2\) space. This did not complete general Itô or Girsanov theory, or every process-level result. A positive application of the finite purpose data retained the actual parent after probing and used the same returned program to generate the complete infinite record. Certificate errors, execution faults, training costs, and future evaluation retained their distinct meanings.

Further process constructions followed. Actual truncation and conditional projection produced a continuous integral on the original natural filtration. The difference between the integral's square and the actual time-integrated energy of the same input was derived to be a martingale using the actual second-order energy on past events and Fubini's theorem. The limit of the original Brownian sums of squared partition increments also received an explicit mesh-error bound. None of this completed partition quadratic variation for every full integral, general Itô or Girsanov theory, controlled stochastic differential equations (SDEs), or provenance for every dimension and filtration. In the serial-population construction, actual parent identities and empirical offspring proportions were connected to the existing lineage and Price results. Mixed-program populations did not automatically share an unchecked maintenance certificate. Provenance obligations for the entire theory were still being pursued at that stage.

Sections 39.22--39.25 gave actual partitions and weighted tables for the same full input, a complete training ledger, and joint risk for mixed-program offspring together with all execution and readout outcomes. Sections 39.26--39.29 then derived scalar Itô and Dynkin results for the same input from actual stochastic sums, under explicit derivative bounds. They derived the full vector predictable domain and a mean-square continuous process from the actual joint past, and constructed the original OU conditional kernel together with actual change-of-measure and Kullback--Leibler (KL) applications for constant controls. Mean-square continuity does not itself imply sample-path continuity; Section 39.31 separately constructed a continuous version from the same joint vector input and actual finite paths. Section 39.30 derived a complete path under change of measure from actual past-dependent feedback tables. Information cost used energy under the same actual final measure. The construction did not replace the existence of one policy with another policy or change the original past read by the feedback. General controlled SDEs and the entire predictable-control class remained incomplete at that stage.

Section 39.32 proved that scale projections retaining active registers, programs, and actual histories transport execution, birth identities, and the whole log. Funding for a continuing population explicitly used a shared pool and an external supply port associated with \texttt{idle}. When the actual initial slot was occupied and initial funds were nonnegative, supply covering the highest work cost was equivalent to an almost-sure funding guarantee for every registered plan under its actual generative law. With lower supply, a plan involving an actual parent could overspend with positive probability, while an empty-slot plan remained feasible. The limit of the uniform guarantee therefore did not make all behavior impossible, nor did a ledger with an open supply establish autonomy. Section 39.33 added useful acceptance under a separation condition for the same measurement and returned program. The original erroneous word, data, and threshold were retained. Mathematical completion and empirical support remained separate assessments.

The same numerical calibration was also connected with actual selection. The setting had an unknown two-dimensional positive rational diagonal gain, two distinct rational latent responses in \([-1,1]\), and explicit permissions for independent readout. The original selector read only the actual rational observations and used the same computed inverse and action scores. The derived guarantee bounded the joint probability of acceptance and error over the whole experiment; it did not mean that every accepted decision was correct. One admissible complete word produced an accepted error with positive probability and was retained rather than discarded as an anomaly. The registered 128 symbolic and synthetic observations reconstructed formal source words item by item. They were not described as new empirical samples and did not alter the original actual batch. By keeping the unknown true ordering in the proof conditions without supplying it to the selector, and retaining errors and failures, the certificate stated the scope of the same actual decision.

Sections 39.34--39.35 further connected the actual change of measure for the same bounded control, the constructed integral under the new measure, the original terminal information cost, and actual partition limits for a fixed joint input. Cost was actual energy under the final measure. The integral was constructed before its substitution identity was proved; the desired equality was not used as a definition. A scalar output driven by joint multidimensional noise did not automatically provide general state calculus or controlled SDEs. Section 39.36 directly checked a submitted program and the coarse program returned by computation. An exact truth table could accept a harmless hidden branch rejected by a static check, while a genuine hidden difference in action remained evidence of nonimplementability under the stated permissions. The word-by-word conclusions of 39.36 concerned finite fault words. Section 39.41 separately constructed an actual countable continuing record from the same retained parent. Costs of internal reading and branching were not silently added.

Section 39.37 proved the joint risk of an actual adaptive cache on the same data. It separately specified a new policy that reversed accepted actions solely according to the direction of the observed anchor. The new policy's word-by-word correctness under explicit noise bounds did not cancel the original policy's failure under a correlated source, and refusal did not count as a completed task. Section 39.38 computed carriers, decoders, and physical cut sets from an actual response table and proved an exponential-size obstruction for complete carrier output. This was not a general hardness result for obtaining a single answer, and it did not identify every physical deletion with death. At that stage, the whole theory and its empirical support still had separate unfinished obligations.

Sections 39.39--39.44 continued to supply actual provenance for earlier applications. States, actions, discounted occupation, a terminal law, and infinite rewards under one fixed measure generated the original control run. An explicit nonconstant evaluative specification was connected with the original optimal-control result. The occupation measure was not normalized, and noise after the end of deployment did not disappear. An Itô result for finite-dimensional states was derived from the same actual integral state, and an actual conditional martingale was proved on a specified closure of integrals under the new measure. This did not complete arbitrary controlled SDEs or all regularity classes.

Program preservation, recovery, and continuation were further distinguished. Finite truth tables checked the submitted program's actual updates and connected that program, together with the same computed return value, to the original recovery theorem and P9. The fact that \texttt{idle} preserves an intact target but cannot repair a damaged one gave a strict separation between two properties. Eventual recovery without faults and complete actual records containing faults were both retained. Programs equivalent for recovery but different in cost could not thereby be called economically or purposively equivalent. Countable continuing execution was no substitute for an infinite population of descendants, endogenous resources, or autonomy of the whole organization. At that stage, the common provenance of all these semantics and completion of the entire theory were still being pursued.

Sections 39.45--39.48 connected actual drift, OU dynamics on the physical clock, the full original predictable domain under \(Q\), and the long-run behavior and purpose of the same actual program. Left-endpoint sampling can fail for measurable drift, so the mathematics had to establish the actual integral rather than replace the object for formal convenience. The physical recovery conclusion entered the earlier application through the physical process's own law. Equality in law still did not mean equality of two paths.

The same submitted program and computed realization could then jointly support admissible frequencies, a discounted objective, GOD, and GOA under an explicit specification. Infinitely many returns under nonzero fault noise did not imply eventual permanent health. The same physical stationary state could have different costs and evaluations. Set-valued optimality allowed ties and several invariant laws at the fine level; it could not be described as a unique purpose or guaranteed survival. The evaluative specification still required justification, while its mathematical relation from mechanism to ranking now had actual execution and records as its source.

These developments did not complete every obligation concerning organization, resources, formation, lineage, or scales of splitting and merging. Nor did they establish a unique purpose for an arbitrary world. Later rotating-body mechanisms and resource and data constructions appear below. At that time, overall mathematical completion remained OPEN. Whether experience supports the model's conditions still requires investigation in the world.

Sections 39.49--39.54 realized time-varying service in the actual whole body, and copying, program inheritance, and parent identity in the same offspring that was actually born and continued to repair itself. A local birth did not guarantee a healthy whole body. Success within a block did not guarantee that an unvisited part was intact. Identity and scale relations had to follow from these actual transitions and histories rather than from a shared name or matching observations.

Under the same policy law, finite reserves, positive fault probability, and positive costs yielded a strict obstruction for healthy occupation. This was not a prediction of extinction by an arbitrary specified calendar deadline. A dynamic evaluation admitted repair, \texttt{idle}, and ties; weak optimality did not guarantee survival. Observation-driven selection was connected with the actual post-training parent, the actual returned program's long-run return, and joint risk. UNKNOWN remained when the original data failed to separate the relevant cases. Unknown truth was not given to the selector; training costs and failed terminal states were not erased.

These proofs reduced the scope for supplying funds, scale compatibility, or a correct program as independent certificates. They did not establish indefinite maintenance without external supply, a unique purpose for arbitrary organizations, or completion of the entire theory. Continued funding, countable blocks, actual hitting times, compatible connections to P12, and the remaining original obligations were pursued under their precise conditions. Mathematical completion and empirical support remained distinct.

Sections 39.55--39.57 reconnected actual payments under external supply, countable block records, and actual hitting times to their respective generative processes. Training did not become free, and neither the returned program nor the post-training parent was reset. The resource result depended on explicitly permitted supply arriving before expenditure; it did not mean that an organization created its required resources from nothing. Infinitely many returns within a block did not imply a healthy whole body. Even the complete block history could omit a passive initial state, and both the strict nonidentifiability result and the positive transport result were retained.

The actual first hit observed at a round's end entered the original recovery application through its own complete path law. Initial-state, all-fault, and empty-body cases had different outcomes. A body intact after one microstep could be damaged again by the end of a round, so an expected recovery time had to specify its observation clock. It was neither a resource lifetime nor a proof of perpetual existence for arbitrary life or all organizations. Subsequent lineage and program constructions and the other required common connections remained in progress; overall completion was still OPEN at that time. Mathematical results continued to be distinguished from empirical evidence.

Sections 39.58--39.60 made ``a repair program was returned'' refer to this actual program, and placed parent history and resources within the same complete training and deployment record. Actual birth readouts sufficed to identify the repair-target bit but not to identify the histories of two observationally equivalent parents. Additional observation was not introduced to remove this contrast. The complete record laws of native deployment and the interpreter of the same program were proved to agree, so actual long-run frequencies no longer depended on a separately supplied source-matching condition. Calibration assumptions for purpose optimality remained distinct from these facts about execution.

Finite linear low rank then supplied both the existence of a two-dimensional factorization and a classification of coordinates in the nondegenerate case. It still could not create the physical meanings of Π and Φ or supply independent interventions. Program provenance, causal identity, lineage, optimality, and life still had to be connected through actual common applications. Their separate local provability did not establish completion of the entire theory. The new six-property and actual-population constructions continued under the existing task, while empirical support continued to require evidence from the world.

Sections 39.61--39.64 connected the same data-returned program with the existing six-property endpoint, actually compared repair with \texttt{idle}, and generated two actual generations and historical Price terms by dispatching one actual founder. Failed transactions, empty requests, zero births, and negative selection terms all remained in the original law. One positive realization did not guarantee a universal adaptive advantage.

A native transaction then had an organizational path law on the same tick, but the organizational representation still forgot information in the complete record, including actual abstract syntax tree (AST) syntax and costs. Source ownership and clock, serial population registration, complete actual costs, and parent-identity discovery still needed connection within the same process. Newly obtained marker information also had to pay its reading cost. Completion of a local endpoint did not establish completion of the entire theory, and empirical support remained separate from formal proof.

Sections 39.65--39.67 proved actual transport of source identity and population registration, explicitly showing that the zero clock of a newborn could not be shifted uniformly with the other clocks. Actual births also supplied readable parent markers and an actual discovery algorithm with a probability of error. Additional reading rights and costs were listed explicitly. The original visible information did not acquire ancestral identity for free.

The original funding ledger then covered every physical writing cost in the same complete training and finite dispatch. Affording training did not mean affording the whole experiment; observation costs also had to enter the same wallet. An infinite-maintenance objective did not automatically become a finite-endpoint optimum, and nonidentifiability through a readout channel did not make the entire world indistinguishable. All the roles in finite generations still needed to be connected through actual common applications. Overall mathematical completion remained OPEN at that time.

Sections 39.68--39.69 connected all training, actual finite dispatch, and every read to one wallet. They linked the actual parent of the first birth with discovery, maintenance, and the specified evaluation. Erroneous markers and the absence of a birth retained their actual probabilities. An upper bound on an error event was not a success guarantee for the entire procedure.

The same Register process then also had canonical predictive and control quotients derived from actual transitions. Lifting the actions of the actual returned program from the quotient recovered the original program. Predictive equivalence did not become identity, and identical physical predictions could not determine purposes across different mechanisms. Section 39.70 constructed funded continuation from the actual terminal state and the same remaining balance. Actual births and Price terms, mean maintenance, and expected accounting followed from the same process. External supply for continuation was explicitly registered; a finite wallet did not become indefinite maintenance without outside supply.

The actual body, program, owner, and clock of the terminal parent came from the run and were not reset after birth. Bounds on organizational means and expected costs came from the actual marginals and accounts of the same process, not directly from stability, frequency, or aspiration. Correct marker discovery still did not guarantee a birth; the Price identity still did not guarantee convergence; optimality belonged to the registered evaluative specification. These conditional results could form a complete mathematical network without replacing normative reasons or evidence from the world.

For the original finite primitive mechanism, a general source construction then yielded positive eigendata and the earlier long-run growth result. The spectral remainder for an actual reversible killed process was also connected with the same conditional law. Each result had its own mechanism, initial state, generative relation, and spectral assumptions. Long-run composition of types did not automatically identify an individual's parent, and convergence of a conditional distribution did not imply infinite survival. These results supplied conditional mathematical content for discussing persistence and purpose; the justification of value remained a philosophical task. Their sources are in Sections 39.72--39.74 of the mathematical volume.

Sections 39.75--39.81 added a joint-decision specification with one persistent unknown source, a single wallet, and a finite endpoint. The completed status of the original set of conditional mathematical specifications was unchanged. New conclusions connected the actual output of an outer finite table to the original Register transactions, permitted receipts, native population, and prepayment at every prefix. Optimality included complete-history randomized policies required to be admissible only on actual positive-probability histories. The target was public, writing fees were positive and fixed, and the unknown parameter was initialized only once. These conditions did not authorize reading a protected program or the actual parent identity.

A complete parameter subfamily for general finite horizons used a wallet containing one writing fee. When the price ruled out paid recovery after failure, the optimal value was exactly the greater of zero and the expected value of first repair followed by receipt-dependent continuation with two objects. The theorem supplied both directions of the price boundary for positive returns, together with a zero-value result when the balance could not cover prepayment. The two-step value \(1/160\) retained the correlation induced by the persistent source. The three-step value \(9/256\) could be attained by a fixed procedure using the same menu. The four-step value \(39525/524288\) exceeded the best fixed procedure under the same conditions by \(20925/524288\). The four-step improvement came from feedback about bodily state. After success, continuation without writing was already simultaneously optimal for every possible value of the unknown source parameter, so further disclosure of that parameter had zero value. Another three-step diagnostic showed how an actual fault for which no fee was charged lowered the posterior repair threshold, making it better to stop retrying than to preserve the old prior unchanged.

These exact rational values were formal consequences of the registered model, not observations from the world or estimates of effects across domains. Birth success, optimal net return, and affordability at every prefix were verified separately. The original environment continued when intervention stopped. The earlier long-run retained-parent result, external-supply assumptions, and the finite counterexample \(-3/40\) were all preserved. Detailed proofs, original tables, and source identities are in the \href{https://github.com/MurphyHoops/UEOT/blob/3d443ec29eb17af1fc00699a74e2d42d73516d81/formalization/ueot-core/docs/compression/theory_completion/unified_closure/joint_decision_closure_v1/README.md}{joint-decision delivery package}, whose scope is stated precisely in the mathematical volume.

Sections 39.82--39.88 derived a receipt observer and count-based posterior from the same process, covering zero priors, endpoint rates, and complete-history policies randomized with real-valued probabilities. A shared graph preserved the joint law of every permitted action, native execution, and actual payments. Its actual node count had the source-specific bound \(24(H+1)^4\); bit length, computation, and complete population records were accounted for separately. Reachable counterexamples showed that bodily state alone was insufficient, some histories at the tail could be merged, and new address permissions could invalidate that merging. Proofs, THIS artifacts, costs, and boundaries for the four earlier examples and the preregistered eight-step application are in the \href{https://github.com/MurphyHoops/UEOT/blob/3d443ec29eb17af1fc00699a74e2d42d73516d81/formalization/ueot-core/docs/compression/theory_completion/unified_closure/decision_representation_closure_v1/README.md}{decision-representation delivery package}. Empirical support remains UNVERIFIED. The specification does not automatically extend to unknown mechanisms, arbitrary addresses, or internal hardware.

Sections 39.89--39.95 added a comparison of local faults S and common faults G under the same permissions. Root-only access, short horizons, and insufficient funds for prepayment gave boundaries of indistinguishability. When \(B\ge c\), the maximal total variation (TV) achievable by permitted three-step experiments was \(r^2(1-r)\); it was zero when prepayment was unaffordable. For the fixed three-step application, minimum worst-case regret over all permitted policies randomized with real-valued probabilities was \(3/128\). An independent private seed uniform on four points selected the same actually returned table and attained this bound. The best deterministic value was \(1/32\). Complete native execution, all costs, certificates covering the full policy class, and exact boundaries are in the \href{https://github.com/MurphyHoops/UEOT/blob/3d443ec29eb17af1fc00699a74e2d42d73516d81/formalization/ueot-core/docs/compression/theory_completion/unified_closure/mechanism_decision_closure_v1/README.md}{mechanism-decision delivery package}. These are conditional mathematical results within synthetic mechanisms, not empirical data. Empirical support remains UNVERIFIED. Discrimination at the terminal time is not described as a gain from subsequent learning already realized within that same three-step horizon.

\section{C.3 Current Mathematical Interfaces for Organization and Feedback within One Process}\label{ph-c-3}

The new connection in this edition concerns how the same actual program restores some of its own execution conditions and has the declared consequences for growth. The actual code bits and main execution gates of both programs can be damaged. Majority backups, a finite trusted boot mechanism, one unit of capital per old member, and five units of explicitly specified environmental supply jointly define the positive class. Chapter 7 describes the two programs and the complete comparison class. The first always performs \texttt{care}. The second performs the birth-favoring action \texttt{bloom} when healthy, and repairs itself before performing \texttt{care} when damaged. Within this class, the same damaged second program pays three units for recovery, reads the restored code to perform \texttt{care}, subsequently performs \texttt{bloom} when healthy, and attains the optimal exponential growth factor of expected population over all permitted history-dependent request policies randomized with real-valued probabilities under the same permissions and resources. A separate repair program has not been joined to the optimal program, nor have the trusted basis, information copies, or environmental materials been presented as self-created from nothing. Complete information erasure, a destroyed boot channel, insufficient funds during the current recovery stage, and the advantage of refusal in the subcritical case each have explicit boundaries.

Actual persistence and birth records retain identities, parent pointers, birth times, and payment accounts together. Population means and marked finite genealogies are derived from them. For the noisy positive class with two faithfully inherited programs, the probability guarantee concerning actual shares uses independent founder cohorts. A finite observation time is chosen first; the required initial size is then selected according to the desired precision and failure probability, including the corresponding materials and supply. This does not guarantee that a fixed finite population will never become extinct or will necessarily select the optimum over an unlimited future. An admissible representation removes only silent labels that the controller cannot observe and that do not affect these records. Finite control histories, actual branch weights, costs, and descendant records are pushed forward by the same map; actual individual identities are not merged. Agreement in ordinary predictions remains insufficient to preserve arbitrary costs or consequences for repair and reproduction.

The shared significance of these results is to reduce independently supplied premises of success, fitness, conditional means, and optimality, yielding a conditional explanation that can be checked. Principal claims, actual applications, and counterexamples are listed in the same package's theorem index and API boundaries. Public entry points, axiom scope, program execution, independent exact checks, and builds of the two books are supported by their respective evidence. Formal results for finite models and agreement with their implementations do not replace empirical evidence about nature. They do not establish agents' intentions, subjective experience, or the legitimacy of values. These questions remain the responsibility of the main text's independent arguments and the relevant evidence from the world.

\chapter{Appendix D: Sources, Levels of Argument, and References}\label{ph-app-d}

\section{D.1 Core Sources}\label{ph-d-1}

\textbf{Original work.} Degui Qian, \emph{A Universal Evolutionary Object Theory (UEOT): A Framework and Research Program for a Unified Science}, 43 pages. The public historical text is \href{https://github.com/MurphyHoops/UEOT/blob/3d443ec29eb17af1fc00699a74e2d42d73516d81/UEOT_v3_top.pdf}{UEOT\_v3\_top.pdf} at the repository root. Its original LaTeX source is preserved in \href{https://github.com/MurphyHoops/UEOT/blob/3d443ec29eb17af1fc00699a74e2d42d73516d81/philosophy/source/original/main.tex}{main.tex}. This edition checked both forms. The original work's intellectual sources, references, acknowledgments, and statement of author contributions remain as recorded in that work.

\textbf{Core account.} \emph{\nolinkurl{UEOT_Core_Mathematics_v3.0_Complete.md}}, original Chapters 0--31 and appendices. It reconstructs the original questions as a protocol-relative theory of prediction, realization, dynamics, identity, evaluation, and composition. This edition draws conceptual constraints from it; the mathematical content is preserved and developed in the companion work.

\textbf{Materials on the interfaces with science.} \nolinkurl{CORE_V3_OPEN_PORT_MATRIX.md}, \nolinkurl{CORE_V3_POST_CORE_SCIENTIFIC_CLOSURE.md}, and \nolinkurl{THEORY_COMPLETION_ROADMAP_V2.md}, all dated 7 October 2026. These materials help distinguish conditional theorems, model realizations, finite certification, and empirical support. Their historical task statuses do not replace the results actually obtained subsequently.

\textbf{Subsequent formalization and mathematical work.} The results on hierarchies, defect calibration, objecthood, representations of purpose, control within a common generative model, repair, statistical deployment, and finite-support feasibility in \href{https://github.com/MurphyHoops/UEOT/blob/3d443ec29eb17af1fc00699a74e2d42d73516d81/formalization/ueot-core/}{UEOT Core} provide the technical sources for the correspondences in Appendix C. The latest statements, proofs, and conditions are developed in the independent \href{https://github.com/MurphyHoops/UEOT/blob/3d443ec29eb17af1fc00699a74e2d42d73516d81/mathematics/UEOT_MATHEMATICS.md}{mathematical work}. The philosophical text discusses their significance in ordinary language.

\section{D.2 How the Book Uses These Sources}\label{ph-d-2}

The book's realism about objects, naturalism, account of meaning, and normative position are claims that require philosophical defense. Examples involving bridges, corridors, and institutional distributions are constructed thought experiments intended to develop concepts; they are not presented as empirical studies. The finite learning examples involving exact probabilities are supported by complete mathematical models in the companion volume.

The neuroscientific, financial, particle-physical, and cosmological studies listed in the original remain part of its historical bibliography. This book does not treat specific empirical claims that have not been independently checked as established premises for new arguments. Universal dual drives in the real world, mechanisms of consciousness, social prediction, and unification in fundamental physics still require independent investigation in the relevant fields. These distinctions allow philosophy to pursue broad questions while making clear what supports each kind of claim.

\section{D.3 Philosophical Works Substantively Discussed in This Edition}\label{ph-d-3}

\begingroup\small\raggedright\setlength{\parskip}{0.25em}

The works below provide identifiable sources for this edition's comparisons and objections. They do not imply their authors' endorsement of UEOT or constitute external peer review. The text and notes specify what is drawn from each work, where the accounts differ, and what was actually read. Verifying publication metadata is distinct from reading an entire work. All historical references in the original paper remain in its PDF. Entries are arranged alphabetically by surname and retain their original titles for retrieval.

Barandiaran, Xabier E., Ezequiel Di Paolo, and Marieke Rohde. 2009. ``Defining Agency: Individuality, Normativity, Asymmetry, and Spatio-temporality in Action.'' \emph{Adaptive Behavior} 17(5): 367--386. \href{https://doi.org/10.1177/1059712309343819}{DOI: 10.1177/1059712309343819}.

Chalmers, David J. 1995. ``Facing Up to the Problem of Consciousness.'' \emph{Journal of Consciousness Studies} 2(3): 200--219. \href{https://consc.net/papers/facing.html}{Author's public text}.

Cummins, Robert. 1975. ``Functional Analysis.'' \emph{The Journal of Philosophy} 72(20): 741--765. \href{https://doi.org/10.2307/2024640}{DOI: 10.2307/2024640}.

Fischer, John Martin, and Mark Ravizza. 1998. \emph{Responsibility and Control: A Theory of Moral Responsibility}. Cambridge: Cambridge University Press. \href{https://doi.org/10.1017/CBO9780511814594}{DOI: 10.1017/CBO9780511814594}.

Longino, Helen E. 1990. \emph{Science as Social Knowledge: Values and Objectivity in Scientific Inquiry}. Princeton, NJ: Princeton University Press. \href{https://press.princeton.edu/books/paperback/9780691020518/science-as-social-knowledge}{Publisher's record}.

Maturana, Humberto R., and Francisco J. Varela. 1980. \emph{Autopoiesis and Cognition: The Realization of the Living}. Dordrecht: D. Reidel. \href{https://doi.org/10.1007/978-94-009-8947-4}{DOI: \nolinkurl{10.1007/978-94-009-8947-4}}.

Mossio, Matteo, Cristian Saborido, and Alvaro Moreno. 2009. ``An Organizational Account of Biological Functions.'' \emph{The British Journal for the Philosophy of Science} 60(4): 813--841. \href{https://doi.org/10.1093/bjps/axp036}{DOI: 10.1093/bjps/axp036}.

Scanlon, T. M. 1998. \emph{What We Owe to Each Other}. Cambridge, MA: Belknap Press of Harvard University Press. \href{https://philosophy.fas.harvard.edu/publications/what-we-owe-each-other}{Author's university record}.

Whitehead, Alfred North. 1929/1978. \emph{Process and Reality: An Essay in Cosmology}. Corrected edition, edited by David Ray Griffin and Donald W. Sherburne. New York: Free Press. \href{https://whiteheadresearch.org/research/contemporary-whitehead-studies/proposals/suggested-editions/}{Corrected-edition record}.

Wolf, Susan. 2010. \emph{Meaning in Life and Why It Matters}. Princeton, NJ: Princeton University Press. \href{https://press.princeton.edu/books/paperback/9780691154503/meaning-in-life-and-why-it-matters}{Publisher's record}. This edition also consulted relevant arguments in the author's \href{https://tannerlectures.org/documents/meaning-in-life-and-why-it-matters/}{2007 Tanner Lectures under the same title}; the 2010 book and the earlier lectures are distinct versions.

Worrall, John. 1989. ``Structural Realism: The Best of Both Worlds?'' \emph{Dialectica} 43(1--2): 99--124. \href{https://doi.org/10.1111/j.1746-8361.1989.tb00933.x}{DOI: \nolinkurl{10.1111/j.1746-8361.1989.tb00933.x}}.

Wright, Larry. 1973. ``Functions.'' \emph{The Philosophical Review} 82(2): 139--168. \href{https://doi.org/10.2307/2183766}{DOI: 10.2307/2183766}.

\par\endgroup

\backmatter
\end{document}
